% =====================================================================
%  QUANTITATIVE PORTFOLIO MANAGEMENT -- EXTENDED LECTURE NOTES
%  EDHEC Business School, 2026-2027 -- Prof. Raman Uppal
%  Part 1: Chapters 0-4 (Sessions 1-4), Assignments 1-2.
%  Lightweight version: only standard packages, fast to compile on Overleaf.
% =====================================================================
\documentclass[11pt,a4paper]{report}

\usepackage[utf8]{inputenc}
\usepackage[T1]{fontenc}
\usepackage{lmodern}
\usepackage[margin=2.3cm]{geometry}
\usepackage{amsmath,amssymb,bm}
\usepackage{booktabs,tabularx,longtable}
\usepackage{graphicx}
\usepackage{enumitem}
\usepackage{listings}
\usepackage{needspace}
\setlist{itemsep=2pt,topsep=4pt}
\emergencystretch=3em
\setlength{\parindent}{0pt}
\setlength{\parskip}{4pt}

% ---------- Simple boxes (no tcolorbox: fast to compile) ----------
% Each "box" is an indented block with a bold title and rules above/below.
\newenvironment{notebox}[1]
  {\par\medskip\Needspace{6\baselineskip}\noindent\rule{\linewidth}{0.4pt}\par\nopagebreak
   \noindent\textbf{#1}\par\nopagebreak
   \begin{list}{}{\leftmargin=0.8em\rightmargin=0em\topsep=2pt}\item[]}
  {\end{list}\nopagebreak\noindent\rule{\linewidth}{0.4pt}\par\medskip}
\newenvironment{keybox}[1][]{\begin{notebox}{$\blacktriangleright$ Key result#1}}{\end{notebox}}
\newenvironment{definitionbox}[1]{\begin{notebox}{Definition: #1}}{\end{notebox}}
\newenvironment{intuition}{\begin{notebox}{Intuition}}{\end{notebox}}
\newenvironment{examtip}{\begin{notebox}{Exam / quiz tip}}{\end{notebox}}
\newenvironment{plainbox}{\begin{notebox}{In plain words}}{\end{notebox}}
\newenvironment{addnote}{\begin{notebox}{Additional note (not on the slides)}}{\end{notebox}}
\newenvironment{workedex}[1][]{\begin{notebox}{Worked example#1}}{\end{notebox}}
\newenvironment{mistakes}{\begin{notebox}{Common mistakes to avoid}}{\end{notebox}}

% ---------- Python listings (plain, no colours) ----------
\lstdefinestyle{py}{language=Python, basicstyle=\ttfamily\footnotesize,
  breaklines=true, frame=single, columns=fullflexible, keepspaces=true,
  showstringspaces=false, upquote=true, tabsize=4, xleftmargin=2pt, xrightmargin=2pt}
\lstset{style=py}
\lstdefinestyle{out}{language={}, basicstyle=\ttfamily\footnotesize, frame=leftline,
  breaklines=true, columns=fullflexible, keepspaces=true, xleftmargin=1em}

% ---------- Math shortcuts ----------
\newcommand{\E}{\mathbb{E}}
\newcommand{\V}{\mathbb{V}}
\newcommand{\one}{\mathbf{1}_N}
\newcommand{\w}{\bm{w}}
\newcommand{\R}{\bm{R}}
\newcommand{\T}{^{\top}}
\newcommand{\Rf}{R_f}
\newcommand{\mub}{\bm{\mu}}
\newcommand{\url}[1]{\texttt{#1}}

\title{\textbf{Quantitative Portfolio Management:\\Theory and Practice}\\[6pt]
\Large Part 1: Chapters 1--4\\[4pt]
\large Extended lecture notes and class companion}
\author{EDHEC Business School \textbullet\ Prof.\ Raman Uppal\\[4pt]
\small Personal study notes based on the course slides, 2026--2027}
\date{Last updated: \today}

% =====================================================================
\begin{document}
\maketitle

\begin{abstract}
\noindent These notes follow the course slides chapter by chapter and are
designed to be used \emph{during} class. Each chapter contains: (i) ``In plain
words'' explanations, (ii) the theory with full formulas and step-by-step
derivations, (iii) the intuition behind every expression, (iv) the Python
exercises from the slides with their expected output and interpretation, (v)
takeaways and self-test questions for the in-class quizzes and the handwritten
final exam. The appendix contains question-by-question guides to the
assignments, a formula sheet and an English--Italian glossary. Grey boxes mark
material that is \emph{not} on the slides.
\end{abstract}

\tableofcontents

% ===== begin intro =====
\chapter*{How to use these notes}
\addcontentsline{toc}{chapter}{How to use these notes}
\markboth{How to use these notes}{}

\paragraph{Purpose.} A companion to keep open during class: the slides say
\emph{what}, these notes also say \emph{why} and \emph{how}. Each chapter
follows the slides in order, so you can find the page the professor is on.

\paragraph{The boxes.}
\begin{center}
\begin{tabularx}{\textwidth}{@{}lX@{}}
\toprule
Box & Meaning \\
\midrule
\textbf{In plain words} & The idea explained simply, without formulas. Read this first if you are lost. \\
\textbf{Key result} ($\blacktriangleright$) & Formula or fact you must know for the quiz and the exam. \\
\textbf{Definition} & A precise definition. \\
\textbf{Intuition} & Why the formula makes economic sense (the professor always asks this). \\
\textbf{Additional note} & Derivations, numerical checks and code explanations that are \emph{not} on the slides. \\
\textbf{Worked example} & A numerical example solved step by step (many are not on the slides). \\
\textbf{Common mistakes} & Errors students typically make in code and in exam answers. \\
Code listings & Python from the slides (cleaned), with the expected output in a separate frame. \\
\bottomrule
\end{tabularx}
\end{center}

\paragraph{The course in one page.}
The whole course is built around \textbf{one formula}, the optimal
mean--variance portfolio
\[
\w=\frac1\gamma\,V^{-1}\,(\mub-R_f\one)
\qquad\text{(weights = risk tolerance $\times$ inverse covariance $\times$ risk premia).}
\]
\begin{enumerate}
  \item \textbf{Chapter 1 (data):} how to get prices, turn them into returns, and
    estimate the two inputs $\mub$ and $V$. Main lesson: $\hat\mub$ is very noisy,
    $\hat V$ is better but hard to invert when $N$ is large.
  \item \textbf{Chapter 2 (performance):} how to judge a strategy honestly
    (out of sample), which numbers to look at (Sharpe ratio, risk measures,
    turnover) and how to test whether a difference is real or luck.
  \item \textbf{Chapter 3 (theory):} what the formula gives if we \emph{knew}
    $\mub$ and $V$: efficient frontier, tangency portfolio, beta relation.
  \item \textbf{Chapter 4 (practice):} what happens with \emph{estimated} inputs
    (disaster), and the fixes: Bayesian shrinkage, constraints, GMV,
    Ledoit--Wolf, norm constraints. Common thread: \textbf{shrinkage}.
  \item \textbf{After Part 1:} the rest of the course uses economic models to get
    better inputs (covered in Part 2 of these notes).
\end{enumerate}


\paragraph{Lectures, assignments and exam questions.}
\begin{center}
\begin{tabular}{@{}lllll@{}}
\toprule
Sessions & Chapters & Assignment & Due & Exam question \\
\midrule
1--2 & 1 Data, 2 Performance & \#1 & 11 Nov & Q1 \\
3--4 & 3 MV theory, 4 Estimation error & \#2 & 16 Nov & Q2 \\
5--6 & 5 Black--Litterman, 6 Parametric policies & \#3 & 25 Nov & Q3 \\
7--8 & 7 Volatility timing, 8 Unsystematic risk & \#4 & 30 Nov & Q4 \\
\bottomrule
\end{tabular}
\end{center}

\paragraph{Study routine that matches the grading.}
\begin{enumerate}
  \item \textbf{Before class:} read the ``In plain words'' boxes of the chapter (10 minutes).
  \item \textbf{During class:} follow the slides here; mark what the professor stresses.
  \item \textbf{After class:} do the self-test questions at the end of the chapter (they mimic the quizzes).
  \item \textbf{Weekly:} rewrite by hand, without looking, the ``write by hand''
    code at the end of each chapter (the exam is code on paper).
  \item \textbf{Before the exam:} the Formula sheet and the Glossary at the end of these notes.
\end{enumerate}

% ===== end intro =====


% =====================================================================
\chapter*{Lecture log}
\addcontentsline{toc}{chapter}{Lecture log}
\markboth{Lecture log}{}

\begin{center}\small
\begin{tabular}{@{}clllc@{}}
\toprule
Sess. & Date & Chapter & Due (22:00 CET) & Notes \\
\midrule
1 & 02 Nov & Ch.\,0 Syllabus; Ch.\,1 Financial data & & \checkmark \\
2 & 03 Nov & Ch.\,2 Performance measurement & \#1: 11 Nov & \checkmark \\
3 & 04 Nov & Ch.\,3 MV portfolios ignoring estimation error & & \checkmark \\
4 & 05 Nov & Ch.\,4 MV portfolios adjusting for estimation error & \#2: 16 Nov & \checkmark \\
5 & 16 Nov & Ch.\,5 Black--Litterman & & \\
6 & 17 Nov & Ch.\,6 Parametric portfolio policies & \#3: 25 Nov & \\
7 & 18 Nov & Ch.\,7 Volatility-timed factor portfolios & & \\
8 & 19 Nov & Ch.\,8 Systematic and unsystematic risk & \#4: 30 Nov & \\
\midrule
Exam & 16 Dec & Final exam, 9:00, closed-book, on paper & & \\
\bottomrule
\end{tabular}
\end{center}

% =====================================================================
\chapter*{Chapter 0 -- Syllabus and course organisation}
\addcontentsline{toc}{chapter}{Chapter 0 -- Syllabus and course organisation}
\markboth{Chapter 0 -- Syllabus}{}

\section*{0.1 Instructor}
\begin{itemize}
  \item Raman \textsc{Uppal}, Professor of Finance, EDHEC Business School.
  \item Email: \texttt{Raman.Uppal@edhec.edu} --- start the subject with
        ``\texttt{QPM-2026-2027}''; replies within 24 hours.
  \item Office hours: by appointment and by email. Homepage: \url{https://www.ramanuppal.com}.
\end{itemize}

\section*{0.2 Objectives and learning outcomes}
The objective is to study cutting-edge methods to construct \textbf{optimal
equity portfolios that perform well out of sample}, combining
\emph{finance} (portfolio management), \emph{econometrics} (empirical finance)
and \emph{data science} (Python).

\begin{enumerate}[label=\textbf{LO\arabic*}]
  \item Use Python to obtain data from Yahoo!\ Finance, Nasdaq Data Link
        (Quandl), Ken French Data Library, Open Source Asset Pricing.
  \item Use Python to evaluate the performance of portfolios.
  \item Use Python to construct mean--variance efficient portfolios with
        short-sale constraints, position limits, LASSO and elastic-net
        constraints, and transaction costs.
  \item Use Python to improve performance by adjusting for estimation error in
        means, variances, covariances, and by using conditioning information.
\end{enumerate}

\textbf{Prerequisites:} calculus and linear algebra; Markowitz optimisation;
basic econometrics; Python. The course is \emph{not} an introduction to Python,
not a machine-learning course, and does not discuss AI/LLMs.

\section*{0.3 Assessment}
\begin{keybox}[: grading]
\begin{itemize}
  \item \textbf{15\%} --- four assignments and eight in-class online quizzes.
    \[
      \text{Grade}_j = [\text{Quiz } j \text{ grade}] \times
      \mathbb{1}\{\text{corresponding assignment submitted}\}.
    \]
    The lowest quiz score is dropped. Assignments are \emph{not graded}
    (AI is allowed): full credit if submitted on time.
  \item \textbf{85\%} --- final exam: handwritten, in class, closed-book,
    four questions matching the four assignments; tests portfolio theory and
    the ability to \textbf{write Python code by hand}.
\end{itemize}
Quizzes cannot be made up; late assignments get 0.
\end{keybox}

\section*{0.4 Assignment rules (failure $\Rightarrow$ grade 0)}
\begin{itemize}
  \item Groups of exactly 3 or 4 students (diversity encouraged).
  \item Email to \texttt{Raman.Uppal@edhec.edu}, with subject line
        ``\texttt{QPM 2026-2027: Assignment n}'', where $n\in\{1,2,3,4\}$.
  \item Executable Jupyter notebook named
        \texttt{QPM-2026-Assignment-n-A-B-C-D.ipynb} (A--D = last names).
  \item Section 0 (Markdown): line 1 submission date; line 2
        ``QPM 2026-2027: Assignment n''; line 3 group members alphabetically,
        LAST NAME in capitals; line 4 comments/challenges.
  \item Section $k$: copy Question $k$, then the answer; include all imports.
  \item Before sending: \emph{Kernel $\to$ Restart Kernel and Run All Cells}.
  \item Mention any outside help (no penalty).
\end{itemize}

\section*{0.5 Class rules}
Attendance mandatory; double-badging $\Rightarrow$ 0 for the whole course for
both students (to miss a class, email the HUB and copy the professor). Bring a
laptop (quizzes) and a name card. No phones; laptop only for the course.
Expected workload: about 2 hours of study per hour of lecture.

\section*{0.6 Course map and difficulty}
\begin{center}
\begin{tabular}{@{}llcc@{}}
\toprule
Part & Chapter & Theory & Coding \\
\midrule
A. Preliminaries & 1 Managing financial data & Low & Low \\
 & 2 Performance measurement & Low & Low \\
B. Modern PM & 3 MV ignoring estimation error & Low & Medium \\
 & 4 MV adjusting for estimation error & Medium & Medium \\
C. Post-modern PM & 5 Black--Litterman & High & Medium \\
 & 6 Parametric portfolio policies & Medium & Medium \\
 & 7 Volatility-timed factor portfolios & Medium & Medium \\
 & 8 Systematic \& unsystematic risk & Medium & Medium \\
\bottomrule
\end{tabular}
\end{center}

\paragraph{Topics in detail.}
\begin{itemize}
  \item \textbf{Ch.\,2}: out-of-sample performance metrics; VaR, Expected
    Shortfall, Maximum Drawdown; backtesting and manipulation-proof measures;
    transaction and price-impact costs; bootstrapped $p$-values (Ledoit--Wolf
    2008); $p$-hacking and multiple testing.
  \item \textbf{Ch.\,3}: Markowitz theory (Merton 1972; Roll 1980;
    Huang--Litzenberger 1988); in-sample performance; short-sale, position
    (LASSO, elastic-net) and leverage constraints; transaction costs; higher
    moments; out-of-sample performance.
  \item \textbf{Ch.\,4}: Bayesian approaches (Jorion 1985--1992); three-fund
    portfolio (Kan--Zhou 2007); GMV portfolio; short-sale constraints
    (Jagannathan--Ma 2003); covariance shrinkage (Ledoit--Wolf 2003--2017);
    norm-constrained portfolios (DeMiguel, Garlappi, Nogales, Uppal 2009).
  \item \textbf{Ch.\,5}: Black--Litterman (1990--1992), intuition (He--Litterman
    1999), derivation, numerical example.
  \item \textbf{Ch.\,6}: single/multi-factor models; fundamental indexation
    (Arnott--Hsu--Moore 2005); risk parity (Roncalli 2013); parametric
    portfolio policies (Brandt--Santa-Clara--Valkanov 2009).
  \item \textbf{Ch.\,7}: volatility timing (Kirby--Ostdiek 2012; Moreira--Muir
    2017, 2019); net of costs (Barroso--Detzel 2021); out of sample
    (Cederburg et al.\ 2020); both (DeMiguel, Mart\'in-Utrera, Uppal 2024).
  \item \textbf{Ch.\,8}: asset pricing with unsystematic risk
    (Dello-Preite, Uppal, Zaffaroni, Zviadadze 2025); portfolio choice with
    unsystematic risk (Raponi, Uppal, Zaffaroni 2023).
\end{itemize}

\section*{0.7 Resources}
\begin{itemize}
  \item Slides (main material). Red = important; blue = links.
  \item \textbf{Tidy Finance with Python} (Scheuch, Voigt, Weiss, Frey 2024) ---
        main Python reference. Tidy-data principles used throughout.
  \item Bali, Engle, Murray (2016), \emph{Empirical Asset Pricing} --- empirical finance.
  \item Moreira (2021), \emph{Quantitative Investing} (online, with Python).
  \item EDHEC MOOC ``Investment Management with Python and Machine Learning''
        (Coursera) --- pre-course Chapter 0.
\end{itemize}

% =====================================================================
\chapter{Managing Financial Data}
\label{ch:1}

\section{What we want to do in Chapter 1}
\begin{itemize}
  \item Take a big-picture view of the whole course.
  \item Study the main kinds of data used to construct optimal portfolios.
  \item Learn to obtain this data with Python and store it efficiently.
\end{itemize}
\begin{plainbox}
Chapter 1 sets the stage. The whole course is about one formula that tells you
how much to put in each stock. The formula needs two inputs: the \emph{expected
returns} and the \emph{covariance matrix}. Both must be estimated from data, and
this chapter shows (i) how to get and clean the data and (ii) why those
estimates are much noisier than you would think. Everything else in the course is
a way to deal with this noise.
\end{plainbox}
The two questions of the course are:
\textbf{Q1.} How should we construct optimal portfolios?
\textbf{Q2.} How should we measure portfolio performance?

% ---------------------------------------------------------------------
\section{Overview of the course: mean--variance efficient portfolios}

The course focuses on \textbf{mean--variance efficient} weights, which
trade off optimally the portfolio's return \emph{mean} and return
\emph{variance}. For most of the course we do \emph{not} focus on higher
moments (skewness, kurtosis), other risk measures (VaR, ES) or general
expected-utility objectives.

\subsection{One risky asset ($N=1$)}
An investor splits wealth between one risky asset (e.g.\ the market) and the
risk-free asset. The optimal weight in the risky asset is
\begin{keybox}[: optimal weight, $N=1$]
\[
  w_1 = \frac{1}{\gamma}\,\frac{\E[R_1]-\Rf}{\V[R_1]},
\]
$w_1$: fraction of wealth in the risky asset; $\gamma$: risk aversion;
$\E[R_1]$: expected gross return; $\Rf$: gross risk-free return;
$\V[R_1]=\sigma^2_{R_1}=\sigma_{R_{11}}$: return variance.
\end{keybox}

\begin{intuition}
\begin{enumerate}
  \item $\gamma\uparrow\ \Rightarrow\ w_1\downarrow$ (more risk-averse, less risky asset).
  \item Risk premium $\E[R_1]-\Rf\uparrow\ \Rightarrow\ w_1\uparrow$.
  \item Risk $\V[R_1]\uparrow\ \Rightarrow\ w_1\downarrow$.
\end{enumerate}
The professor insists: \emph{every formula in the course must make intuitive sense.}
\end{intuition}

\begin{addnote}
\textbf{Where the formula comes from.} The mean--variance investor maximises
\[
  \max_{w}\; \underbrace{\Rf + w(\E[R_1]-\Rf)}_{\text{portfolio mean}}
  \;-\; \frac{\gamma}{2}\,\underbrace{w^2\,\V[R_1]}_{\text{portfolio variance}}.
\]
First-order condition: $\E[R_1]-\Rf-\gamma w\,\V[R_1]=0 \Rightarrow
w=\frac{1}{\gamma}\frac{\E[R_1]-\Rf}{\V[R_1]}$. The objective is concave in
$w$ (second derivative $-\gamma\V[R_1]<0$), so this is a maximum.
\emph{Example:} $\E[R_1]-\Rf=6\%$, $\sigma=20\%$, $\gamma=3$:
$w=\frac{1}{3}\cdot\frac{0.06}{0.04}=0.5$, i.e.\ 50\% in the market.
\end{addnote}

\subsection{Many risky assets ($N>1$)}
\begin{keybox}[: the key expression of the course]
\[
  \w = \frac{1}{\gamma}\,\V[\R]^{-1}\big(\E[\R]-\Rf\,\one\big),
\]
with
\[
\w=\begin{pmatrix}w_1\\ \vdots\\ w_N\end{pmatrix},\quad
\E[\R]=\begin{pmatrix}\E[R_1]\\ \vdots\\ \E[R_N]\end{pmatrix},\quad
\one=\begin{pmatrix}1\\ \vdots\\ 1\end{pmatrix},\quad
\V[\R]=\begin{pmatrix}
\sigma_{R_{11}} & \cdots & \sigma_{R_{1N}}\\
\vdots & \ddots & \vdots\\
\sigma_{R_{N1}} & \cdots & \sigma_{R_{NN}}
\end{pmatrix},
\]
where $\sigma_{R_{nn}}=\sigma^2_{R_n}$ is the variance of asset $n$ and
$\sigma_{R_{nm}}$ the covariance between $n$ and $m$. The weight $1-\w\T\one$
goes to the risk-free asset.
\end{keybox}

\begin{plainbox}
Compared with $N=1$, the only new thing is $V^{-1}$ instead of $1/\sigma^2$. It
does two jobs: it penalises volatile assets (like $1/\sigma^2$) and it accounts
for \emph{correlations}: if two assets move together, holding both is like holding
one, so the formula avoids double-counting them, and it may even short one to
hedge the other.
\end{plainbox}

\begin{workedex}[: two risky assets by hand]
Excess means $\mub-R_f\one=(0.06,\ 0.04)\T$, volatilities $(0.20,\ 0.10)$,
correlation $0.5$, risk aversion $\gamma=3$.

\textbf{Step 1: covariance matrix.} $\sigma_{12}=\rho\sigma_1\sigma_2=0.5\times0.2\times0.1=0.01$:
\[
V=\begin{pmatrix}0.04&0.01\\0.01&0.01\end{pmatrix}.
\]
\textbf{Step 2: inverse of a $2\times2$ matrix.} Swap the diagonal, change the
sign of the off-diagonal, divide by the determinant $ad-bc=0.04\times0.01-0.01^2=0.0003$:
\[
V^{-1}=\frac{1}{0.0003}\begin{pmatrix}0.01&-0.01\\-0.01&0.04\end{pmatrix}
=\begin{pmatrix}33.33&-33.33\\-33.33&133.33\end{pmatrix}.
\]
\textbf{Step 3: multiply.}
$V^{-1}(\mub-R_f\one)=(33.33\cdot0.06-33.33\cdot0.04,\ -33.33\cdot0.06+133.33\cdot0.04)=(0.667,\ 3.333)$.

\textbf{Step 4: divide by $\gamma$.} $\w=(0.222,\ 1.111)$. The weights sum to
$1.333$, so the investor \emph{borrows} 33\% at the risk-free rate
($w_f=1-1.333=-0.333$).

\textbf{Reading the result.} Asset 1 has the higher premium (6\% vs.\ 4\%) but gets
the \emph{smaller} weight: it is four times as risky ($0.04$ vs.\ $0.01$) and,
because the two assets are correlated, asset 2 already captures much of asset 1's
exposure. This is why you cannot choose weights by looking at assets one at a time.
\end{workedex}
\begin{addnote}
Same derivation in matrix form:
$\max_{\w}\ \w\T(\E[\R]-\Rf\one)-\frac{\gamma}{2}\w\T\V[\R]\w$.
FOC: $(\E[\R]-\Rf\one)-\gamma\V[\R]\w=\bm 0$ $\Rightarrow$ the formula above.
Note that the $N=1$ case is just the scalar version: $\V^{-1}$ becomes $1/\sigma^2$.
\end{addnote}

% ---------------------------------------------------------------------
\section{Estimation error in mean--variance weights (focus)}

Since $\w=\frac1\gamma\V[\R]^{-1}(\E[\R]-\Rf\one)$, the whole challenge is to
\textbf{estimate $\E[\R]$ and $\V[\R]$ precisely}.

\begin{plainbox}
If the inputs are wrong, the weights are wrong. Expected returns are tiny
compared with the noise around them (a stock may average 8\% a year but move
$\pm30\%$), so after 10 years you still do not know the mean well. The covariance
matrix has thousands of entries for a large $N$, and inverting it magnifies
every small error.
\end{plainbox}
\subsection{Why $\E[\R]$ is hard to estimate}
Precision of the mean does \emph{not} improve by sampling more often over a
fixed span:
\[
  R_{0,100}=\frac{P_{100}}{P_0}
  =\frac{P_{100}}{P_{99}}\times\frac{P_{99}}{P_{98}}\times\cdots\times\frac{P_1}{P_0}
  =\frac{P_{100}}{P_0}.
\]
All intermediate prices cancel: only the first and last prices matter.
What helps is a \emph{longer} sample (more years), not a \emph{denser} one.

\subsection{Why $\V[\R]$ is hard to estimate}
Number of parameters:
\[
  \underbrace{N}_{\text{variances}}+\underbrace{\tfrac{N^2-N}{2}}_{\text{covariances}}
  =\frac{N(N+1)}{2}. \qquad N=100\ \Rightarrow\ 5050 \text{ parameters.}
\]

\subsection{Why poor estimates are disastrous}
\begin{itemize}
  \item $\widehat{\E[\R]}$ has large error;
  \item $\widehat{\V[\R]}^{-1}$ has large error (high \textbf{condition number});
  \item multiplying by $\V^{-1}$ \textbf{magnifies} the error in $\E[\R]$;
  \item $\Rightarrow$ badly estimated $\w$ $\Rightarrow$ poor out-of-sample performance.
\end{itemize}

\begin{definitionbox}{Condition number}
$\kappa(\V)=\lambda_{\max}/\lambda_{\min}$, the ratio of the largest to the
smallest eigenvalue of a symmetric positive-definite matrix. A large $\kappa$
means $\V$ is close to singular, and small errors in inputs produce large
errors in $\V^{-1}\bm{x}$.
\end{definitionbox}

\begin{addnote}
\textbf{Why small eigenvalues cause big weights.} Write $\V=\sum_i\lambda_i
\bm{q}_i\bm{q}_i\T$, so $\V^{-1}=\sum_i\lambda_i^{-1}\bm{q}_i\bm{q}_i\T$. The
directions with the smallest (and most under-estimated) eigenvalues get
multiplied by the largest factors $1/\lambda_i$. With $T$ close to $N$, the
sample eigenvalues spread out: the smallest are biased towards zero, so
$\widehat{\V}^{-1}$ explodes in those directions. If $N>T$, $\widehat{\V}$ is
singular and cannot even be inverted.
\end{addnote}


\begin{workedex}[: condition number of a $2\times2$ matrix]
Two assets with the same variance $\sigma^2$ and correlation $\rho$. The
eigenvalues of $V=\sigma^2\begin{pmatrix}1&\rho\\\rho&1\end{pmatrix}$ are
$\sigma^2(1+\rho)$ (direction ``both assets together'') and $\sigma^2(1-\rho)$
(direction ``long one, short the other''). So
\[
\kappa=\frac{1+\rho}{1-\rho}:\qquad \rho=0.5\Rightarrow\kappa=3,\qquad
\rho=0.9\Rightarrow\kappa=19,\qquad \rho=0.99\Rightarrow\kappa=199 .
\]
The more similar the assets, the smaller the variance of the long--short
direction, and $V^{-1}$ multiplies that direction by $1/(\sigma^2(1-\rho))$, which is
huge. Any small difference in estimated means along that direction is turned into
an enormous long--short bet. This is exactly the $+635\%/-535\%$ example of
Chapter 4 (correlation 0.99).
\end{workedex}
\subsection*{Python exercise 1: why the sample $\V[\R]$ is ill-conditioned}
$N$ identical assets, 20\% annual volatility, pairwise correlation 0.3,
$T=60$ months. True GMV weights $=1/N$ (assets identical). Compare
$\kappa(\V)$ vs.\ $\kappa(\widehat{\V})$ and estimated GMV weights.

\begin{lstlisting}
import numpy as np
import pandas as pd
rng = np.random.default_rng(1)
T = 60                   # 5 years of monthly returns
sd = 0.20 / np.sqrt(12)  # 20% annual volatility
rho = 0.3                # correlation between all pairs
Ns = [5, 10, 25, 50]     # numbers of assets
# True covariance matrices, simulated returns, sample covariances
V = {N: sd**2 * (rho*np.ones((N, N)) + (1-rho)*np.eye(N)) for N in Ns}
R = {N: rng.multivariate_normal(np.zeros(N), V[N], size=T) for N in Ns}
Vhat = {N: np.cov(R[N], rowvar=False) for N in Ns}

# GMV weights from the sample covariance matrix
w = {N: np.linalg.solve(Vhat[N], np.ones(N)) for N in Ns}
w = {N: w[N] / w[N].sum() for N in Ns}
table = pd.DataFrame({
    "cond(V)":    [np.linalg.cond(V[N]) for N in Ns],
    "cond(Vhat)": [np.linalg.cond(Vhat[N]) for N in Ns],
    "true w":     [1/N for N in Ns],
    "min w":      [w[N].min() for N in Ns],
    "max w":      [w[N].max() for N in Ns]}, index=Ns)
print(table.round(2))
\end{lstlisting}
\begin{lstlisting}[style=out]
    cond(V)  cond(Vhat)  true w  min w  max w
5      3.14        7.22    0.20   0.05   0.38
10     5.29       22.12    0.10  -0.04   0.28
25    11.71       61.01    0.04  -0.13   0.26
50    22.43     1864.26    0.02  -0.39   0.47
\end{lstlisting}
\textbf{Interpretation:} with $N=50$, $T=60$, $\kappa(\widehat{\V})$ is about
80 times $\kappa(\V)$; true weights are all 0.02 but estimated ones range from
$-0.39$ to $0.47$.

\begin{addnote}
\textbf{Code points to remember for the handwritten exam.}
\begin{itemize}
  \item Equicorrelation matrix: $\sigma^2[\rho\,\bm{1}\bm{1}\T+(1-\rho)I]$.
  \item GMV weights: $\w_{GMV}=\dfrac{\V^{-1}\one}{\one\T\V^{-1}\one}$. In code,
    \texttt{np.linalg.solve(V, ones)} computes $\V^{-1}\one$ without explicitly
    inverting (more stable than \texttt{np.linalg.inv}); dividing by the sum
    normalises to $\one\T\w=1$.
  \item \texttt{np.cov(X, rowvar=False)}: columns are variables (assets).
  \item \texttt{np.linalg.cond}: condition number (2-norm).
\end{itemize}
\end{addnote}

\paragraph{The solution (rest of the course).} Find better estimators of
$\E[\R]$ and $\V[\R]$ that reduce estimation error and dimension. $\E[\R]$
and $\V[\R]$ matter beyond portfolios too: cost of capital needs $\E[\R]$,
risk management needs $\V[\R]$.

% ---------------------------------------------------------------------
\section{Portfolio weights over time and out-of-sample returns}

We choose weights at each date $t\in\{1,\dots,T\}$:
$\w_t=(w_{t,1},\dots,w_{t,N})\T$, with transpose $\w_t\T=[w_{t,1},\dots,w_{t,N}]$.
\[
  \sum_{n=1}^N w_{t,n}=\w_t\T\one,\qquad
  \E[R_{t,p}]=\sum_{n=1}^N w_{t,n}\E[R_{t,n}]=\w_t\T\E[\R_t].
\]

\begin{keybox}[: out-of-sample return]
Weights $\w_t$ use information up to $t$; performance depends on returns at $t+1$:
\[
  R_{t+1,p}=\w_t\T\R_{t+1}=\sum_{n=1}^N w_{t,n}R_{t+1,n}.
\]
This is \textbf{out-of-sample} performance. Our choice of $\w_t$ depends on our
views about $R_{t+1,n}$.
\end{keybox}

\begin{plainbox}
``Out of sample'' means: decide today using only what you know today, then see
what happens tomorrow. Any backtest that uses tomorrow's information to decide
today's weights is cheating (look-ahead bias) and will look far better than reality.
\end{plainbox}
\paragraph{What can weights depend on?} The asset's expected return and risk
(variance, skewness, kurtosis, downside risk, VaR, ES, tail risk), which in turn
may depend on: past returns of the asset and of other assets; firm
characteristics (size, profitability, \dots); covariances/correlations with
other assets; sensitivity to macroeconomic factors.

\paragraph{Data dimensions.} Across $N$ assets (\textbf{cross-section}), over
$T$ dates (\textbf{time series}), or both (\textbf{panel}).

% ---------------------------------------------------------------------
\section{Notation}
\begin{itemize}
  \item $N$, $n\in\{1,\dots,N\}$: number of assets.
  \item $T$, $t\in\{1,\dots,T\}$: number of observations.
  \item $K$, $k\in\{1,\dots,K\}$: number of factors.
  \item $T^{\text{est}}\le T$: length of the \textbf{estimation window}. Not
    standard across papers: also $M$ (months, usually 60 or 120), $D$ (days,
    usually 30, 60, 90), or $\tau$.
\end{itemize}
\begin{center}
\begin{tabular}{@{}ll@{}}
\toprule
Our notation & Alternatives in the literature \\
\midrule
$\w$ & $\bm{x}$, $\bm\theta$ \\
$\gamma$ & $\gamma=1/\tau$, $\tau$ = risk tolerance \\
$\V[\R]$ or $V_R$ & $V$, $\Sigma_R$, $\Sigma$, $C$ \\
$\E[\R]$ or $E_R$ & $\mu_R$, $\mu$ \\
$\Rf$ (gross) & $r_f$ (net) \\
$\one$ & $\bm 1$, $\bm e$, $\bm\iota$ \\
$T^{\text{est}}$ & $M$, $D$, $\tau$ \\
\bottomrule
\end{tabular}
\end{center}

% ---------------------------------------------------------------------
\section{Convention regarding data: tidy vs.\ wide}
Tidy-data principles (Tidy Finance): \textbf{(1)} every column is a variable;
\textbf{(2)} every row is an observation; \textbf{(3)} every cell is a single value.

\begin{center}
\begin{minipage}{0.38\textwidth}\centering
\textbf{Tidy (long)}: $T\times N$ rows\\[3pt]
\begin{tabular}{ccc}
\toprule
date & asset & price\\\midrule
1 & 1 & $P_{1,1}$\\
1 & 2 & $P_{1,2}$\\
$\vdots$ & $\vdots$ & $\vdots$\\
$T$ & $N$ & $P_{T,N}$\\\bottomrule
\end{tabular}
\end{minipage}\hfill
\begin{minipage}{0.55\textwidth}\centering
\textbf{Wide}: dates in rows, assets in columns\\[3pt]
\begin{tabular}{cccc}
\toprule
 & $n=1$ & $\cdots$ & $n=N$\\\midrule
$t=1$ & $P_{1,1}$ & $\cdots$ & $P_{1,N}$\\
$\vdots$ & $\vdots$ & $\ddots$ & $\vdots$\\
$t=T$ & $P_{T,1}$ & $\cdots$ & $P_{T,N}$\\\bottomrule
\end{tabular}
\end{minipage}
\end{center}
\begin{keybox}[: rule]
Store data in \textbf{tidy} format; \textbf{pivot to wide} only for matrix
calculations (covariances, correlations, $\V^{-1}$).
\end{keybox}

\subsection*{Python exercise 2: wide and tidy data}
\begin{lstlisting}
import numpy as np
import pandas as pd
pd.set_option("display.precision", 4)
rng = np.random.default_rng(2)
dates = pd.date_range("2024-01-31", periods=24, freq="ME")
ret = rng.normal(0.01, 0.06, size=(24, 3))
prices_wide = pd.DataFrame(100*np.cumprod(1 + ret, axis=0),
    index=dates, columns=["AAA", "BBB", "CCC"])
prices_wide.index.name = "date"
# Wide -> long (tidy): one row per (date, symbol)
prices = (prices_wide.reset_index()
    .melt(id_vars="date", var_name="symbol", value_name="price")
    .sort_values(["symbol", "date"]))

# Wrong: pct_change() runs across the stocks
prices["ret_wrong"] = prices["price"].pct_change()
# Right: compute returns within each stock
g = prices.groupby("symbol")["price"]
prices["ret"] = g.pct_change()
print(prices.iloc[22:26])
# Long -> wide, as needed for covariance matrices
ret_wide = prices.pivot(index="date", columns="symbol", values="ret")
print(ret_wide.dropna().corr().round(2))
\end{lstlisting}
\begin{lstlisting}[style=out]
row  date        symbol  price     ret_wrong  ret
22   2025-11-30  AAA     72.2815   -0.0818    -0.0818
23   2025-12-31  AAA     73.3244    0.0144     0.0144
24   2024-01-31  BBB     97.8635    0.3347     NaN
25   2024-02-29  BBB    109.4097    0.1180     0.1180
\end{lstlisting}
\textbf{Interpretation:} in row 24 the ``wrong'' return of 33.47\% is the change
from the last price of AAA to the first of BBB --- not a return. With
\texttt{groupby}, the first return of each stock is correctly \texttt{NaN}.

\begin{addnote}
Key pandas verbs: \texttt{melt} (wide$\to$long), \texttt{pivot}
(long$\to$wide), \texttt{groupby(...).pct\_change()} (returns within each
asset), \texttt{sort\_values(["symbol","date"])} (needed before
\texttt{pct\_change}). With \texttt{yfinance} multi-ticker downloads,
\texttt{.stack(level="Ticker")} also converts wide to long.
\end{addnote}

% ---------------------------------------------------------------------
\section{Prices and returns}

\begin{keybox}[: return definitions]
\begin{align*}
\text{Gross return:}\quad & R_{t,n}=\frac{P_{t,n}}{P_{t-1,n}}\\
\text{Net return:}\quad & r_{t,n}=R_{t,n}-1=\frac{P_{t,n}}{P_{t-1,n}}-1\\
\text{Excess return:}\quad & R_{t,n}-R_{t,0}=\frac{P_{t,n}}{P_{t-1,n}}-\frac{P_{t,0}}{P_{t-1,0}}
\end{align*}
$R_{t,0}$: reference asset, usually risk-free ($R_{t,f}$). Excess gross returns
$=$ excess net returns (the $-1$'s cancel).
\end{keybox}

\subsection{Multiperiod (compound) returns}
\[
R_{t,n}(h)=\frac{P_{t,n}}{P_{t-h,n}}
=\frac{P_{t,n}}{P_{t-1,n}}\times\cdots\times\frac{P_{t-h+1,n}}{P_{t-h,n}}
=\prod_{i=0}^{h-1}R_{t-i,n}.
\]

\subsection{Annualised returns}
Returns are quoted per year. Monthly net returns are annualised by
multiplying by 12; annualised gross $=1+$ annualised net. With one-year
returns over $h$ years:
\[
\underbrace{R_{t,n}(h)^{1/h}=\Big[\prod_{i=0}^{h-1}R_{t-i,n}\Big]^{1/h}}_{\text{geometric mean (exact)}}
\;\approx\;
\underbrace{\frac1h\sum_{i=0}^{h-1}R_{t-i,n}}_{\text{arithmetic mean}}.
\]

\begin{addnote}
\textbf{Arithmetic vs.\ geometric mean.} Always (AM--GM inequality)
arithmetic $\ge$ geometric, and approximately
\[
\bar r_{\text{arith}}-\bar r_{\text{geom}}\approx \tfrac12\sigma^2 .
\]
\emph{Classic example:} $+50\%$ then $-50\%$: arithmetic mean $=0\%$, but
wealth goes $1\to1.5\to0.75$, geometric mean $=\sqrt{0.75}-1\approx-13.4\%$.
\end{addnote}

\subsection{Continuous compounding (log returns)}
Let $p_{t,n}=\ln P_{t,n}$. Then $\ln R_{t,n}=p_{t,n}-p_{t-1,n}$ and
\begin{align}
\ln R_{t,n}(h) &= \ln R_{t,n}+\ln R_{t-1,n}+\cdots+\ln R_{t-h+1,n} \label{eq:sum}\\
&= p_{t,n}-p_{t-h,n}. \label{eq:cancel}
\end{align}

\begin{keybox}[: two insights]
\begin{enumerate}
  \item \textbf{Eq.~\eqref{eq:cancel}}: only the first and last prices matter.
    The 100-year return is identical using two prices, annual, monthly or daily
    data. $\Rightarrow$ extra observations do \emph{not} help estimate the mean
    (they \emph{do} help for the variance). Estimates of expected returns are
    very imprecise; estimates of variances are much more precise.
  \item \textbf{Eq.~\eqref{eq:sum}}: multiperiod log return $=$ \emph{sum} of
    one-period log returns $\Rightarrow$ easy statistics (sums, not products).
\end{enumerate}
\end{keybox}

\begin{plainbox}
Two kinds of returns, two jobs. Simple (gross) returns are what you earn: a
portfolio's simple return is the weighted average of the assets' simple returns.
Log returns are convenient over time: the 10-year log return is just the sum of
the yearly ones. Rule: across assets use simple returns; over time use log returns.
\end{plainbox}
\paragraph{Disadvantage of log returns.} Gross portfolio returns aggregate
linearly, log portfolio returns do not:
\[
R_{t,p}=\sum_n w_{t,n}R_{t,n}, \qquad
\ln R_{t,p}\neq\sum_n w_{t,n}\ln R_{t,n},
\]
but over short intervals $\ln R_{t,p}\approx\sum_n w_{t,n}\ln R_{t,n}$.
\begin{intuition}
\begin{itemize}
  \item \textbf{Cross-section} (allocating across assets) $\Rightarrow$ use
    gross returns $R_{t,p}$ or gross excess returns $R_{t,p}-R_{t,0}$.
  \item \textbf{Time series} $\Rightarrow$ use log returns $\ln R_{t,p}$ or log
    excess returns.
  \item This course focuses on the \textbf{cross-section}.
\end{itemize}
\end{intuition}

\subsection*{Python exercise 3: sampling more often does not help estimate the mean}
2{,}000 samples of 10 years of daily log returns, true $\mu=8\%$, $\sigma=20\%$ p.a.
\begin{lstlisting}
import numpy as np
import pandas as pd
rng = np.random.default_rng(3)
M = 2000                 # number of simulated 10-year samples
years, d = 10, 252       # years per sample, days per year
mu, sigma = 0.08, 0.20   # true annual mean and volatility
r = rng.normal(mu/d, sigma/np.sqrt(d), size=(M, years*d))
# k-day log returns: sum k consecutive daily log returns
k = {"daily": 1, "monthly": 21, "annual": 252}
rk = {f: r.reshape(M, -1, k[f]).sum(axis=2) for f in k}
n = {f: d // k[f] for f in k}   # observations per year

# Annualized estimates in each of the M samples
mean_hat = {f: rk[f].mean(axis=1) * n[f] for f in k}
vol_hat = {f: rk[f].std(axis=1, ddof=1) * np.sqrt(n[f]) for f in k}
table = pd.DataFrame({
    "T": [rk[f].shape[1] for f in k],
    "SE of mean": [mean_hat[f].std() for f in k],
    "SE of vol":  [vol_hat[f].std() for f in k]}, index=list(k))
print(table.round(4))
\end{lstlisting}
\begin{lstlisting}[style=out]
            T  SE of mean  SE of vol
daily    2520      0.0632     0.0028
monthly   120      0.0632     0.0126
annual     10      0.0632     0.0460
\end{lstlisting}
\begin{itemize}
  \item SE of the mean $=6.32\%$ at every frequency $=\sigma/\sqrt{\text{years}}=0.20/\sqrt{10}$.
  \item SE of volatility falls from 4.6\% (annual) to 0.28\% (daily).
  \item Even with 10 years of daily data, a 95\% CI for the mean is
    $8\%\pm12.6\%$: we cannot tell if the mean is positive.
\end{itemize}


\begin{workedex}[: how many years do we need?]
The standard error of an annual mean is $\sigma/\sqrt{\text{years}}$. For a single
stock ($\sigma=30\%$) and for the market ($\sigma=16\%$):
\begin{center}
\begin{tabular}{@{}lcccc@{}}
\toprule
Years of data & 5 & 10 & 25 & 100 \\
\midrule
SE, single stock ($\sigma=30\%$) & 13.4\% & 9.5\% & 6.0\% & 3.0\% \\
SE, market ($\sigma=16\%$) & 7.2\% & 5.1\% & 3.2\% & 1.6\% \\
\bottomrule
\end{tabular}
\end{center}
A 95\% confidence interval is about $\pm2$ SE. Even with 100 years of data, the
market's mean is known only to within $\pm3.2\%$, while the equity premium itself
is about 5--7\%. For volatility, by contrast, a few years of daily data are enough.
\end{workedex}
\begin{addnote}
\textbf{The math behind it.} With $T$ years sampled $n$ times per year, each
period has mean $\mu/n$ and variance $\sigma^2/n$. The annualised mean
estimator $n\cdot\frac{1}{nT}\sum r_i$ has variance
$n^2\cdot\frac{1}{(nT)^2}\cdot nT\cdot\frac{\sigma^2}{n}=\frac{\sigma^2}{T}$:
independent of $n$. For volatility, $\text{SE}(\hat\sigma)\approx
\sigma/\sqrt{2nT}$, which shrinks as $n$ grows. Check:
$0.20/\sqrt{2\cdot 2520}=0.0028$. \checkmark
\end{addnote}

\subsection*{Python exercise 4: simple and log returns}
10 years of monthly gross returns, two assets, mean 0.8\%, vol 5\% per month.
\begin{lstlisting}
import numpy as np
rng = np.random.default_rng(4)
T = 120   # 10 years of monthly returns on 2 assets
R = 1 + rng.normal(0.008, 0.05, size=(T, 2))  # gross
r = np.log(R)                                 # log
# 1. Product of gross returns = exp(sum of log returns)
print("prod(R):", R[:, 0].prod().round(6))
print("exp(sum r):", np.exp(r[:, 0].sum()).round(6))
# 2. Annualized arithmetic mean vs geometric mean
arith = 12 * (R[:, 0] - 1).mean()
geom = R[:, 0].prod()**(12/T) - 1
half_var = 12 * (R[:, 0] - 1).var() / 2
print("arithmetic:", arith.round(4), "geometric:", geom.round(4),
      "sigma^2/2:", half_var.round(4))
# 3. Portfolio log return vs weighted sum of log returns
w = np.array([0.6, 0.4])
gap_m = np.abs(np.log(R @ w) - r @ w).mean()
Ra = R.reshape(10, 12, 2).prod(axis=1)  # annual gross
gap_a = np.abs(np.log(Ra @ w) - np.log(Ra) @ w).mean()
print("mean gap, monthly:", gap_m.round(5), "annual:", gap_a.round(5))
\end{lstlisting}
\begin{lstlisting}[style=out]
prod(R): 2.079249
exp(sum r): 2.079249
arithmetic: 0.0892 geometric: 0.0759 sigma^2/2: 0.0157
mean gap, monthly: 0.00066 annual: 0.00742
\end{lstlisting}
\begin{enumerate}
  \item Product of gross returns $=\exp(\sum$ log returns$)$ exactly.
  \item Arithmetic (8.92\%) exceeds geometric (7.59\%) by 1.33\%, close to $\sigma^2/2=1.57\%$.
  \item Log-aggregation gap: 0.07\% monthly vs.\ 0.74\% annual (about 11 times larger).
\end{enumerate}

\subsection{Returns with dividends}
\[
R_{t,n}=\frac{P_{t,n}+D_{t,n}}{P_{t-1,n}},\qquad
\ln R_{t,n}=\ln(P_{t,n}+D_{t,n})-\ln P_{t-1,n}.
\]

\subsection{Distribution of returns}
Returns are random variables, characterised by their \textbf{joint},
\textbf{conditional} or \textbf{unconditional} distribution.

\paragraph{Unconditional: IID lognormal.} Convenient assumption: gross returns
IID lognormal $\Leftrightarrow$ $\ln R_{t,n}\sim\mathcal N(\mu_n,\sigma_n^2)$.
Lognormal has lower bound 0 $\Rightarrow$ respects \textbf{limited liability};
$\ln R$ ranges over $(-\infty,+\infty)$.

\begin{keybox}[: lognormal moments]
If $\ln X\sim\mathcal N(\mu,\sigma^2)$: \quad
$\E[X^m]=\exp\!\big(m\mu+\tfrac12 m^2\sigma^2\big)$. Hence
\begin{align*}
\E[R_n] &= \exp\!\big(\mu_n+\tfrac12\sigma_n^2\big)\\
\E[R_n^2] &= \exp\!\big(2\mu_n+2\sigma_n^2\big)\\
\V[R_n] &= \E[R_n^2]-(\E[R_n])^2=\exp(2\mu_n+\sigma_n^2)\big(\exp(\sigma_n^2)-1\big).
\end{align*}
\end{keybox}
\begin{intuition}
$\E[R]=e^{\mu+\sigma^2/2}>e^{\mu}$: the mean of the gross return exceeds the
exponential of the mean log return by the factor $e^{\sigma^2/2}$ (Jensen's
inequality). Taking logs: $\ln\E[R]=\mu+\sigma^2/2$ --- the same
$\sigma^2/2$ gap between arithmetic and geometric means seen in Exercise 4.
\end{intuition}

\subsection{Are returns Normal? Skewness and kurtosis}
At short horizons returns are \emph{not} Normal: weak skewness, strong
\textbf{excess kurtosis} (fat tails).
\[
S[X]=\E\!\left[\frac{(X-\mu)^3}{\sigma^3}\right],\qquad
K[X]=\E\!\left[\frac{(X-\mu)^4}{\sigma^4}\right].
\]
Normal: $S=0$, $K=3$; excess kurtosis $=K-3$. Sample estimators (dividing by $T$):
\[
\hat S=\frac{1}{\hat\sigma^3}\frac1T\sum_{t=1}^T(x_t-\hat\mu)^3,\quad
\hat K=\frac{1}{\hat\sigma^4}\frac1T\sum_{t=1}^T(x_t-\hat\mu)^4,\quad
\hat\mu=\frac1T\sum x_t,\quad \hat\sigma^2=\frac1T\sum(x_t-\hat\mu)^2.
\]
Under Normality, for large $T$: $\hat S\sim\mathcal N(0,6/T)$,
$\hat K\sim\mathcal N(3,24/T)$.

\begin{addnote}
\textbf{Jarque--Bera test} combines both:
\[
JB=T\left[\frac{\hat S^2}{6}+\frac{(\hat K-3)^2}{24}\right]\ \overset{H_0}{\sim}\ \chi^2_2 .
\]
Large $JB$ (small $p$-value) $\Rightarrow$ reject Normality. It is just the sum
of the two squared $z$-statistics implied by the asymptotic distributions above.
\end{addnote}

\begin{plainbox}
Real daily returns have many more ``crash'' and ``boom'' days than a bell curve
predicts (fat tails). Kurtosis measures this: 3 for a Normal, often 5--10 for
daily stock returns. Over a month, many daily shocks add up and the distribution
looks more Normal, but not completely.
\end{plainbox}
\subsection*{Python exercise 5: testing whether returns are Normal}
10{,}500 daily returns from Student-$t$ with 5 d.o.f.\ (fat tails), scaled to
1\% daily vol; monthly $=$ sum of 21 daily log returns.
\begin{lstlisting}
import numpy as np
import pandas as pd
from scipy import stats
rng = np.random.default_rng(5)
T, nu = 21*500, 5    # 500 months of days; t with 5 df
z = rng.standard_t(nu, size=T) / np.sqrt(nu/(nu-2))
daily = 0.01 * z     # 1% daily volatility, fat tails
monthly = daily.reshape(-1, 21).sum(axis=1)
x = {"daily": daily, "monthly": monthly}

# Moves larger than 4 standard deviations
n4 = {f: np.sum(np.abs(x[f] - x[f].mean()) > 4*x[f].std()) for f in x}
table = pd.DataFrame({
    "T": [len(x[f]) for f in x],
    "skew": [stats.skew(x[f]) for f in x],
    "kurt": [stats.kurtosis(x[f], fisher=False) for f in x],
    "JB p-value": [stats.jarque_bera(x[f]).pvalue for f in x],
    "4-sigma": [n4[f] for f in x],
    "if Normal": [len(x[f])*2*stats.norm.sf(4) for f in x]},
    index=list(x))
print(table.round(3))
\end{lstlisting}
\begin{lstlisting}[style=out]
             T    skew   kurt  JB p-value  4-sigma  if Normal
daily    10500  -0.104  7.256       0.000       40      0.665
monthly    500   0.076  2.932       0.751        0      0.032
\end{lstlisting}
\begin{itemize}
  \item Daily: kurtosis 7.3 (SE under Normality $\sqrt{24/T}\approx0.05$);
    40 four-sigma days vs.\ fewer than one expected.
  \item Monthly: kurtosis 2.9, JB $p=0.75$ $\Rightarrow$ cannot reject Normality:
    sums of many fat-tailed draws are close to Normal (CLT).
  \item Caveat: in real data volatility is time-varying, so monthly returns
    are also fat-tailed (less than daily).
\end{itemize}
\begin{addnote}
\texttt{np.sqrt(nu/(nu-2))}: the variance of a Student-$t$ with $\nu$ d.o.f.\ is
$\nu/(\nu-2)$, so dividing standardises it to unit variance.
\texttt{stats.kurtosis(..., fisher=False)} returns $K$ (Normal $=3$);
the default \texttt{fisher=True} returns \emph{excess} kurtosis (Normal $=0$).
\end{addnote}

% ---------------------------------------------------------------------
\section{Obtaining financial and macroeconomic data}
Three kinds of data: (1) stock prices, (2) stock characteristics,
(3) macroeconomic data. Code adapted from an earlier edition of Tidy Finance
(the current edition uses its own \texttt{tidyfinance} package and
\texttt{polars}).

\subsection{Stock-price data (Yahoo!\ Finance)}
\begin{lstlisting}
import pandas as pd
import numpy as np
import yfinance as yf

prices = (yf.download(tickers="AAPL", start="2000-01-01", end="2024-12-31",
                      auto_adjust=False,        # keep "Adj Close"
                      multi_level_index=False)  # one level of column names
  .reset_index()
  .assign(symbol="AAPL")
  .rename(columns={"Date": "date", "Open": "open", "High": "high",
                   "Low": "low", "Close": "close",
                   "Adj Close": "adjusted", "Volume": "volume"}))
prices.head()
\end{lstlisting}
\begin{lstlisting}[style=out]
   date        open      high      low       close     adjusted  volume     symbol
0  2000-01-03  0.936384  1.004464  0.907924  0.999442  0.848323  535796800  AAPL
1  2000-01-04  0.966518  0.987723  0.903460  0.915179  0.776801  512377600  AAPL
\end{lstlisting}

\begin{keybox}[: adjusted prices]
Adjusted prices correct for corporate actions (stock splits, dividends) that
change quoted prices but have no direct impact on investors. \textbf{Use
adjusted prices} to compute returns. Yahoo rescales adjusted prices after every
dividend/split, so values downloaded today differ from older downloads.
\end{keybox}
Apple splits: 2:1 (1987, 2000, 2005), 7:1 (2014), 4:1 (2020). Opening price
on 2000-01-04: \$0.9665 (split-adjusted); on 2023-01-04: \$126.98.

\subsection*{Python exercise 6: adjusting prices for splits and dividends}
250 days, start at \$400, 4-for-1 split on day 150, \$1 dividend on day 200.
Adjustment factor on each event date multiplies \emph{all earlier} prices:
split $\to 1/4$; dividend $\to P_t/(P_t+D)$.
\begin{lstlisting}
import numpy as np
import pandas as pd
rng = np.random.default_rng(6)
n = 250                                       # trading days
true_ret = rng.normal(0.0004, 0.015, size=n)  # total returns
true_ret[0] = 0.0
split_day, ratio = 150, 4                     # 4-for-1 split
div_day, div = 200, 1.00                      # $1 dividend per (new) share
G = np.cumprod(1 + true_ret)
day = np.arange(n)
shares = np.where(day >= split_day, ratio, 1)
close = (400*G/shares - div*(G/G[div_day])*(day >= div_day))

# Adjustment factor on each event date; multiplies all prices BEFORE it
f = np.ones(n)
f[split_day] = 1 / ratio
f[div_day] = close[div_day] / (close[div_day] + div)
adj = np.append(np.cumprod(f[::-1])[::-1][1:], 1.0)
df = pd.DataFrame({"close": close, "adjusted": close*adj})
df["raw_ret"] = df["close"].pct_change()
df["adj_ret"] = df["adjusted"].pct_change()
print(df.loc[[split_day, div_day]].round(4))
err = np.abs(df["adj_ret"].values[1:] - true_ret[1:]).max()
print("Largest error of adjusted returns:", err)
\end{lstlisting}
\begin{lstlisting}[style=out]
        close  adjusted  raw_ret  adj_ret
150  102.2968  101.4969  -0.7536  -0.0143
200  126.8846  126.8846  -0.0280  -0.0203
Largest error of adjusted returns: 3.8e-16
\end{lstlisting}
\begin{itemize}
  \item Split date: raw ``return'' $-75.4\%$; adjusted $-1.43\%$ (true).
  \item Dividend date: raw $-2.80\%$ misses the \$1 dividend; adjusted
    $-2.03\%$ is the true total return.
  \item Yahoo uses the similar factor $1-D/P_{t-1}$: very close, not identical,
    to total returns.
\end{itemize}
\begin{addnote}
\texttt{np.cumprod(f[::-1])[::-1]} is a \emph{reverse cumulative product}:
entry $t$ is the product of all factors from $t$ to the end. Shifting by one
(\texttt{[1:]}, append 1) means each price is multiplied by the factors of all
events that happen \emph{after} it --- exactly the backward adjustment.
\end{addnote}

\subsection{From prices to returns}
\[ r_t=\frac{P_t}{P_{t-1}}-1. \]
\begin{lstlisting}
returns = (prices
  .sort_values("date")
  .assign(ret=lambda x: x["adjusted"].pct_change())
  .get(["symbol", "date", "ret"]))
\end{lstlisting}
First return is \texttt{NaN}; e.g.\ AAPL 2000-01-04: $-8.43\%$.

\paragraph{Many stocks.}
\begin{enumerate}
  \item Download the tickers of the index constituents (e.g.\ DJIA).
  \item Pass the list to \texttt{yf.download(tickers=[...])}.
  \item Convert to long format with \texttt{.stack(level="Ticker")}.
  \item Compute returns with \texttt{groupby("symbol")}.
\end{enumerate} Learn to do it for S\&P 500, Nasdaq, DAX
(see Brugi\`ere 2020, Sec.~9.5).

\subsection*{Python exercise 7: survivorship bias}
Using \emph{today's} constituents to compute past returns keeps only survivors.
1{,}000 stocks, 120 months, true mean 0.8\%, vol 12\% per month; delisted once
price $<30\%$ of start (the return in that month is included).
\begin{lstlisting}
import numpy as np
rng = np.random.default_rng(7)
N, T = 1000, 120
ret = rng.normal(0.008, 0.12, size=(N, T))
price = np.cumprod(1 + ret, axis=1)
dead = np.maximum.accumulate(price < 0.3, axis=1)   # once dead, always dead
listed = np.hstack([np.ones((N, 1), bool), ~dead[:, :-1]])
survivor = ~dead[:, -1]
ret_listed = np.where(listed, ret, np.nan)
all_firms = np.nanmean(ret_listed, axis=0).mean()
survivors = ret[survivor].mean()
print("Share of firms surviving:", survivor.mean().round(3))
print("Mean return, all listed firms:", all_firms.round(4))
print("Mean return, survivors only:", survivors.round(4))
\end{lstlisting}
\begin{lstlisting}[style=out]
Share of firms surviving:      0.663
Mean return, all listed firms: 0.0079   # ~ true 0.008
Mean return, survivors only:   0.0132   # survivorship bias
\end{lstlisting}
\begin{keybox}[: survivorship bias]
Survivors earn 1.32\%/month vs.\ true 0.8\%: overstatement of about 6\% per
year. For backtests use the constituents \textbf{at each past date} (e.g.\
CRSP), not today's.
\end{keybox}

\subsection{Stock-characteristics data}
Sample moments (simplest estimators of $\E[\R]$ and $\V[\R]$):
\[
\hat\mu_{R_n}=\frac1T\sum_{t=1}^T R_{t,n},\quad
\hat\sigma_{R_{nn}}=\frac1T\sum_{t=1}^T(R_{t,n}-\hat\mu_{R_n})^2,\quad
\hat\sigma_{R_{nm}}=\frac1T\sum_{t=1}^T(R_{t,n}-\hat\mu_{R_n})(R_{t,m}-\hat\mu_{R_m}).
\]
\textbf{Note:} pandas \texttt{.var()} and \texttt{.cov()} divide by $T-1$, not $T$
(\texttt{ddof=1} by default; NumPy \texttt{np.var} uses \texttt{ddof=0}).

Alternatively, use firm \textbf{characteristics} --- size, value (book-to-market),
profitability, investment --- to estimate $\E[R_n]$ and $\V[R_n]$. The
source is \textbf{Ken French's data library}, which we read with the package
\texttt{pandas-datareader}.

\begin{lstlisting}
import pandas_datareader as pdr
start_date = "1960-01-01"
end_date = "2022-12-31"

factors_ff3_monthly_raw = pdr.DataReader(
  name="F-F_Research_Data_Factors", data_source="famafrench",
  start=start_date, end=end_date)[0]
factors_ff3_monthly = (factors_ff3_monthly_raw
  .divide(100)                          # percent -> decimal
  .reset_index(names="month")
  .assign(month=lambda x: pd.to_datetime(x["month"].astype(str)))
  .rename(str.lower, axis="columns")
  .rename(columns={"mkt-rf": "mkt_excess"}))
\end{lstlisting}
\begin{itemize}
  \item FF3: \texttt{mkt\_excess}, \texttt{smb} (size), \texttt{hml} (value), \texttt{rf}.
  \item FF5 (\texttt{"F-F\_Research\_Data\_5\_Factors\_2x3"}): adds \texttt{rmw}
    (profitability) and \texttt{cma} (investment). Same code, different name.
  \item Also available: daily FF factors, 10 industry portfolios, $q$-factors
    (Hou, Xue, Zhang 2015).
\end{itemize}

\paragraph{Other sources.} PyAnomaly, getFamaFrenchFactors; \textbf{Open Source
Asset Pricing} (330+ predictors; Chen--Zimmermann 2022); \texttt{investpy}
(Investing.com, non-US stocks, funds, ETFs, FX, bonds, crypto); Gu--Kelly--Xiu
(2020) for ML; Open Source Bond Asset Pricing (Dickerson--Mueller--Robotti 2023);
PyCleanTrace for TRACE bond data (Dick-Nielsen--Poulsen 2019).

\paragraph{EDGAR.} SEC filing database, free. Python APIs: SEC Edgar Downloader,
python-edgar, OpenEDGAR, EDGAR Tools, SEC EDGAR API; commercial sec-api.io.
Non-US equivalents: France --- Registre du commerce et des soci\'et\'es; UK ---
Companies House; Canada --- SEDAR+; others for Germany, Ireland, Italy, Switzerland.

\subsection{Macroeconomic data}
Macro variables are used to predict the market equity premium.
\textbf{Welch and Goyal (2008)} re-examine many predictors:
\begin{center}
\begin{tabularx}{\textwidth}{@{}lX@{}}
\toprule
Code & Definition \\
\midrule
\texttt{dp} & log dividend--price ratio: $\ln D_{12}-\ln P$ \\
\texttt{dy} & log dividend yield: $\ln D_{12}-\ln P_{t-1}$ \\
\texttt{ep} & log earnings--price ratio: $\ln E_{12}-\ln P$ \\
\texttt{de} & log dividend payout: $\ln D_{12}-\ln E_{12}$ \\
\texttt{tbl}, \texttt{lty}, \texttt{ltr} & 3-month T-bill rate, long-term yield, long-term govt bond return \\
\texttt{tms} & term spread: \texttt{lty} $-$ \texttt{tbl} \\
\texttt{dfy} & default spread: BAA $-$ AAA yields \\
\texttt{svar}, \texttt{infl}, \texttt{ntis}, \texttt{b/m} & stock variance, inflation, net equity expansion, book-to-market \\
\texttt{rp\_div} & target: next-month log excess return incl.\ dividends \\
\bottomrule
\end{tabularx}
\end{center}
\begin{lstlisting}
sheet_id = "1bM7vCWd3WOt95Sf9qjLPZjoiafgF_8EG"
sheet_name = "macro_predictors.xlsx"
macro_predictors_link = ("https://docs.google.com/spreadsheets/d/" + sheet_id
                         + "/gviz/tq?tqx=out:csv&sheet=" + sheet_name)
macro_predictors = (
  pd.read_csv(macro_predictors_link, thousands=",")
  .assign(
    month=lambda x: pd.to_datetime(x["yyyymm"], format="%Y%m"),
    IndexDiv=lambda x: x["Index"] + x["D12"],
    logret=lambda x: np.log(x["IndexDiv"]) - np.log(x["IndexDiv"].shift(1)),
    Rfree=lambda x: np.log(x["Rfree"] + 1),
    rp_div=lambda x: (x["logret"] - x["Rfree"]).shift(-1),
    dp=lambda x: np.log(x["D12"]) - np.log(x["Index"]),
    dy=lambda x: np.log(x["D12"]) - np.log(x["Index"].shift(1)),
    ep=lambda x: np.log(x["E12"]) - np.log(x["Index"]),
    de=lambda x: np.log(x["D12"]) - np.log(x["E12"]),
    tms=lambda x: x["lty"] - x["tbl"],
    dfy=lambda x: x["BAA"] - x["AAA"])
  .get(["month", "rp_div", "dp", "dy", "ep", "de", "svar", "b/m",
        "ntis", "tbl", "lty", "ltr", "tms", "dfy", "infl"])
  .query("month >= @start_date and month <= @end_date")
  .dropna())
\end{lstlisting}
\begin{addnote}
\texttt{.shift(-1)} on \texttt{rp\_div} aligns the \emph{next} month's excess
return with \emph{this} month's predictors: a regression
$rp_{t+1}=a+b\,x_t+\varepsilon_{t+1}$ is then a genuine predictive regression
(no look-ahead bias).
\end{addnote}

\paragraph{FRED} (St.\ Louis Fed, 817{,}000 series):
\begin{lstlisting}
cpi_monthly = (pdr.DataReader(name="CPIAUCNS", data_source="fred",
                              start=start_date, end=end_date)
  .reset_index(names="month")
  .rename(columns={"CPIAUCNS": "cpi"})
  .assign(cpi=lambda x: x["cpi"] / x["cpi"].iloc[-1]))  # rebase to last month = 1
\end{lstlisting}
Any series: just change the key (e.g.\ gold-ore PPI: \texttt{PCU2122212122210}).

\paragraph{WRDS / CRSP / Compustat} (optional, not covered in class): SQL
(PostgreSQL) database; Tidy Finance explains access and storage in SQLite.

% ---------------------------------------------------------------------
\section{Takeaways of Chapter 1}
\begin{enumerate}
  \item Key formula: $\w=\frac1\gamma\V[\R]^{-1}(\E[\R]-\Rf\one)$; the whole
    course is about estimating $\E[\R]$ and $\V[\R]$ well.
  \item $\E[\R]$: imprecise, because only the first and last prices matter;
    SE $=\sigma/\sqrt{\text{years}}$ regardless of frequency.
  \item $\V[\R]$: $N(N+1)/2$ parameters; the sample matrix is ill-conditioned
    when $N$ is large relative to $T$; $\V^{-1}$ magnifies errors in means.
  \item Performance must be judged \textbf{out of sample}: $R_{t+1,p}=\w_t\T\R_{t+1}$.
  \item Tidy data; pivot to wide only for matrix algebra; compute returns
    \emph{within} each asset (\texttt{groupby}).
  \item Gross returns aggregate across assets; log returns aggregate over time.
    Arithmetic $-$ geometric mean $\approx\sigma^2/2$.
  \item Daily returns are fat-tailed; monthly closer to Normal.
  \item Use adjusted prices; beware survivorship bias.
\end{enumerate}


\begin{mistakes}
\begin{enumerate}
  \item Computing returns with \texttt{pct\_change()} on a tidy column without
    \texttt{groupby("symbol")}: the first ``return'' of each stock is computed
    from the previous stock's last price.
  \item Using the raw \texttt{Close} price instead of the \texttt{Adj Close}: splits
    show up as $-50\%$ or $-75\%$ ``returns'' and dividends are missed.
  \item Annualising volatility by multiplying by 12 instead of $\sqrt{12}$ (only
    means and variances scale with time).
  \item Forgetting that \texttt{pandas .std()} uses $T-1$ and \texttt{np.std} uses $T$.
  \item Averaging log returns across assets to get a portfolio return (only simple
    returns aggregate across assets).
  \item Using today's index constituents to study the past (survivorship bias).
  \item Using \texttt{np.linalg.inv} on a near-singular matrix without noticing:
    always look at the condition number or use \texttt{np.linalg.solve}.
\end{enumerate}
\end{mistakes}
\section{Self-test questions (quiz/exam practice)}
\begin{enumerate}
  \item If risk aversion doubles, what happens to $w_1$? \emph{(It halves.)}
  \item How many parameters in $\V[\R]$ for $N=50$? \emph{(1275.)}
  \item Going from monthly to daily data over the same 10 years, does the SE of
    the mean fall? \emph{(No; only the SE of volatility falls.)}
  \item Why is $\ln R_{t,p}\neq\sum_n w_{t,n}\ln R_{t,n}$? When is the
    approximation good? \emph{(Log is non-linear; good over short intervals.)}
  \item Excess gross vs.\ excess net returns? \emph{(Identical.)}
  \item Kurtosis of a Normal? SE of $\hat K$ with $T=10{,}000$?
    \emph{(3; $\sqrt{24/10000}\approx0.049$.)}
  \item A stock has a 2-for-1 split. What is the raw return on that day if
    the true return is 0? \emph{($-50\%$.)}
  \item Write by hand the Python code to compute returns from a tidy
    DataFrame with columns \texttt{symbol, date, adjusted}.
\begin{lstlisting}
df = df.sort_values(["symbol", "date"])
df["ret"] = df.groupby("symbol")["adjusted"].pct_change()
\end{lstlisting}
  \item Write the code for GMV weights given a covariance matrix \texttt{V}.
\begin{lstlisting}
ones = np.ones(V.shape[0])
w = np.linalg.solve(V, ones)
w = w / w.sum()
\end{lstlisting}
  \item Write the code for MV weights given \texttt{mu}, \texttt{V},
    \texttt{rf}, \texttt{gamma}.
\begin{lstlisting}
w = np.linalg.solve(V, mu - rf * np.ones(len(mu))) / gamma
\end{lstlisting}
\end{enumerate}

\section{To do for next class}
\begin{itemize}
  \item Read the \emph{Financial Data} section of Tidy Finance with Python.
  \item Start Assignment \#1 (covers Ch.\,1 and Ch.\,2; due 11 Nov, 22:00 CET).
\end{itemize}

% =====================================================================
%  CHAPTER 2
% =====================================================================
\chapter{Performance Measurement}
\label{ch:2}

\section{What we want to do in Chapter 2}
\begin{itemize}
  \item Learn how to measure the performance of a portfolio strategy and why it
    must be \textbf{out of sample}: in-sample returns suffer from
    \textbf{look-ahead bias}.
  \item Study many performance metrics. The key one, used for the rest of the
    course, is the \textbf{Sharpe ratio}.
  \item Test whether the difference between the Sharpe ratios of two portfolios
    is \textbf{statistically significant} (delta method and bootstrap).
\end{itemize}

% ---------------------------------------------------------------------
\begin{plainbox}
Chapter 2 is about \emph{how to keep score}. Before comparing strategies we must
agree on (1) how to compute returns honestly (out of sample), (2) which numbers
to look at (mainly the Sharpe ratio, plus risk measures and turnover), and (3)
how to tell whether one strategy is \emph{really} better or just lucky (statistical tests).
\end{plainbox}
\section{Static, myopic, and dynamic portfolios}
\begin{center}
\begin{tabularx}{\textwidth}{@{}lXX@{}}
\toprule
Type & Definition & Example (25-year-old saving for retirement at 65) \\
\midrule
\textbf{Static} & Form the portfolio at $t=0$ and hold it until $T$. & Chooses a portfolio once and holds it for 40 years. \\
\textbf{Myopic} & Rebalance at each date, optimising \emph{only one period ahead} (ignore future rebalancing). & Rebalances every year, optimising only for the next year. \\
\textbf{Dynamic} & Rebalance at each date, \emph{anticipating} future rebalancing. & Rebalances every year, knowing she will rebalance again next year. \\
\bottomrule
\end{tabularx}
\end{center}

\begin{keybox}[: the course studies myopic portfolios]
\begin{itemize}
  \item \textbf{Not static}: static portfolios are suboptimal relative to myopic (and dynamic) ones.
  \item \textbf{Not dynamic}: they require estimating many more parameters, so
    they perform poorly out of sample (see Merton 1971, 1990; Campbell--Viceira 2002).
  \item Ignoring estimation error, myopic is suboptimal relative to dynamic,
    but in practice estimation error dominates.
\end{itemize}
\end{keybox}

% ---------------------------------------------------------------------
\section{Backtesting: in-sample vs.\ out-of-sample performance}

\paragraph{Notation.} $R_{t,p}$ is the gross return at date $t$ of portfolio $p$
(convention: $R_{\text{date},\text{asset}}$). A myopic portfolio $\w_{0,p}$
chosen at $t=0$ is evaluated by its return $R_{1,p}$.

\begin{keybox}[: weights must depend only on past information]
$\w_{0,p}$ may depend on information at $t=0,-1,-2,\dots$ but \textbf{never} on
information from $t=1,2,\dots$
\begin{itemize}
  \item If $\w_{0,p}$ uses only past information $\Rightarrow$ $R_{1,p}$ is an
    \textbf{out-of-sample} return.
  \item If $\w_{0,p}$ also uses future information $\Rightarrow$ $R_{1,p}$ is an
    \textbf{in-sample} return (not realistic: look-ahead bias).
\end{itemize}
The sequence of weights $\w_p=\{\w_{0,p},\w_{1,p},\dots,\w_{T-1,p}\}$ generates
the out-of-sample returns $\bm R_p=\{R_{1,p},R_{2,p},\dots,R_{T,p}\}$.
\end{keybox}

\begin{center}
\begin{tabular}{c|c|c|c|c|c}
$t=0$ & $t=1$ & $t=2$ & $\cdots$ & $t=T-1$ & $t=T$\\\hline
$\w_{0,p}$ & $R_{1,p}$, $\w_{1,p}$ & $R_{2,p}$, $\w_{2,p}$ & $\cdots$ & $R_{T-1,p}$, $\w_{T-1,p}$ & $R_{T,p}$
\end{tabular}
\end{center}

\subsection{Rolling vs.\ expanding window}
\begin{itemize}
  \item \textbf{Rolling window}: fixed-length lookback. E.g.\ a 5-year window: the
    decision in 2020 uses 2015--2019, in 2021 uses 2016--2020, in 2022 uses 2017--2021.
    Use it if the distant past is \emph{not} relevant or is different from the recent past.
  \item \textbf{Expanding window}: fixed start, growing length. The decision in
    2020 uses 2015--2019, in 2021 uses 2015--2020, in 2022 uses 2015--2021.
    Use it if the distant past is \emph{similar} to the recent past.
  \item \textbf{Exponentially declining weights} on past data: a compromise.
\end{itemize}
\begin{lstlisting}
# pandas: rolling, expanding, exponentially weighted (EWMA) estimates
mu_roll = ret.rolling(window=60).mean()        # last 60 months
mu_exp  = ret.expanding(min_periods=60).mean() # all data since the start
vol_ewm = ret.ewm(halflife=12).std()           # weights decay with half-life 12 months
# IMPORTANT: to use them as forecasts at t for t+1, shift by one period
mu_forecast = ret.rolling(60).mean().shift(1)
\end{lstlisting}
\begin{addnote}
\texttt{rolling(60).mean()} at date $t$ \emph{includes} the return at $t$.
To form weights at $t-1$ for the return at $t$ without look-ahead bias, use
\texttt{.shift(1)}, or loop over \texttt{R.iloc[t-60:t]} as in the backtest
skeleton of Notebook 00.
\end{addnote}

% ---------------------------------------------------------------------
\section{Measures of portfolio performance}
Three groups: (1) mean returns; (2) risk of returns; (3) mean and risk combined.

\subsection{Measures of mean returns}
\begin{enumerate}
  \item \textbf{Total mean return} of the portfolio.
  \item \textbf{Systematic component}: the reward for bearing systematic risk, from a factor model.
  \item \textbf{Alpha}: the part of the mean return unexplained by the factor
    model. If the model is correct, $\alpha=0$. (In Chapter 8: part of alpha
    is a reward for unsystematic risk.)
  \item \textbf{Outperformance frequency}: the average fraction of times the
    portfolio has a higher cumulative return than the benchmark within twelve
    months from the start of each such period.
\end{enumerate}

\paragraph{Return factors} (long--short portfolios):
$R_{mkt}-R_f$ (market); $R_{smb}$ small minus big (size); $R_{hml}$ high minus low
book-to-market (value); $R_{umd}$ up minus down (momentum); $R_{rmw}$ robust
minus weak profitability; $R_{cma}$ conservative minus aggressive investment.

\begin{keybox}[: factor models of expected returns]
\begin{align*}
\text{CAPM (Sharpe 1964):}\quad & E[R_p]-R_f=\beta_{p,mkt}E[R_{mkt}-R_f]\\
\text{FF3 (1993):}\quad & \ldots+\beta_{p,smb}E[R_{smb}]+\beta_{p,hml}E[R_{hml}]\\
\text{Carhart 4F (1997):}\quad & \text{FF3}+\beta_{p,umd}E[R_{umd}]\\
\text{FF5 (2015):}\quad & \text{FF3}+\beta_{p,rmw}E[R_{rmw}]+\beta_{p,cma}E[R_{cma}]\\
\text{FF6 (2018):}\quad & \text{FF5}+\beta_{p,umd}E[R_{umd}]
\end{align*}
Also: $q$-factor models (Hou, Xue, Zhang 2015, 2017; Hou, Mo, Xue, Zhang 2019);
$K$ statistical (\emph{latent}) factors from PCA of the covariance matrix
(Connor--Goldberg--Korajczyk 2010; Kozak--Nagel--Santosh 2018, 2020).
Dello-Preite, Uppal, Zaffaroni, Zviadadze (2025): once mean returns are adjusted
for \emph{unsystematic} risk, the only systematic factor needed is the market
(or market and size).
\end{keybox}

\begin{addnote}
\textbf{How betas and alphas are estimated in practice:} a time-series OLS
regression of the portfolio's excess return on the factors,
\[
R_{t,p}-R_{t,f}=\alpha_p+\beta_{p,mkt}(R_{t,mkt}-R_{t,f})+\beta_{p,smb}R_{t,smb}+\beta_{p,hml}R_{t,hml}+\varepsilon_{t,p}.
\]
The intercept is the alpha; its $t$-statistic tests $\alpha_p=0$.
\begin{lstlisting}
import statsmodels.api as sm
X = sm.add_constant(ff3[["mkt_excess", "smb", "hml"]])
fit = sm.OLS(port_excess, X).fit()
alpha, betas = fit.params["const"], fit.params[1:]
print(fit.tvalues["const"])        # t-stat of alpha
\end{lstlisting}
\end{addnote}

\subsection{Measures of the riskiness of returns}
\begin{enumerate}
  \item \textbf{Volatility}: standard deviation of returns.
  \item \textbf{Semi-volatility}: like volatility, but only observations below the mean:
    \[ \text{Semi-vol}=\sqrt{\frac1T\sum_{t:\,x_t<\hat\mu}(x_t-\hat\mu)^2}. \]
  \item \textbf{Skewness} and 4.\ \textbf{Kurtosis} (Chapter 1).
  \setcounter{enumi}{4}
  \item \textbf{Average maximum drawdown}: time-series average of the maximum
    percentage loss of portfolio value $V(\tau)$ over any period
    $\tau_1\to\tau_2$ within the last twelve months:
    \[
    MDD=\frac{1}{T-12}\sum_{t=12}^{T-1}\ \max_{t-11\le\tau_1<\tau_2\le t}
    \Big\{0,\,1-\frac{V(\tau_2)}{V(\tau_1)}\Big\}\times100 .
    \]
  \item \textbf{Value at Risk (VaR)}: the loss that is exceeded with a given
    probability (1\% or 5\%) over a given horizon, in normal market conditions.
    E.g.\ a one-week 99\% VaR of \$10m: a 1\% chance of losing more than \$10m
    in a week, i.e.\ once every 100 weeks, about once every two years.
  \item \textbf{Conditional VaR (CVaR)} at $q\%$: the expected return in the worst
    $q\%$ of cases. Also called average VaR, \textbf{expected shortfall},
    expected tail loss, superquantile. More sensitive than VaR to the
    \emph{shape} of the tail.
\end{enumerate}

\begin{addnote}
In terms of the return quantile $q_\alpha$ (e.g.\ $\alpha=1\%$):
$\text{VaR}_\alpha=-q_\alpha$ and $\text{CVaR}_\alpha=-E[R\mid R\le q_\alpha]$,
so $\text{CVaR}\ge\text{VaR}$ always. VaR says nothing about \emph{how much} you
lose beyond the threshold; CVaR does. Two tails with the same VaR can have very
different CVaRs.
\end{addnote}


\begin{workedex}[: VaR and CVaR under Normality]
Monthly excess return with mean $0.5\%$ and volatility $4\%$, Normal. The 1\%
quantile of a standard Normal is $z_{0.01}=-2.326$.
\begin{itemize}
  \item 1\% quantile of the return: $0.005-2.326\times0.04=-8.8\%$, so
    $\text{VaR}_{1\%}=8.8\%$ (``in 1 month out of 100 we lose more than 8.8\%'').
  \item CVaR for a Normal: $\text{CVaR}_\alpha=\sigma\,\varphi(z_\alpha)/\alpha-\mu$,
    with $\varphi$ the standard Normal density: $0.04\times0.02665/0.01-0.005=10.2\%$
    (``when we are in the worst 1\% of months, we lose 10.2\% on average'').
\end{itemize}
Compare with Python exercise 1: strategy A (Normal) has VaR 9.0\% and CVaR 10.4\%,
close to these values; strategy B (crash risk) has VaR 20.7\%: the Normal
formula would underestimate its risk by more than half.
\end{workedex}
\subsection{Measures combining mean and risk}

\begin{keybox}[: the Sharpe ratio (main measure of the course)]
\[
SR_p=\frac{E[R_p-R_f]}{\sqrt{V[R_p-R_f]}}=\frac{\text{excess mean return}}{\text{total portfolio risk}}.
\]
For IID returns over $h$ periods the numerator scales with $h$ and the denominator
with $\sqrt h$, so $SR$ scales with $\sqrt h$: \textbf{annualise a monthly Sharpe
ratio by multiplying by $\sqrt{12}\approx3.46$.}
\end{keybox}

\begin{plainbox}
The Sharpe ratio answers: ``how much extra return do I get for each unit of risk
I take?'' An annual SR of 0.5 means an expected excess return of half a
standard deviation per year. The stock market's long-run SR is around 0.4--0.5.
Because returns grow with time and volatility with its square root, the SR
grows with the square root of the horizon.
\end{plainbox}
\paragraph{Correct annualisation (Lo 2002).} $\sqrt{12}$ is correct only if returns
have no serial correlation. With first-order autocorrelation $\rho$ (AR(1)),
the multiplier for $q$ periods is
\[
\eta(q)=\frac{q}{\sqrt{q+2\sum_{k=1}^{q-1}(q-k)\rho^k}}.
\]
Positive autocorrelation $\Rightarrow$ multiplier smaller than $\sqrt{12}$
(the $\sqrt{12}$ rule overstates the SR); negative $\Rightarrow$ larger.

\begin{addnote}
\textbf{Why.} The variance of the sum of $q$ correlated returns is
$V\big[\sum_{i=1}^q x_i\big]=\sigma^2\big[q+2\sum_{k=1}^{q-1}(q-k)\rho_k\big]$:
there are $q-k$ pairs at lag $k$. The mean of the sum is $q\mu$. So the
$q$-period SR is $\frac{q\mu}{\sigma\sqrt{q+2\sum(q-k)\rho_k}}=\eta(q)\,SR_1$.
For an AR(1), $\rho_k=\rho^k$. With $\rho=0$, $\eta(q)=\sqrt q$.
\end{addnote}

\paragraph{Limitations of the Sharpe ratio.}
\begin{itemize}
  \item Appropriate for \emph{efficient portfolios}, not individual assets (it ignores correlations).
  \item Not appropriate when returns have high skewness and kurtosis (only two moments).
  \item Ignores the full distribution (the \textbf{Omega} measure uses it).
  \item Still, it is the main measure used throughout the course.
\end{itemize}

\begin{center}
\renewcommand{\arraystretch}{1.6}
\begin{tabularx}{\textwidth}{@{}l>{\centering\arraybackslash}p{5.2cm}X@{}}
\toprule
Measure & Formula & Notes \\
\midrule
\textbf{Sortino} & $\dfrac{E[R_p]-R_{\text{target}}}{\sqrt{DV}}$,\ \ $DV=\frac1T\sum\min(R_{t,p}-R_{\text{target}},0)^2$ & Penalises only \emph{downside} deviations from the target. \\
\textbf{Treynor} & $\dfrac{E[R_p-R_f]}{\beta_{p,mkt}}$,\ \ $\beta_{p,mkt}=\dfrac{C[R_p,R_{mkt}]}{V[R_{mkt}]}$ & Risk = systematic only; ignores unsystematic risk. SR and TR rank, but do not say \emph{by how much} one is better. \\
\textbf{M$^2$ / RAP} & $E[R_p-R_f]\times\dfrac{\sqrt{V[R_B]}}{\sqrt{V[R_p-R_f]}}=SR_p\sqrt{V[R_B]}$ & Modigliani--Modigliani. Excess return rescaled to the benchmark's volatility; same ranking as SR; useful to plot cumulative returns at equal volatility. \\
\textbf{Information ratio} & $\dfrac{E[R_p-R_B]}{\sqrt{V[R_p-R_B]}}$ & Benchmark is a risky index; denominator is \emph{tracking error}. SR for absolute, IR for relative performance. \\
\textbf{Jensen's alpha} & $(R_p-R_f)-(R_{\text{model}}-R_f)$ & Realised minus model-implied excess return (CAPM alpha, FFC four-factor alpha, \dots). \\
\textbf{Certainty equivalent} & depends on utility & Adjusts also for higher moments, but requires choosing the risk aversion. Conclusions usually similar to SR. \\
\bottomrule
\end{tabularx}
\end{center}

\begin{addnote}
\textbf{Certainty-equivalent return} for a mean--variance investor (as in
DeMiguel, Garlappi and Uppal 2009, which comes back in Chapter 4):
$\text{CEQ}_p=\hat\mu_p-\frac{\gamma}{2}\hat\sigma_p^2$, the risk-free return
that the investor finds as good as the risky strategy.
\end{addnote}

\subsection*{Python exercise 1: same Sharpe ratio, different risks}
50 years of monthly excess returns. A: Normal. B: small gains most months, rare
large losses (like selling insurance). Both rescaled to mean 0.5\% and vol 4\%.
\begin{lstlisting}
import numpy as np
import pandas as pd
from scipy import stats
rng = np.random.default_rng(21)
T = 600   # 50 years of monthly excess returns
a = rng.normal(0, 1, T)                               # Strategy A: Normal
crash = rng.random(T) < 0.03                          # Strategy B: 3% crash months
b = np.where(crash, -8.0, 0.5) + rng.normal(0, 0.5, T)
x = {"A": a, "B": b}
x = {k: 0.005 + 0.04*(x[k] - x[k].mean())/x[k].std() for k in x}  # same mean & vol

wealth = {k: np.cumprod(1 + x[k]) for k in x}
peak = {k: np.maximum.accumulate(wealth[k]) for k in x}
q = {k: np.quantile(x[k], 0.01) for k in x}
table = pd.DataFrame({
    "Sharpe":  [np.sqrt(12)*x[k].mean()/x[k].std() for k in x],
    "Sortino": [np.sqrt(12)*x[k].mean() / np.sqrt(np.mean(np.minimum(x[k], 0)**2)) for k in x],
    "skew":    [stats.skew(x[k]) for k in x],
    "VaR 1%":  [-q[k] for k in x],
    "CVaR 1%": [-x[k][x[k] <= q[k]].mean() for k in x],
    "max DD":  [(1 - wealth[k]/peak[k]).max() for k in x]}, index=list(x))
print(table.round(3))
\end{lstlisting}
\begin{lstlisting}[style=out]
   Sharpe  Sortino    skew  VaR 1%  CVaR 1%  max DD
A   0.433    0.661  -0.162   0.090    0.104   0.593
B   0.433    0.469  -4.584   0.207    0.218   0.452
\end{lstlisting}
\begin{itemize}
  \item Identical Sharpe ratios by construction.
  \item Sortino, skewness, VaR and CVaR all reveal B's \textbf{crash risk}: its 1\% VaR is more than twice A's.
  \item Max drawdown happens to be larger for A in this sample: \textbf{no single risk measure tells the whole story}.
\end{itemize}

\begin{addnote}
\textbf{Code to remember.} Maximum drawdown in two lines:
\texttt{wealth = np.cumprod(1+r)}, \texttt{peak = np.maximum.accumulate(wealth)},
then \texttt{(1 - wealth/peak).max()}. Note that in this exercise the Sortino
target is 0 (\texttt{np.minimum(x, 0)}).
\end{addnote}

\subsection*{Python exercise 2: annualising the Sharpe ratio}
100{,}000 years of monthly AR(1) excess returns with $\rho\in\{-0.2,0,0.3\}$.
\begin{lstlisting}
import numpy as np
import pandas as pd
from scipy.signal import lfilter
rng = np.random.default_rng(22)
T = 1_200_000
eps = rng.normal(0, 0.04, T)
rhos = [-0.2, 0.0, 0.3]
x = {r: 0.005 + lfilter([1], [1, -r], eps) for r in rhos}   # AR(1): x_t = rho x_{t-1} + eps_t
xa = {r: x[r].reshape(-1, 12).sum(axis=1) for r in rhos}     # annual returns
sr_m = {r: x[r].mean()/x[r].std() for r in rhos}

k = np.arange(1, 12)
lo = {r: 12/np.sqrt(12 + 2*np.sum((12-k)*r**k)) for r in rhos}  # Lo (2002) multiplier
table = pd.DataFrame({
    "sqrt(12) x monthly SR": [np.sqrt(12)*sr_m[r] for r in rhos],
    "Lo multiplier x SR":    [lo[r]*sr_m[r] for r in rhos],
    "actual annual SR":      [xa[r].mean()/xa[r].std() for r in rhos]}, index=rhos)
print(table.round(3))
\end{lstlisting}
\begin{lstlisting}[style=out]
      sqrt(12) x monthly SR  Lo multiplier x SR  actual annual SR
-0.2                  0.428               0.515             0.517
 0.0                  0.437               0.437             0.439
 0.3                  0.419               0.317             0.317
\end{lstlisting}
With $\rho=0.3$, $\sqrt{12}$ overstates the annual SR by about a third
(0.419 vs.\ 0.317); with $\rho<0$ it understates it. The Lo multiplier is close
to the truth in both cases.

\begin{addnote}
\texttt{lfilter([1], [1, -r], eps)} computes the recursion
$y_t=\varepsilon_t+\rho\,y_{t-1}$ quickly (a loop would be slow for 1.2 million
observations). Positive autocorrelation is typical of illiquid assets and hedge
funds with smoothed returns, so their annualised Sharpe ratios are often
overstated.
\end{addnote}

% ---------------------------------------------------------------------
\section{Transaction costs and turnover}

\begin{definitionbox}{Turnover}
When rebalancing at $t+1$, the trade in asset $n$ is $|w_{t+1,n}-w_{t^+,n}|$, where
$w_{t+1,n}$ is the \textbf{new} weight and $w_{t^+,n}$ is the \textbf{old} date-$t$
weight evaluated at date-$(t+1)$ prices (it has drifted with returns).
\[
\text{Turnover}_{t+1}=\sum_{n=1}^N\big|w_{t+1,n}-w_{t^+,n}\big|,\qquad
\overline{\text{Turnover}}=\frac1T\sum_{t=0}^{T-1}\sum_{n=1}^N\big|w_{t+1,n}-w_{t^+,n}\big|\times12 .
\]
\end{definitionbox}

\begin{keybox}[: weights just before rebalancing]
After returns realise, the old weights drift to
\[
w_{t^+,n}=\frac{w_{t,n}\,(1+r_{t+1,n})}{1+r_{t+1,p}},\qquad r_{t+1,p}=\sum_n w_{t,n}r_{t+1,n}.
\]
\end{keybox}
\begin{addnote}
\textbf{Why.} Dollar holding in $n$ grows by $(1+r_{t+1,n})$, total wealth by
$(1+r_{t+1,p})$; the new share of wealth is the ratio. Consequence: even
\textbf{1/N has positive turnover}, because you must sell winners and buy losers
to return to equal weights. A buy-and-hold portfolio has zero turnover.
\end{addnote}

\begin{plainbox}
\textbf{Numerical example.} 1/N with two assets: 50/50. Next month asset A gains
10\%, asset B is flat. Now A is $0.55/1.05=52.4\%$ and B is $47.6\%$. To go back to
50/50 you sell 2.4\% of A and buy 2.4\% of B: turnover $=|0.5-0.524|+|0.5-0.476|=4.8\%$.
With a cost of 50 bp per unit traded you pay $0.005\times0.048=0.024\%$ of wealth.
\end{plainbox}
\paragraph{Industry turnover ratio.}
$\dfrac{\min(\text{securities sold},\ \text{securities purchased})}{\text{average net assets}}\times100$.
Example: bought \$10m, sold \$9m, average assets \$100m $\Rightarrow$ 9\%.
It counts only one side, while the academic turnover counts both.

\begin{keybox}[: returns net of proportional transaction costs]
With proportional cost $\kappa$ (no fixed costs):
\begin{gather*}
W_{t+1,p}=W_{t,p}\,R_{t+1,p}\Big(1-\kappa\sum_{n=1}^N|w_{t+1,n}-w_{t^+,n}|\Big),\\
r^{\text{net}}_{t+1,p}=R_{t+1,p}\Big(1-\kappa\sum_{n=1}^N|w_{t+1,n}-w_{t^+,n}|\Big)-1 .
\end{gather*}
Typical $\kappa$: 50 bp (Balduzzi--Lynch 1999, based on NYSE studies);
10--20 bp for large stocks today; 0--5 bp for a hedge fund. See French (2008).
\end{keybox}

\subsection*{Python exercise 3: turnover and break-even transaction cost}
10 assets, 20 years; true means change every month; we observe a noisy forecast.
A = 1/N; B = 1/N + tilt towards assets with the highest forecasts.
\begin{lstlisting}
import numpy as np
import pandas as pd
rng = np.random.default_rng(23)
N, T = 10, 240
mu = rng.normal(0.008, 0.01, size=(T, N))          # true means
f = rng.normal(0, 0.04, size=(T, 1))               # market factor
R = mu + f + rng.normal(0, 0.05, size=(T, N))      # returns
signal = mu + rng.normal(0, 0.01, size=(T, N))     # noisy forecast
w = {"A": np.full((T, N), 1/N),
     "B": 1/N + 2*(signal - signal.mean(axis=1, keepdims=True))}
rp = {s: (w[s]*R).sum(axis=1) for s in w}
w_plus = {s: w[s]*(1 + R)/(1 + rp[s])[:, None] for s in w}   # drifted weights
turn = {s: np.abs(w[s][1:] - w_plus[s][:-1]).sum(axis=1) for s in w}

kappas = [0, 10, 25, 50, 100]                                # in basis points
net = {(s, k): (1 - k/1e4*turn[s])*(1 + rp[s][1:]) - 1 for s in w for k in kappas}
sr = {key: np.sqrt(12)*net[key].mean()/net[key].std() for key in net}
print("Average annual turnover:", {s: round(float(12*turn[s].mean()), 2) for s in w})
table = pd.DataFrame({s: [sr[(s, k)] for k in kappas] for s in w}, index=kappas)
table.index.name = "cost (bp)"
print(table.round(3))
\end{lstlisting}
\begin{lstlisting}[style=out]
Average annual turnover: {'A': 0.47, 'B': 3.69}
cost (bp)   SR_A   SR_B
0          0.784  0.958
10         0.781  0.935
25         0.777  0.900
50         0.770  0.842
100        0.755  0.725
\end{lstlisting}
B trades about 8 times more than 1/N. Before costs B wins (0.96 vs.\ 0.78);
the advantage disappears between 50 and 100 bp. \textbf{The value of a strategy
depends on its turnover.}

\paragraph{Limitations of this three-step approach} (optimise $\to$ measure
turnover $\to$ multiply by $\bar\kappa$):
\begin{enumerate}
  \item Costs should be included \emph{when optimising}, not after;
  \item the relevant turnover is that of the cost-aware portfolio (smaller);
  \item use the cost of each stock at each date, $\kappa_{t,n}$, not an average $\bar\kappa$.
\end{enumerate}
Later in the course: optimise with costs, use $\kappa_{t,n}$, and \textbf{net
trades across factors} before computing costs (DeMiguel, Mart\'in-Utrera,
Nogales, Uppal 2020; DeMiguel, Mart\'in-Utrera, Uppal 2024). Turnover also
creates capital-gains taxes (not studied; see Dybvig--Koo 1996, DeMiguel--Uppal 2005).

\paragraph{Turnover by style.} Higher for active (vs.\ passive), growth (vs.\
value), quantitative (vs.\ qualitative). \textbf{Controlling turnover}: account
for costs when rebalancing; long horizon; buy-and-hold; tax-efficient vehicles.

% ---------------------------------------------------------------------
\section{Review of hypothesis testing}
\paragraph{Steps.} (1) State the hypothesis; (2) choose a test statistic and its
distribution; (3) set the significance level; (4) state the decision rule;
(5) collect data; (6) compute the statistic; (7) statistical decision;
(8) investment decision.

\paragraph{Test of a mean.} $H_0:\mu=\mu_0$ vs.\ $H_1:\mu\ne\mu_0$. If
$\bar X\sim\mathcal N(\mu_0,\sigma^2/T)$ under $H_0$:
\[
z=\frac{\bar X-\mu_0}{\sigma/\sqrt T}\sim\mathcal N(0,1),\qquad
P\big(-1.96\le z\le1.96\big)=0.95 .
\]
Reject $H_0$ at 5\% (two-sided) if $|z|>1.96$; one-sided ($H_1:\mu<\mu_0$) if
$z<-1.645$.

\begin{definitionbox}{p-value, Type I/II errors, power}
\begin{itemize}
  \item \textbf{p-value}: probability of a result at least as extreme as the
    observed one, assuming $H_0$ is true; the lowest level at which $H_0$ can be
    rejected. Two-sided: $p=2\,\Phi(-|z|)$.
  \item \textbf{Type I error}: reject $H_0$ when it is true (probability $\alpha$, the level of the test).
  \item \textbf{Type II error}: do not reject $H_0$ when it is false (probability $\beta$).
  \item \textbf{Power} $=1-\beta$: probability of rejecting $H_0$ when an
    alternative is true. Fix $\alpha$ (usually 5\%) and maximise power. Low power
    $\Rightarrow$ the test cannot discriminate between $H_0$ and $H_1$.
\end{itemize}
\end{definitionbox}

\begin{plainbox}
A test asks: ``if the boring hypothesis were true (e.g.\ no difference), how
surprising would my data be?'' The $z$-statistic measures the surprise in
standard errors. Beyond about 2 standard errors (1.96) we call it too surprising
and reject. The $p$-value is the probability of being at least that surprised by
pure chance.
\end{plainbox}
\paragraph{Worked example.} $H_0:\mu=15$, $\bar X=13$, $\sigma=3.0456$, $T=30$:
\[
z=\frac{13-15}{3.0456/\sqrt{30}}=-3.597,\qquad
p=2\times P[z<-3.597]=2\times0.000161=0.000322 .
\]
$|z|>1.96$ $\Rightarrow$ reject $H_0$. Reporting the $p$-value avoids arbitrary
conventional levels.

% ---------------------------------------------------------------------
\section{Test of the difference in Sharpe ratios (focus)}
We test $H_0: SR_n-SR_m=0$ (portfolio $n$ vs.\ benchmark $m$). Same logic as the
test of a mean, with $\bar X\to \widehat{SR}_n-\widehat{SR}_m$ and $\mu_0\to0$:
\[
z=\frac{\widehat{SR}_n-\widehat{SR}_m}{\sqrt{V[\widehat{SR}_n-\widehat{SR}_m]}} .
\]
\textbf{All we need is the variance of the difference.} Two methods: the delta
method and the bootstrap.

\subsection{The delta method}
\begin{keybox}[: delta method]
If $\sqrt T(\hat\theta-\theta_0)\xrightarrow{d}\mathcal N(0,\Omega)$ and $f$ is
differentiable with gradient $\nabla f(\theta_0)$, then
\[
\sqrt T\{f(\hat\theta)-f(\theta_0)\}\xrightarrow{d}\mathcal N(0,V),\qquad
V=\nabla f(\theta_0)\T\,\Omega\,\nabla f(\theta_0),
\qquad \widehat{se}\{f(\hat\theta)\}=\sqrt{\hat V/T}.
\]
Test $H_0:f(\theta_0)=c$ with $z=\frac{f(\hat\theta)-c}{\widehat{se}}\approx\mathcal N(0,1)$, $p=2\Phi(-|z|)$.
\end{keybox}
\begin{intuition}
In a small enough range, any differentiable $f$ is approximately \emph{linear}
(first-order Taylor expansion), and a linear function of a Normal variable is
Normal. As $T$ grows the estimator's variance shrinks, so the linear
approximation is applied over a smaller range and becomes more accurate.
\end{intuition}

\begin{plainbox}
We know how noisy a sample mean and a sample variance are. The Sharpe ratio is
a \emph{function} of the two (mean divided by square root of variance). The delta
method converts the noise of the inputs into the noise of the output by
linearising the function (a first-order Taylor expansion). Steps 1--3 do exactly
that, first for one SR, then for the difference of two.
\end{plainbox}
\subsection{Step 1: distribution of $\hat\mu_n$ and $\hat\sigma^2_n$}
For IID Normal excess returns, with $\bm u=(\mu_n,\sigma_n^2)\T$:
\[
\sqrt T(\hat{\bm u}-\bm u)\to\mathcal N\!\left(\bm 0_2,\ \Omega\right),\qquad
\Omega=\begin{pmatrix}\sigma_n^2&0\\0&2\sigma_n^4\end{pmatrix}
\quad\Rightarrow\quad
V[\hat\mu_n]\overset{a}{=}\frac{\sigma_n^2}{T},\quad V[\hat\sigma_n^2]\overset{a}{=}\frac{2\sigma_n^4}{T}.
\]

\subsection{Step 2: distribution of $\widehat{SR}_n$}
$f(\bm u)=\mu_n/\sqrt{\sigma_n^2}$. Gradient and quadratic form:
\[
f_u=\begin{pmatrix}\dfrac1{\sigma_n}\\[6pt]-\dfrac{\mu_n}{2\sigma_n^3}\end{pmatrix},\qquad
f_u\T\Omega f_u=\frac{\sigma_n^2}{\sigma_n^2}+\frac{\mu_n^2}{4\sigma_n^6}\,2\sigma_n^4
=1+\frac12 SR_n^2 .
\]
\begin{keybox}[: standard error and CI of the Sharpe ratio (IID Normal)]
\[
\sqrt T\big(\widehat{SR}_n-SR_n\big)\to\mathcal N\!\Big(0,\,1+\tfrac12SR_n^2\Big),\qquad
SE[\widehat{SR}_n]=\sqrt{\frac1T\Big(1+\frac12\widehat{SR}_n^2\Big)},
\]
\[
z=\frac{\widehat{SR}_n-SR_0}{SE[\widehat{SR}_n]}\approx\mathcal N(0,1),\qquad
CI_{95\%}=\widehat{SR}_n\pm1.96\sqrt{\frac1T\Big(1+\frac12\widehat{SR}_n^2\Big)}.
\]
Robust variants replace $\sigma^2/T$ by a HAC long-run variance.
\end{keybox}
\begin{addnote}
\textbf{Numerical example.} Monthly $SR=0.15$ (annual $\approx0.52$), $T=120$ months:
$SE=\sqrt{(1+0.5\cdot0.0225)/120}=0.0918$ monthly, i.e.\ $\approx0.32$ annualised.
The 95\% CI for the annual SR is about $0.52\pm0.62$: it includes zero. Ten years
of data are not enough to prove that a Sharpe ratio of 0.5 is positive.
(Note: the $SR$ inside the formula must be at the \emph{same frequency} as $T$.)
\end{addnote}

\subsection{Step 3: distribution of $\widehat{SR}_n-\widehat{SR}_m$}
IID Normal excess returns $R_{t,n},R_{t,m}$ with covariance $\sigma_{nm}$ and
correlation $\rho_{nm}=\sigma_{nm}/(\sigma_n\sigma_m)$.
$\bm u=(\mu_n,\mu_m,\sigma_n^2,\sigma_m^2)\T$:
\[
\Omega=\begin{pmatrix}
\sigma_n^2&\sigma_{nm}&0&0\\
\sigma_{nm}&\sigma_m^2&0&0\\
0&0&2\sigma_n^4&2\sigma_{nm}^2\\
0&0&2\sigma_{nm}^2&2\sigma_m^4
\end{pmatrix},\qquad
f(\bm u)=\frac{\mu_n}{\sigma_n}-\frac{\mu_m}{\sigma_m},\qquad
f_u=\begin{pmatrix}1/\sigma_n\\-1/\sigma_m\\-\mu_n/(2\sigma_n^3)\\ \mu_m/(2\sigma_m^3)\end{pmatrix}.
\]
\begin{addnote}
\textbf{Doing the multiplication.} $\Omega$ is block-diagonal, so the quadratic
form splits into a ``means'' block and a ``variances'' block:
\begin{align*}
\text{means: }&\frac{\sigma_n^2}{\sigma_n^2}+\frac{\sigma_m^2}{\sigma_m^2}-2\frac{\sigma_{nm}}{\sigma_n\sigma_m}=2-2\rho_{nm},\\
\text{variances: }&\frac{\mu_n^2}{4\sigma_n^6}2\sigma_n^4+\frac{\mu_m^2}{4\sigma_m^6}2\sigma_m^4
-2\frac{\mu_n\mu_m}{4\sigma_n^3\sigma_m^3}2\sigma_{nm}^2
=\tfrac12\big(SR_n^2+SR_m^2-2SR_nSR_m\rho_{nm}^2\big).
\end{align*}
\end{addnote}

\begin{keybox}[: Jobson--Korkie (1981) test with Memmel (2003) correction]
\begin{gather*}
V\big[\widehat{SR}_n-\widehat{SR}_m\big]=\frac1T\Big[2-2\hat\rho_{nm}+\frac12\big(\widehat{SR}_n^2+\widehat{SR}_m^2-2\widehat{SR}_n\widehat{SR}_m\hat\rho_{nm}^2\big)\Big],\\
z=\frac{\widehat{SR}_n-\widehat{SR}_m}{\sqrt{V[\widehat{SR}_n-\widehat{SR}_m]}} .
\end{gather*}
If $\rho_{nm}=0$: $V=\frac1T\big[(1+\frac12SR_n^2)+(1+\frac12SR_m^2)\big]$, the sum of the two variances.
\end{keybox}
\begin{intuition}
The term $2-2\rho_{nm}$ dominates. \textbf{The higher the correlation between the
two strategies, the smaller the variance of the difference}, so the easier it is
to detect a given difference in Sharpe ratios. Comparing two highly correlated
strategies (e.g.\ MV vs.\ a slight variant) is much easier than comparing two
unrelated ones.
\end{intuition}


\begin{workedex}[: applying the JKM test]
Ten years of monthly data ($T=120$). Monthly Sharpe ratios $\widehat{SR}_n=0.15$
(annual 0.52) and $\widehat{SR}_m=0.10$ (annual 0.35); correlation $\hat\rho=0.8$.
\[
V=\frac{1}{120}\Big[2-1.6+\tfrac12\big(0.0225+0.01-2\times0.15\times0.10\times0.64\big)\Big]
=\frac{0.4+0.0067}{120}=0.00339,
\]
so $SE=\sqrt{0.00339}=0.0582$.
$z=(0.15-0.10)/0.0582=0.86$, $p=2\Phi(-0.86)=0.39$: \textbf{not significant}.
An annual SR difference of 0.17 over ten years cannot be distinguished from luck.

With $\hat\rho=0$ instead: $SE=0.130$ and $z=0.39$, even weaker. High correlation
between the two strategies is what makes the test more powerful (the $2-2\rho$ term).
\end{workedex}
\paragraph{Limitations of the delta method (JKM).} Not valid with
\textbf{non-normality} (heavy tails, skewness) or \textbf{time dependence}
(autocorrelation, volatility clustering). See the QuantPy article recommended in
the slides. Alternative: \textbf{bootstrap}.

\subsection*{Python exercise 4: when is the JKM test reliable?}
5{,}000 samples of 10 years; two strategies with the \emph{same} true SR
($H_0$ true), correlation 0.5. Case IID vs.\ AR(1) with $\rho=0.3$.
\begin{lstlisting}
import numpy as np
from scipy.signal import lfilter
rng = np.random.default_rng(24)
M, T = 5000, 120
rho_nm = 0.5
z1 = rng.normal(size=(M, T))
z2 = rho_nm*z1 + np.sqrt(1-rho_nm**2)*rng.normal(size=(M, T))
ar = {"IID": 0.0, "AR(1)": 0.3}
n = {c: 0.005 + 0.04*lfilter([1], [1, -ar[c]], z1, axis=1) for c in ar}
m = {c: 0.005 + 0.04*lfilter([1], [1, -ar[c]], z2, axis=1) for c in ar}

srn = {c: n[c].mean(1)/n[c].std(1, ddof=1) for c in ar}
srm = {c: m[c].mean(1)/m[c].std(1, ddof=1) for c in ar}
rho = {c: np.array([np.corrcoef(n[c][i], m[c][i])[0, 1] for i in range(M)]) for c in ar}
var = {c: (2 - 2*rho[c] + 0.5*(srn[c]**2 + srm[c]**2
       - 2*srn[c]*srm[c]*rho[c]**2))/T for c in ar}
zst = {c: (srn[c] - srm[c])/np.sqrt(var[c]) for c in ar}
print("Rejection rate at 5% (should be 0.05):",
      {c: round(float(np.mean(np.abs(zst[c]) > 1.96)), 3) for c in ar})
\end{lstlisting}
\begin{lstlisting}[style=out]
Rejection rate at 5% (should be 0.05): {'IID': 0.053, 'AR(1)': 0.15}
\end{lstlisting}
IID: rejects $\approx5\%$, as it should. AR(1): rejects 15\%, \textbf{three
times too often} (size distortion). Hence the need for tests robust to serial
dependence: the block bootstrap of Ledoit and Wolf (2008).

\subsection{Bootstrapping}
\begin{definitionbox}{Bootstrap}
Estimate the sampling distribution of an estimator by \textbf{resampling the
data with replacement}. Useful when parametric assumptions are in doubt, when
parametric inference is impossible, or when standard-error formulas are complex.
\end{definitionbox}
\paragraph{Intuition.} From one sample of French heights you get one mean. Draw a
new sample of the same size $N$ from your data with replacement (e.g.\ from
$[1,2,3,4,5]$ you might get $[2,5,4,4,1]$), compute its mean, repeat 1{,}000--10{,}000
times. The histogram of these bootstrap means estimates the sampling
distribution of the mean. Works for almost any statistic. Reference: Efron and
Tibshirani (1994).


\begin{workedex}[: one bootstrap round by hand]
Data (5 monthly returns, in \%): $[2,\ -1,\ 3,\ 0,\ 1]$, mean $1.0$.
Draw 5 positions with replacement, e.g.\ $(3,1,5,5,2)$: the bootstrap sample is
$[3,\ 2,\ 1,\ 1,\ -1]$, mean $1.2$. Repeat 5{,}000 times: the standard deviation
of the 5{,}000 means is the bootstrap standard error; the 2.5\% and 97.5\%
quantiles give a 95\% confidence interval. For a Sharpe ratio difference you do
the same with \emph{pairs} of returns and compute $\widehat{SR}_n-\widehat{SR}_m$
in each round.
\end{workedex}
\subsection{Block bootstraps for time series}
The IID bootstrap destroys the time ordering, so it cannot reproduce serial
correlation. Solution: resample \textbf{blocks} of consecutive observations.
\begin{center}
\begin{tabularx}{\textwidth}{@{}lX@{}}
\toprule
Method & How blocks are formed \\
\midrule
Block bootstrap & Non-overlapping blocks, e.g.\ $(X_1,X_2),(X_3,X_4),\dots,(X_9,X_{10})$; draw blocks with replacement (e.g.\ 3,1,5,2,5). \\
Moving block (MBB) & Overlapping blocks $(X_1,X_2),(X_2,X_3),\dots,(X_9,X_{10})$. \\
Stationary & \textbf{Random} block lengths (geometric, mean $\ell$) and random starting points (Politis--Romano). \\
Circular block & Data treated as a circle: blocks wrap around from the end to the start, so edge points are sampled as often as middle points (Politis--Romano 1992). \textbf{Recommended}; used by Ledoit--Wolf (2008). \\
\bottomrule
\end{tabularx}
\end{center}
\begin{plainbox}
Financial returns have memory (calm and turbulent periods cluster). If you
shuffle single months you destroy that memory and underestimate uncertainty. So
you shuffle \emph{blocks} of consecutive months: inside each block the memory
survives. And you always move the two strategies' returns together, month by
month, to keep their correlation.
\end{plainbox}
Resample \textbf{pairs} $(R_{t,n},R_{t,m})$ together, to preserve the correlation
between the two strategies.

\subsection{Ledoit--Wolf (2008) bootstrap $p$-value for $\Delta SR$}
$H_0:SR_n=SR_m$ vs.\ $H_1:SR_n\ne SR_m$. Studentised statistic
$S=\widehat{\Delta SR}/\widehat{SE}(\widehat{\Delta SR})$, with a HAC or bootstrap SE.

\begin{keybox}[: practical recipe]
\begin{enumerate}
  \item Compute $\widehat{SR}_n,\widehat{SR}_m$, $\widehat{\Delta SR}$ and
    $S_{obs}=\widehat{\Delta SR}/\widehat{SE}$, with $\widehat{SE}$ from the
    delta method using a HAC (Newey--West) long-run covariance.
  \item Choose block length $\ell$ (rule of thumb $\ell\propto T^{1/3}$; LW
    calibrate it); circular or stationary bootstrap.
  \item For $b=1,\dots,B$: draw a bootstrap series of \emph{pairs} by
    concatenating random blocks, truncated at $T$; compute
    $\widehat{\Delta SR}^{*(b)}$ and the \textbf{centred} studentised statistic
    \[ S_b^*=\frac{\widehat{\Delta SR}^{*(b)}-\widehat{\Delta SR}}{\widehat{SE}^{*(b)}} . \]
    Centring is needed because the bootstrap samples come from data in which
    $H_0$ does not hold exactly.
  \item Two-sided $p$-value:
    $\hat p=\Big(1+\sum_{b=1}^B\mathbb 1\{|S_b^*|\ge|S_{obs}|\}\Big)\big/(B+1)$.
\end{enumerate}
\end{keybox}

\subsection*{Sample code (optional in the slides; key functions)}
\begin{lstlisting}
import numpy as np
from numpy.random import default_rng

def sharpe(x):
    x = np.asarray(x); s = x.std(ddof=1)
    return np.nan if s == 0 else x.mean()/s

# IID bootstrap of the mean and percentile CI
rng = default_rng(42)
data = rng.normal(0.0, 1.0, size=200)
B = 5000
boot_means = np.array([rng.choice(data, size=data.size, replace=True).mean()
                       for _ in range(B)])
se_mean = boot_means.std(ddof=1)
ci_mean = np.quantile(boot_means, [0.025, 0.975])

def block_bootstrap_pairs(a, b, ell, rng):
    # Circular block bootstrap: blocks of length ell may wrap around the end
    n = len(a)
    num_blocks = int(np.ceil(n/ell))
    starts = rng.integers(0, n, size=num_blocks)
    idx = np.concatenate([np.arange(s, s+ell) % n for s in starts])[:n]
    return a[idx], b[idx]

def _newey_west_variance_of_mean(series, max_lag=None):
    x = np.asarray(series); n = x.size; x = x - x.mean()
    if max_lag is None:
        max_lag = int(np.floor(4*(n/100)**(2/9)))
    gamma0 = np.dot(x, x)/n
    cov_k = lambda k: np.dot(x[:-k], x[k:])/n
    w = lambda k: 1 - k/(max_lag + 1.0)              # Bartlett weights
    s = gamma0 + 2.0*sum(w(k)*cov_k(k) for k in range(1, max_lag+1))
    return s/n

def delta_sr_se_hac(a, b, max_lag=None):
    a = np.asarray(a); b = np.asarray(b)
    muA, muB = a.mean(), b.mean(); sA, sB = a.std(ddof=1), b.std(ddof=1)
    # influence function of SR_a - SR_b
    d = (a - muA)/sA - (b - muB)/sB \
        - (muA/(2*sA**3))*((a - muA)**2 - sA**2) \
        + (muB/(2*sB**3))*((b - muB)**2 - sB**2)
    return np.sqrt(_newey_west_variance_of_mean(d, max_lag))

def delta_sr_pvalue_block(a, b, ell, B=5000, rng=None):
    rng = default_rng() if rng is None else rng
    dObs = sharpe(a) - sharpe(b)
    TObs = dObs/delta_sr_se_hac(a, b)
    TBoot = []
    for _ in range(B):
        rA, rB = block_bootstrap_pairs(a, b, ell, rng)
        seB = delta_sr_se_hac(rA, rB)
        if np.isfinite(seB) and seB != 0:
            TBoot.append((sharpe(rA) - sharpe(rB) - dObs)/seB)   # centred
    return np.mean(np.abs(np.asarray(TBoot)) >= np.abs(TObs))

# ell = max(5, round(n**(1/3)))
\end{lstlisting}
\begin{addnote}
\begin{itemize}
  \item \textbf{Influence function.} The vector \texttt{d} is, observation by
    observation, the first-order (delta-method) contribution of each date to
    $\widehat{SR}_a-\widehat{SR}_b$. Its sample mean has variance
    $\approx V[\widehat{\Delta SR}]$; using Newey--West on \texttt{d} makes the SE
    robust to autocorrelation and heteroskedasticity. With IID Normal data it
    reproduces the JKM formula.
  \item \textbf{Newey--West.} Long-run variance $=\gamma_0+2\sum_{k=1}^{L}
    \big(1-\frac{k}{L+1}\big)\gamma_k$: autocovariances up to lag $L$ with
    declining (Bartlett) weights, which guarantee a positive variance.
  \item The slides also show bootstrapping a quantile, the \textbf{BCa}
    (bias-corrected and accelerated) interval, and the stationary bootstrap
    (restart a new block with probability $p=1/\ell$ at each step).
  \item Challenge in the slides: reproduce Tables 1 and 2 of Ledoit--Wolf (2008).
\end{itemize}
\end{addnote}

\subsection{Multiple testing and $p$-hacking}
\subsection*{Python exercise 5: the best of 200 useless strategies}
200 strategies, 10 years of monthly returns, all with \textbf{true mean zero}.
Pick the one with the highest $t$-statistic.
\begin{lstlisting}
import numpy as np
from scipy import stats
rng = np.random.default_rng(25)
K, T = 200, 120
x = rng.normal(0, 0.04, size=(K, T))                 # no skill at all
t = x.mean(axis=1) / (x.std(axis=1, ddof=1)/np.sqrt(T))
best = np.argmax(t)
p_naive = 2*stats.norm.sf(t[best])
p_bonf = min(1, K*p_naive)                           # Bonferroni correction
print("Best strategy: annual SR =", round(np.sqrt(12)*t[best]/np.sqrt(T), 2),
      " t-stat =", round(t[best], 2))
print("Naive p-value:", round(p_naive, 4))
print("Bonferroni p-value:", round(p_bonf, 3))
print("Strategies with |t| > 1.96:", np.sum(np.abs(t) > 1.96))
\end{lstlisting}
\begin{lstlisting}[style=out]
Best strategy: annual SR = 0.73   t-stat = 2.31
Naive p-value: 0.0208
Bonferroni p-value: 1
Strategies with |t| > 1.96: 11
\end{lstlisting}
\begin{keybox}[: multiple testing]
The best useless strategy has an annual SR of 0.73 and naive $p=0.02$: it looks
``significant''. 11 of 200 pass at 5\%, about the 10 expected by chance. After
\textbf{Bonferroni} ($p\times K$), $p=1$: no evidence of skill. See Harvey, Liu,
and Zhu (2015).
\end{keybox}
\begin{addnote}
Note the identity used in the code: annual SR $=\sqrt{12}\,t/\sqrt T$, since
$t=\widehat{SR}_{\text{monthly}}\sqrt T$. Harvey, Liu and Zhu (2015) argue that,
given the hundreds of factors already tried, a new factor should clear
$t>3$, not $t>1.96$. Bonferroni is conservative (it controls the probability of
\emph{even one} false discovery); less strict alternatives such as Holm or
Benjamini--Hochberg exist.
\end{addnote}

% ---------------------------------------------------------------------
\section{Takeaways of Chapter 2}
\begin{enumerate}
  \item The course studies \textbf{myopic} portfolios evaluated \textbf{out of
    sample}: weights at $t$ use only information up to $t$; performance is
    $R_{t+1,p}=\w_t\T\R_{t+1}$.
  \item Rolling windows when the past is less relevant; expanding when it is
    similar; EWMA in between.
  \item Many metrics; the \textbf{Sharpe ratio} is the main one. Annualise with
    $\sqrt{12}$ only if returns are not autocorrelated (else Lo 2002).
  \item Same Sharpe ratio does not mean same risk: look at Sortino, skewness,
    VaR, CVaR, drawdowns.
  \item Turnover uses \emph{drifted} weights $w_{t^+}$; net return
    $=R_{t+1,p}(1-\kappa\,\text{Turnover})-1$. High-turnover strategies can lose
    their edge after costs.
  \item $SE(\widehat{SR})=\sqrt{(1+SR^2/2)/T}$: Sharpe ratios are estimated imprecisely.
  \item JKM test for $SR_n-SR_m$: the variance shrinks with the correlation
    between the strategies; it over-rejects with autocorrelation $\Rightarrow$ use
    the Ledoit--Wolf circular block bootstrap.
  \item Trying many strategies inflates false discoveries: correct for multiple testing.
\end{enumerate}


\begin{mistakes}
\begin{enumerate}
  \item Computing the Sharpe ratio on total returns instead of \emph{excess} returns.
  \item Using weights estimated with data that include the evaluation month (look-ahead).
  \item Using $\sqrt{12}$ to annualise when returns are clearly autocorrelated.
  \item Computing turnover with $|w_{t+1}-w_t|$ instead of $|w_{t+1}-w_{t^+}|$: this
    ignores the drift of weights and gives 1/N zero turnover (it is not zero).
  \item Concluding that strategy A ``beats'' B because its SR is higher, without a
    test: with 10 years the SE of an annual SR is about 0.3.
  \item In the bootstrap, resampling the two strategies independently (this
    destroys their correlation) or forgetting to centre the statistic.
  \item Reporting the best of many backtests without a multiple-testing correction.
\end{enumerate}
\end{mistakes}
\section{Self-test questions (quiz/exam practice)}
\begin{enumerate}
  \item Why does the course not study dynamic portfolios? \emph{(Too many
    parameters to estimate $\Rightarrow$ poor out-of-sample performance.)}
  \item A monthly Sharpe ratio is 0.2 with IID returns. Annual SR?
    \emph{($0.2\sqrt{12}=0.69$.)} If $\rho=0.3$, is the true annual SR higher or
    lower? \emph{(Lower.)}
  \item Which measure uses a risky benchmark and tracking error? \emph{(Information ratio.)}
  \item Which measure uses only systematic risk in the denominator? \emph{(Treynor.)}
  \item Why does 1/N have positive turnover? \emph{(Weights drift with returns;
    rebalancing back to $1/N$ requires trading.)}
  \item Strategy gross return 1\%, turnover 0.5, $\kappa=50$ bp. Net return?
    \emph{($1.01\times(1-0.005\times0.5)-1=0.7475\%$.)}
  \item $T=240$ months, monthly $SR=0.1$. SE of $\widehat{SR}$? \emph{($\sqrt{1.005/240}=0.0647$.)}
  \item Two strategies with $\rho_{nm}=0.9$ vs.\ $\rho_{nm}=0$: in which case is a
    given $\Delta SR$ easier to detect? \emph{($\rho=0.9$: variance of the difference much smaller.)}
  \item Why centre the bootstrap statistic at $\widehat{\Delta SR}$? \emph{(The
    bootstrap data do not satisfy $H_0$ exactly.)}
  \item You test 50 strategies and the best has $p=0.01$. Bonferroni $p$? \emph{(0.5.)}
  \item Write by hand: Sharpe ratio, Sortino ratio, max drawdown, 5\% VaR and CVaR of a return array \texttt{r}.
\begin{lstlisting}
sr = np.sqrt(12) * r.mean() / r.std(ddof=1)
sortino = np.sqrt(12) * r.mean() / np.sqrt(np.mean(np.minimum(r, 0)**2))
wealth = np.cumprod(1 + r); mdd = (1 - wealth / np.maximum.accumulate(wealth)).max()
q = np.quantile(r, 0.05); var5 = -q; cvar5 = -r[r <= q].mean()
\end{lstlisting}
  \item Write by hand the JKM $z$-statistic for two return arrays \texttt{a}, \texttt{b}.
\begin{lstlisting}
T = len(a)
sa, sb = a.mean()/a.std(ddof=1), b.mean()/b.std(ddof=1)
rho = np.corrcoef(a, b)[0, 1]
var = (2 - 2*rho + 0.5*(sa**2 + sb**2 - 2*sa*sb*rho**2)) / T
z = (sa - sb) / np.sqrt(var)
p = 2 * stats.norm.sf(abs(z))
\end{lstlisting}
  \item Write by hand the turnover of a strategy with weight matrix \texttt{W} ($T\times N$) and returns \texttt{R}.
\begin{lstlisting}
rp = (W * R).sum(axis=1)
W_plus = W * (1 + R) / (1 + rp)[:, None]
turnover = np.abs(W[1:] - W_plus[:-1]).sum(axis=1)
\end{lstlisting}
\end{enumerate}

\section{To do for next class}
\begin{itemize}
  \item Read ``Working with Data'' in Moreira (2021), \emph{Quantitative Investing}, Chapter 5.
  \item Work on Assignment \#1 (Ch.\,1--2; due 11 Nov, 22:00 CET).
\end{itemize}

\section*{References (Chapter 2)}
\begin{itemize}[leftmargin=*]
\small
\item Balduzzi, P., and A.\,W.~Lynch. 1999. Transaction costs and predictability. \emph{JFE} 52(1): 47--78.
\item Campbell, J.\,Y., and L.\,M.~Viceira. 2002. \emph{Strategic Asset Allocation}. Oxford UP.
\item Carhart, M.\,M. 1997. On persistence in mutual fund performance. \emph{JF} 52(1): 57--82.
\item Dello-Preite, M., R.~Uppal, P.~Zaffaroni, and I.~Zviadadze. 2025. Cross-sectional asset pricing with unsystematic risk. SSRN 4135146.
\item DeMiguel, V., A.~Mart\'in-Utrera, F.\,J.~Nogales, and R.~Uppal. 2020. A transaction-cost perspective on the multitude of firm characteristics. \emph{RFS} 33(5): 2180--2222.
\item DeMiguel, V., A.~Mart\'in-Utrera, and R.~Uppal. 2024. A multifactor perspective on volatility-managed portfolios. \emph{JF} 79(6): 3859--3891.
\item Efron, B., and R.\,J.~Tibshirani. 1994. \emph{An Introduction to the Bootstrap}. CRC.
\item Fama, E.\,F., and K.\,R.~French. 1993; 2015; 2018. \emph{JFE} 33: 3--56; 116: 1--22; 128: 234--252.
\item French, K.\,R. 2008. The cost of active investing. \emph{JF} 63(4): 1537--1573.
\item Harvey, C.\,R., Y.~Liu, and H.~Zhu. 2015. \dots and the cross-section of expected returns. \emph{RFS} 29(1): 5--68.
\item Jobson, J.\,D., and B.\,M.~Korkie. 1981. Performance hypothesis testing with the Sharpe and Treynor measures. \emph{JF}: 889--908.
\item Ledoit, O., and M.~Wolf. 2008. Robust performance hypothesis testing with the Sharpe ratio. \emph{JEF} 15: 850--859.
\item Lo, A.\,W. 2002. The statistics of Sharpe ratios. \emph{FAJ} 58(4): 36--52.
\item Memmel, C. 2003. Performance hypothesis testing with the Sharpe ratio. SSRN 412588.
\item Politis, D.\,N., and J.\,P.~Romano. 1992. A circular block-resampling procedure for stationary data. Wiley.
\item Sharpe, W. 1964. Capital asset prices. \emph{JF} 19(3): 425--442.
\end{itemize}


% ===== begin ch3 =====
\chapter{Mean--Variance Portfolios that Ignore Estimation Error}
\label{ch:3}

\begin{plainbox}
This is the ``textbook'' theory of portfolio choice (Markowitz, Tobin, Sharpe).
We \emph{pretend} that we know the true expected returns $\mub$ and the true
covariance matrix $V$. Under this assumption we find the best possible
portfolios. Chapter 4 then asks what happens when $\mub$ and $V$ must be
estimated from data (answer: things go badly wrong).
\end{plainbox}

\section{What we want to do in Chapter 3}
\begin{itemize}
  \item Derive the \textbf{efficient frontier} \emph{without} a risk-free asset (Markowitz).
  \item Derive it \emph{with} a risk-free asset (Tobin): the \textbf{tangency portfolio}.
  \item Understand \textbf{diversification}, the role of covariances, and where the
    \textbf{beta} relation (CAPM-like) comes from.
  \item Extensions: optimisation relative to a \textbf{benchmark}, extra
    \textbf{constraints}, \textbf{transaction costs}.
\end{itemize}

% ---------------------------------------------------------------------
\section{Universe selection}
Before choosing assets we must choose the \textbf{universe} we pick them from.
It depends on:
\begin{itemize}
  \item \textbf{Region}: a French investor may use France, the EU or developed
    markets; regulation may restrict an Indian investor to Indian assets; some
    focus on emerging markets.
  \item \textbf{Breadth}: large caps only, riskier firms only, or everything.
    Norway's sovereign fund (NBIM) holds over 10{,}000 companies, about 1.5\% of
    all listed shares in the world, across equities, fixed income, real estate
    and infrastructure.
  \item \textbf{Asset class}: equities, fixed income, derivatives (futures,
    options), commodities (gold, oil, corn), real estate (REITs), crypto.
  \item \textbf{Risk preferences} (treasuries vs.\ high credit risk),
    \textbf{liquidity} (insurers can hold illiquid assets), views on
    \textbf{inflation and growth}.
\end{itemize}
\textbf{In this course the universe is equities} (mostly US evidence, but the
results apply to other developed markets too).

% ---------------------------------------------------------------------
\section{Portfolio choice in pictures: the Nobel-prize ideas}
All pictures are in the mean ($\mu$) vs.\ standard deviation ($\sigma$) plane.
Investors want to go \emph{up} (more return) and \emph{left} (less risk).
\begin{center}
\begin{tabularx}{\textwidth}{@{}lX@{}}
\toprule
Who & Idea \\
\midrule
\textbf{Markowitz} (1990) & Do not evaluate assets in isolation: look at their
\emph{correlations}. Combining assets gives a \textbf{frontier} that lies to the
left of individual assets. More assets $\Rightarrow$ a better (further left) frontier. \\
\textbf{Tobin} (1981) & Add a risk-free asset: every efficient portfolio is a mix
of (1) the risk-free asset and (2) one portfolio of risky assets, the
\textbf{tangency portfolio} $T$. This is \textbf{two-fund separation}. \\
\textbf{Sharpe} (1990) & In equilibrium the tangency portfolio is the
\textbf{market portfolio} ($T=M$) $\Rightarrow$ CAPM: \emph{only systematic risk is priced}. \\
\textbf{Samuelson} (1970) & Multi-period: if $R_f$ and the slope of the
tangency line (the Sharpe ratio) do not change, the optimal portfolio does not change. \\
\textbf{Merton} (1997) & If $R_f$ and the Sharpe ratio change over time, the
optimal portfolio $=$ Markowitz portfolio $+$ \textbf{intertemporal hedging}
demands against these changes. \\
\bottomrule
\end{tabularx}
\end{center}

\begin{plainbox}
\textbf{The picture to keep in mind} (draw it yourself: it is the most important
graph of the course). Horizontal axis: volatility $\sigma$; vertical axis:
expected return $\mu$.
\begin{itemize}
  \item Single assets are scattered dots.
  \item The \textbf{frontier} is a curve shaped like a ``$\subset$'' lying on its side
    (a hyperbola), to the \emph{left} of the dots: combining assets gives less risk
    for the same return. Its leftmost point is the \textbf{GMV} portfolio; the
    upper half is the \textbf{efficient frontier}.
  \item From the risk-free rate $R_f$ on the vertical axis, draw the steepest
    straight line that still touches the frontier: it touches at the
    \textbf{tangency portfolio} $T$. This line is the \textbf{capital market line};
    its slope is the maximum Sharpe ratio.
  \item Every investor holds a point on this line: closer to $R_f$ if
    risk-averse, beyond $T$ (borrowing) if risk-tolerant.
\end{itemize}
\end{plainbox}

% ---------------------------------------------------------------------
\section{Portfolio diversification}

\subsection{Systematic vs.\ unsystematic risk}
\begin{itemize}
  \item \textbf{Systematic} risk: common to many assets $\Rightarrow$ cannot be diversified away.
  \item \textbf{Unsystematic} (idiosyncratic) risk: specific to one asset $\Rightarrow$ diversifiable.
  \item Assume (for now; revisited in Chapter 8) that unsystematic risk earns no
    reward $\Rightarrow$ hold diversified portfolios $\Rightarrow$ only
    systematic risk matters and only it is compensated.
\end{itemize}

\subsection{Mean and variance of a portfolio}
With dollar amounts $V_1,\dots,V_N$ and total $V=\sum V_n>0$, the weights are
$w_n=V_n/V$ with $\sum_n w_n=\w\T\one=1$ (a portfolio with $V=0$ is an
\emph{arbitrage portfolio}).
\begin{keybox}[: portfolio moments]
\begin{gather*}
R_p=\w\T\R,\qquad
\mu_p=\w\T\mub=\sum_n w_n\E[R_n],\\
\sigma_p^2=\w\T V\w=\sum_{n=1}^N\sum_{m=1}^N w_nw_m\sigma_{nm},\qquad
\sigma_p=\sqrt{\sigma_p^2}.
\end{gather*}
Two assets: $\sigma_p^2=w_1^2\sigma_1^2+w_2^2\sigma_2^2+2w_1w_2\sigma_{12}$.
\end{keybox}

\begin{plainbox}
\textbf{The ``table trick''.} Write a table with $w_1R_1,\dots,w_NR_N$ on the
rows and on the columns. Cell $(n,m)$ is $w_nw_m\sigma_{nm}$. The portfolio
variance is the \textbf{sum of all the cells}. The diagonal cells are variances
($N$ of them); the off-diagonal cells are covariances ($N^2-N$ of them). So the
variance of a sum is \emph{not} the sum of the variances: you need the covariances.
\end{plainbox}

\begin{center}
\begin{tabular}{c|cccc}
 & $w_1R_1$ & $w_2R_2$ & $\cdots$ & $w_NR_N$\\\hline
$w_1R_1$ & $w_1^2\sigma_1^2$ & $w_1w_2\sigma_{12}$ & $\cdots$ & $w_1w_N\sigma_{1N}$\\
$w_2R_2$ & $w_2w_1\sigma_{21}$ & $w_2^2\sigma_2^2$ & $\cdots$ & $w_2w_N\sigma_{2N}$\\
$\vdots$ & $\vdots$ & $\vdots$ & $\ddots$ & $\vdots$\\
$w_NR_N$ & $w_Nw_1\sigma_{N1}$ & $w_Nw_2\sigma_{N2}$ & $\cdots$ & $w_N^2\sigma_N^2$
\end{tabular}
\end{center}


\begin{workedex}[: the minimum-variance mix of two assets]
With two assets and weight $w$ in asset 1,
$\sigma_p^2(w)=w^2\sigma_1^2+(1-w)^2\sigma_2^2+2w(1-w)\sigma_{12}$. Setting the
derivative to zero gives
\[
w^*=\frac{\sigma_2^2-\sigma_{12}}{\sigma_1^2+\sigma_2^2-2\sigma_{12}} .
\]
Example: $\sigma_1=20\%$, $\sigma_2=30\%$, $\rho=0.2$, so $\sigma_{12}=0.012$:
$w^*=(0.09-0.012)/(0.04+0.09-0.024)=0.736$, and $\sigma_p=18.1\%$, \emph{lower than
the safer asset alone} (20\%). Diversification creates a portfolio less risky
than any of its components, as long as $\rho$ is not too high.
(If $\rho=1$ there is no benefit; if $\rho=-1$ risk can be eliminated completely.)
\end{workedex}
\subsection{Diversification with $N$ assets}
Equal weights $w_n=1/N$: each of the $N$ variance cells is $\frac{1}{N^2}\sigma_n^2$;
each of the $N^2-N$ covariance cells is $\frac1{N^2}\sigma_{nm}$. Adding them up:
\begin{keybox}[: the diversification formula]
\[
\sigma_p^2=\frac1N\,\underbrace{\Big(\frac1N\sum_n\sigma_n^2\Big)}_{\text{average variance}}
+\Big(1-\frac1N\Big)\underbrace{\Big(\frac{1}{N^2-N}\sum_n\sum_{m\ne n}\sigma_{nm}\Big)}_{\text{average covariance}}
\ \xrightarrow{N\to\infty}\ \text{average covariance}.
\]
For a well-diversified portfolio, \textbf{covariances determine risk; individual
variances barely matter}. When adding an asset, look at its covariance
(correlation) with the portfolio, not at its variance.
\end{keybox}
\begin{plainbox}
Each stock's own ``noise'' gets weight $1/N^2$ and there are only $N$ of them, so
in total they count $1/N$ and vanish. But the ``common movement'' (all stocks go
down when the market goes down) is in all $N^2-N$ pairs and never disappears.
That floor is market risk.
\end{plainbox}

\subsection*{Python exercise 1: diversification in practice}
500 stocks, 10 years monthly; each return $=$ market $\times$ beta (0.5--1.5) $+$ firm-specific noise.
\begin{lstlisting}
import numpy as np
import pandas as pd
rng = np.random.default_rng(31)
N_all, T = 500, 120
beta = rng.uniform(0.5, 1.5, N_all)
f = rng.normal(0.006, 0.045, T)                  # market
e = rng.normal(0, 0.09, size=(T, N_all))         # firm-specific
R = 0.002 + np.outer(f, beta) + e                # stock returns
V = np.cov(R, rowvar=False)
avg_var = np.mean(np.diag(V))
avg_cov = (V.sum() - np.trace(V)) / (N_all*(N_all - 1))   # off-diagonal average

Ns = [1, 2, 5, 10, 25, 50, 100, 500]
vol = {N: np.mean([R[:, rng.choice(N_all, N, replace=False)].mean(axis=1).std(ddof=1)
                   for _ in range(200)]) for N in Ns}
formula = {N: np.sqrt(avg_var/N + (1 - 1/N)*avg_cov) for N in Ns}
table = pd.DataFrame({"simulated vol": [np.sqrt(12)*vol[N] for N in Ns],
                      "formula vol": [np.sqrt(12)*formula[N] for N in Ns]}, index=Ns)
print(table.round(3))
print("Limit (sqrt of average covariance):", round(float(np.sqrt(12*avg_cov)), 3))
\end{lstlisting}
\begin{lstlisting}[style=out]
N    simulated vol  formula vol
1            0.358        0.359
2            0.279        0.281
5            0.221        0.221
10           0.197        0.197
25           0.181        0.181
50           0.175        0.176
100          0.173        0.173
500          0.170        0.170
Limit (sqrt of average covariance): 0.17
\end{lstlisting}
Volatility falls from 36\% (one stock) to 20\% (10 stocks) but only to 17\% (500
stocks). \textbf{The 17\% floor is $\sqrt{\text{average covariance}}$}:
diversification removes firm-specific risk, not market risk. Most of the benefit
comes from the first 10--20 stocks.

% ---------------------------------------------------------------------
\section{Mean--variance frontier without a risk-free asset (focus)}

\begin{definitionbox}{Frontier portfolio}
Given a target expected return, the portfolio with the smallest risk is a
\textbf{frontier portfolio}. The set of all of them in the $(\sigma,\mu)$ plane
is the \textbf{portfolio frontier}; its upper half is the \textbf{efficient frontier}.
\end{definitionbox}

\paragraph{Notation.} $\one$ vector of ones; $I_N$ identity; $\w$ weights;
$\mub=(\E[R_1],\dots,\E[R_N])\T$; $V$ the symmetric $N\times N$ covariance
matrix (often written $\Sigma$ in papers).

\subsection{The optimisation problem}
\[
\min_{\w}\ \w\T V\w\qquad\text{s.t.}\quad (1)\ \w\T\one=1,\qquad (2)\ \w\T\mub=\mu_{targ}.
\]
\begin{plainbox}
Find the weights with the lowest variance such that (1) they add up to 100\% and
(2) the portfolio earns the target return, e.g.\ 15\%. It is a quadratic
objective with linear constraints, so it has a closed-form solution (when there
are no other constraints).
\end{plainbox}

\subsection{Solving it step by step (Lagrangian)}
\textbf{Step 1: Lagrangian.} With multipliers $\lambda_w,\lambda_R$:
\[
L=\w\T V\w+\lambda_w(1-\w\T\one)+\lambda_R(\mu_{targ}-\w\T\mub).
\]
\textbf{Step 2: FOC in $\w$.}
\[
\frac{\partial L}{\partial\w}=2V\w-\lambda_w\one-\lambda_R\mub=\bm 0
\ \Rightarrow\
\w=V^{-1}\Big(\frac{\lambda_w}{2}\one+\frac{\lambda_R}{2}\mub\Big). \tag{1}
\]
\textbf{Step 3: FOCs in the multipliers} give back the constraints:
$1=\w\T\one$ (2) and $\mu_{targ}=\w\T\mub$ (3).

\textbf{Step 4: substitute (1) into (2) and (3).} Define the three scalars
\begin{keybox}[: the three ``a'' numbers]
\[
a_1=\mub\T V^{-1}\mub,\qquad a_2=\mub\T V^{-1}\one,\qquad a_3=\one\T V^{-1}\one,\qquad D=a_1a_3-a_2^2 .
\]
\end{keybox}
Then
\[
\begin{pmatrix}1\\ \mu_{targ}\end{pmatrix}=
\begin{pmatrix}a_3&a_2\\a_2&a_1\end{pmatrix}
\begin{pmatrix}\lambda_w/2\\ \lambda_R/2\end{pmatrix}
\ \Rightarrow\
\frac{\lambda_w}{2}=\frac{a_1-a_2\mu_{targ}}{D},\qquad
\frac{\lambda_R}{2}=\frac{-a_2+a_3\mu_{targ}}{D}.
\]
\textbf{Step 5: plug back into (1).}
\begin{plainbox}
\textbf{What do $a_1,a_2,a_3$ mean?} They are just numbers that summarise the
assets, but each has an interpretation:
\begin{itemize}
  \item $a_3=\one\T V^{-1}\one=1/\sigma^2_{GMV}$: the higher $a_3$, the lower the
    risk of the safest portfolio.
  \item $a_2/a_3=\mu_{GMV}$: the expected return of the safest portfolio.
  \item $D/a_3=a_1-a_2^2/a_3$ measures how much the expected returns differ across
    assets (after accounting for risk): it controls how steep the frontier is.
    If all assets had the same mean, $D=0$ and the frontier would be a single point
    in the $\mu$ direction.
\end{itemize}
The frontier variance is a quadratic function of the target return:
$\sigma_p^2=\frac{1}{a_3}+\frac{a_3}{D}\big(\mu_{targ}-\mu_{GMV}\big)^2$. Each extra
unit of target return above the GMV costs more and more variance.
\end{plainbox}

\begin{keybox}[: frontier weights for target return $\mu_{targ}$]
\[
\w=\frac1D\big(a_1V^{-1}\one-a_2V^{-1}\mub\big)+\frac1D\big(a_3V^{-1}\mub-a_2V^{-1}\one\big)\,\mu_{targ}.
\]
Closed form only without extra constraints. With $w_n\ge0$ or $w_n\le w_{\max}$
we must solve numerically (Chapter 4 appendix).
\end{keybox}

\begin{addnote}
\textbf{How to remember the $2\times2$ inverse.}
$\begin{pmatrix}a&b\\c&d\end{pmatrix}^{-1}=\frac{1}{ad-bc}\begin{pmatrix}d&-b\\-c&a\end{pmatrix}$.
Here $\begin{pmatrix}a_3&a_2\\a_2&a_1\end{pmatrix}^{-1}=\frac1D\begin{pmatrix}a_1&-a_2\\-a_2&a_3\end{pmatrix}$,
which gives the two multipliers directly. $D>0$ as long as the assets do not all have the same mean.
\end{addnote}

\subsection{Two-fund separation and the GMV portfolio}
Define $\w_0$ (frontier portfolio with $\mu_{targ}=0$) and $\w_1$ ($\mu_{targ}=1$):
\[
\w_0=\frac1D\big(a_1V^{-1}\one-a_2V^{-1}\mub\big),\qquad
\w_1=\w_0+\frac1D\big(a_3V^{-1}\mub-a_2V^{-1}\one\big).
\]
\begin{keybox}[: two-fund separation (no risk-free asset)]
Every frontier portfolio is a combination of $\w_0$ and $\w_1$:
\[ \w_p=\w_0+(\w_1-\w_0)\mu_{targ}=(1-\mu_{targ})\w_0+\mu_{targ}\w_1 . \]
Any two frontier portfolios generate the entire frontier.
\end{keybox}
Its variance is
$\sigma_p^2=(1-\mu_{targ})^2\w_0\T V\w_0+2\mu_{targ}(1-\mu_{targ})\w_0\T V\w_1+\mu_{targ}^2\w_1\T V\w_1$;
minimising it over $\mu_{targ}$ gives the \textbf{global minimum-variance (GMV)} portfolio.

\begin{addnote}
\textbf{Closed forms you can derive from the ``a'' numbers} (Merton 1972):
\[
\sigma_p^2(\mu_{targ})=\frac{a_3\mu_{targ}^2-2a_2\mu_{targ}+a_1}{D}
\quad\text{(a parabola in $(\mu,\sigma^2)$, a hyperbola in $(\mu,\sigma)$)},
\]
\[
\w_{GMV}=\frac{V^{-1}\one}{a_3},\qquad \mu_{GMV}=\frac{a_2}{a_3},\qquad \sigma^2_{GMV}=\frac1{a_3}.
\]
(Set $d\sigma_p^2/d\mu_{targ}=0$: $\mu_{targ}=a_2/a_3$; plugging in gives $\sigma^2=1/a_3$.)
\end{addnote}

\subsection{Worked example with $N=3$}
\[
\mub=\begin{pmatrix}0.20\\0.30\\0.40\end{pmatrix},\quad
V=\begin{pmatrix}0.0625&0.0700&0.1050\\0.0700&0.1225&0.0840\\0.1050&0.0840&0.3600\end{pmatrix},\quad
V^{-1}=\begin{pmatrix}85.06&-37.61&-16.03\\-37.61&26.35&4.82\\-16.03&4.82&6.33\end{pmatrix}.
\]
$a_1=0.865$, $a_2=2.40$, $a_3=20.09$, $D=11.63$, and
\[
\w_0=\begin{pmatrix}2.48\\-0.96\\-0.52\end{pmatrix},\quad
\w_1=\begin{pmatrix}-5.18\\4.37\\1.82\end{pmatrix},\quad
\w_1-\w_0=\begin{pmatrix}-7.66\\5.32\\2.34\end{pmatrix},\qquad
\w=\w_0+\mu_{targ}(\w_1-\w_0).
\]
\begin{center}
\begin{tabularx}{\textwidth}{@{}lccX@{}}
\toprule
$\mu_{targ}$ & weights $(w_1,w_2,w_3)$ & $\sigma_p$ & meaning \\
\midrule
35\% & $(-0.20,\ 0.90,\ 0.30)$ & 0.376 & short \$0.20 of asset 1, buy \$0.90 of 2 and \$0.30 of 3 \\
30\% & $(0.18,\ 0.64,\ 0.18)$ & 0.3258 & lower risk than asset 2 alone ($\sigma=0.35$) for the same 30\% \\
GMV ($\approx$11.9\%) & $(1.56,\ -0.32,\ -0.24)$ & 0.223 & lowest possible risk \\
\bottomrule
\end{tabularx}
\end{center}
\textbf{Conclusion:} asset 2 alone gives 30\% with $\sigma=0.35$; diversifying
gives 30\% with $\sigma=0.3258$. To trace the whole frontier, loop over many $\mu_{targ}$.

\begin{lstlisting}
import numpy as np
ones = np.ones((3, 1))
mu = np.array([[0.20], [0.30], [0.40]])
V = np.array([[0.0625, 0.0700, 0.1050],
              [0.0700, 0.1225, 0.0840],
              [0.1050, 0.0840, 0.3600]])
Vinverse = np.linalg.inv(V)
a1 = mu.T @ Vinverse @ mu          # ".T" transpose, "@" matrix product
a2 = mu.T @ Vinverse @ ones
a3 = ones.T @ Vinverse @ ones
D = a1*a3 - a2**2                  # 0.8650, 2.3980, 20.0930, 11.6306
w0 = (a1*Vinverse@ones - a2*Vinverse@mu)/D
w1 = w0 + (a3*Vinverse@mu - a2*Vinverse@ones)/D

def solution(mu_targ):
    weights = w0 + mu_targ*(w1 - w0)
    vol = np.sqrt(weights.T @ V @ weights)[0, 0]
    return weights, vol
print(solution(0.35))   # [-0.2034, 0.9067, 0.2966], vol 0.3764
print(solution(0.30))   # [ 0.1797, 0.6407, 0.1797], vol 0.3258
\end{lstlisting}
\begin{addnote}
Column vectors \texttt{np.ones((3,1))} produce $1\times1$ \emph{matrices} for
$a_1,a_2,a_3$ (hence \texttt{[0,0]} to extract the number). With 1-D arrays
(\texttt{np.ones(3)}) they are plain numbers; both are fine, but be consistent.
\end{addnote}

% ---------------------------------------------------------------------
\section{Portfolio frontier with a risk-free asset}

\paragraph{Two views of the same portfolio.} \$40 risk-free, \$30 IBM, \$30 Merck.
View 1: 40\% $R_f$, 30\% IBM, 30\% Merck. View 2: 40\% in the risk-free
sub-portfolio, 60\% in a risky sub-portfolio that is 50/50 IBM/Merck.

If a fraction $a$ goes to a risky portfolio $q$ and $1-a$ to $R_f$:
$\mu_p=(1-a)R_f+a\E[R_q]$, $\sigma_p=|a|\sigma_q$ $\Rightarrow$ a \textbf{straight
line} from $R_f$ through $q$. The steepest line touches the frontier at the
\textbf{tangency portfolio}.

\subsection{The problem}
With $w_f=1-\w\T\one$ in the risk-free asset:
\[
\min_{\w}\ \w\T V\w\qquad\text{s.t.}\qquad \w\T\mub+(1-\w\T\one)R_f=\mu_{targ}
\ \Longleftrightarrow\ \w\T(\mub-R_f\one)=\mu_{targ}-R_f .
\]
\begin{keybox}[: frontier weights with a risk-free asset]
\[
\w=\underbrace{\left[\frac{\mu_{targ}-R_f}{(\mub-R_f\one)\T V^{-1}(\mub-R_f\one)}\right]}_{\text{scaling: how much risk}}
\ \underbrace{V^{-1}(\mub-R_f\one)}_{\text{direction: which risky portfolio}} . \tag{8}
\]
The higher $\mu_{targ}$, the larger the position in the risky portfolio and
the smaller in $R_f$. The scalar plays the role of risk tolerance $1/\gamma$.
\end{keybox}

\begin{addnote}
\textbf{Derivation (one constraint only).} Let $\bm e=\mub-R_f\one$ and
$H=\bm e\T V^{-1}\bm e$. Lagrangian $L=\w\T V\w+\lambda(\mu_{targ}-R_f-\w\T\bm e)$.
FOC: $2V\w=\lambda\bm e\Rightarrow\w=\frac\lambda2V^{-1}\bm e$. Plug into the
constraint: $\frac\lambda2H=\mu_{targ}-R_f\Rightarrow\frac\lambda2=\frac{\mu_{targ}-R_f}{H}$.
\textbf{Bonus:} $\sigma_p=\frac{\mu_{targ}-R_f}{\sqrt H}$, so
$\frac{\mu_{targ}-R_f}{\sigma_p}=\sqrt H$: every portfolio on the line has the
same Sharpe ratio $\sqrt H=\sqrt{(\mub-R_f\one)\T V^{-1}(\mub-R_f\one)}$, the
\textbf{maximum Sharpe ratio}.
\end{addnote}

\subsection{Worked example: three risky assets $+$ $R_f=6.5\%$, $\mu_{targ}=30\%$}
$\mub-R_f\one=(0.135,0.235,0.335)\T$, $H=0.6382$, $V^{-1}(\mub-R_f\one)=(-2.7268,\ 2.7299,\ 1.0889)\T$:
\[
\w=\frac{0.30-0.065}{0.6382}\begin{pmatrix}-2.7268\\2.7299\\1.0889\end{pmatrix}
=\begin{pmatrix}-1.0041\\1.0052\\0.4010\end{pmatrix},
\]
\[
w_f=1-(-1.0041+1.0052+0.4010)=0.5979,\qquad \sigma_p=0.2942 .
\]
\begin{center}
\begin{tabular}{@{}lccc@{}}
\toprule
Case for $\mu_p=30\%$ & Asset 2 alone & 3 risky assets & 3 risky $+$ $R_f$\\
\midrule
$\sigma_p$ & 0.3500 & 0.3256 & 0.2942 \\
\bottomrule
\end{tabular}
\end{center}
\textbf{More assets $\Rightarrow$ same return with less risk.} Check:
$\sigma_p=(0.30-0.065)/\sqrt{0.6382}=0.235/0.7989=0.2942$ \checkmark.


\begin{keybox}[: capital market line]
All efficient portfolios with a risk-free asset lie on the line
\[
\mu_p=R_f+SR_T\,\sigma_p,\qquad SR_T=\frac{\E[R_T]-R_f}{\sigma_T}=\sqrt{(\mub-R_f\one)\T V^{-1}(\mub-R_f\one)} .
\]
In the example: $\mu_p=0.065+0.799\,\sigma_p$. With $\sigma_p=0.2942$ we get
$\mu_p=0.30$ \checkmark. An investor who wants 40\% needs $\sigma_p=(0.40-0.065)/0.799=0.419$
and will borrow at $R_f$ to lever the tangency portfolio.
\end{keybox}
\paragraph{Tangency portfolio.} Rescale the risky weights to sum to one:
\[
\w_{tang}=\frac{\w}{\w\T\one}=\frac{1}{0.4021}\begin{pmatrix}-1.0041\\1.0052\\0.4010\end{pmatrix}
=\begin{pmatrix}-2.4971\\2.4998\\0.9972\end{pmatrix}.
\]
It depends on $R_f$ but \textbf{not} on $\mu_{targ}$. Its Sharpe ratio is $\sqrt H=0.799$.

\begin{lstlisting}
Rf = 0.065
x = Vinverse @ (mu - Rf*ones)                 # direction V^{-1}(mu - Rf 1)
def solutionWithRf(mu_targ, Rf):
    H = ((mu - Rf*ones).T @ Vinverse @ (mu - Rf*ones))[0, 0]
    w = (mu_targ - Rf)/H * x
    return w, np.sqrt(w.T @ V @ w)[0, 0]
w, vol = solutionWithRf(0.30, 0.065)          # [-1.0041, 1.0052, 0.4010], 0.2942
w_f = 1 - w.sum()                             # 0.5979
weightsTangency = x / (ones.T @ x)            # [-2.4972, 2.4999, 0.9972]
\end{lstlisting}

\subsection*{Python exercise 2: the estimated tangency portfolio}
10 assets with known $\mub$, $V$; estimate the tangency portfolio from 10 years of data.
\begin{lstlisting}
import numpy as np
import pandas as pd
rng = np.random.default_rng(32)
N, T = 10, 120
mu = np.linspace(0.04, 0.12, N) / 12                     # true excess means
beta = np.linspace(0.6, 1.4, N)
V = 0.045**2*np.outer(beta, beta) + np.diag(np.full(N, 0.06**2))
x = np.linalg.solve(V, mu); w_true = x / x.sum()         # true tangency
R = rng.multivariate_normal(mu, V, size=T)
mu_hat, V_hat = R.mean(axis=0), np.cov(R, rowvar=False)
x = np.linalg.solve(V_hat, mu_hat); w_hat = x / x.sum()  # estimated tangency
sr = lambda w, m, S: np.sqrt(12) * (w @ m) / np.sqrt(w @ S @ w)
print("true max SR:", round(sr(w_true, mu, V), 2))
print("estimated, in-sample SR:", round(sr(w_hat, mu_hat, V_hat), 2))
print("estimated, true SR:", round(sr(w_hat, mu, V), 2))
\end{lstlisting}
\begin{lstlisting}[style=out]
true weights      -0.02  0.01  0.04  0.06  0.09  0.11  0.14  0.16  0.19  0.22
estimated weights -0.00 -0.21  0.40 -0.38 -0.18  0.83 -0.01  0.35  0.05  0.15
true max SR: 0.49   estimated, in-sample SR: 0.95   estimated, true SR: 0.34
\end{lstlisting}
\begin{keybox}[: promise vs.\ reality]
\begin{itemize}
  \item \textbf{True SR} (true weights, true parameters): 0.49 --- the best possible.
  \item \textbf{In-sample SR} (estimated weights, \emph{estimated} parameters): 0.95 --- \textbf{overfitting}: what the investor \emph{believes}.
  \item \textbf{True SR of the estimated portfolio} (estimated weights, \emph{true} parameters): 0.34 --- what the investor \emph{gets}.
\end{itemize}
The estimated weights look nothing like the true ones ($-0.38$ vs.\ $+0.06$ for
asset 4). This gap is the subject of Chapter 4.
\end{keybox}

\subsection{Equivalent approach: maximise mean--variance utility}
\begin{keybox}[: MVU formulation]
\begin{gather*}
\max_{\w}\ \mu_p-\frac\gamma2\sigma_p^2=R_f+\w\T(\mub-R_f\one)-\frac\gamma2\w\T V\w\\
\Rightarrow\quad
(\mub-R_f\one)-\gamma V\w=\bm 0
\quad\Rightarrow\quad
\boxed{\w=\frac1\gamma V^{-1}(\mub-R_f\one)} . \tag{10}
\end{gather*}
Varying $\gamma$ traces the whole frontier. Equations (8) and (10) coincide if
\[
\frac1\gamma=\frac{\mu_{targ}-R_f}{(\mub-R_f\one)\T V^{-1}(\mub-R_f\one)} .
\]
\end{keybox}
\begin{plainbox}
\textbf{Slide question: why do both sides represent the same economic quantity?}
Both measure how much risk the investor is willing to take. Choosing a high
target return $\mu_{targ}$ is the same as being risk tolerant (small $\gamma$):
in both cases you move further out along the capital market line. In the
example, $\mu_{targ}=30\%$ corresponds to $\gamma=0.6382/0.235=2.7157$ --- exactly
the value used in the Chapter 4 appendix code.
\end{plainbox}

\paragraph{The tangency portfolio does not depend on $\gamma$.}
\[
\w_{tang}=\frac{V^{-1}(\mub-R_f\one)}{\one\T V^{-1}(\mub-R_f\one)} .
\]
$\gamma$ cancels in numerator and denominator $\Rightarrow$ \textbf{all investors
hold risky assets in the same proportions}, whatever their risk aversion. In
equilibrium demand $=$ supply, so the tangency portfolio is the
\textbf{market portfolio}.

\paragraph{Example (used again in Chapter 4).} $\gamma=5$, $R_f=1.01$, two assets with
gross mean 1.09, vol 20\%, correlation 0.99: MV weights $(0.201,0.201)$,
$w_f=0.598$, tangency $(0.5,0.5)$.

\subsection{The market portfolio}
\[
w_n^{mkt}=\frac{\text{market cap of }n}{\text{total market cap}}=\frac{num_n\times P_n}{\sum_{m=1}^N num_m\times P_m}.
\]
\begin{itemize}
  \item If the CAPM holds, hold only the market (plus $R_f$): \textbf{passive investing};
    little benefit from active funds.
  \item In practice one holds a \textbf{cap-weighted index} that tracks the market
    (origin of index funds and ETFs). Benefits: broad exposure, little trading
    (low costs), large firms are liquid, high capacity.
  \item Limitations: the CAPM may not hold; other factors (size, momentum) may
    earn more $\Rightarrow$ factor indexes, other assets (Chapters 5--8).
\end{itemize}

% ---------------------------------------------------------------------
\section{Risk--return of individual assets and frontier portfolios (focus)}
Questions: how does one asset contribute to portfolio risk? What exactly is
systematic risk? How is it related to expected return? Four steps (with a risk-free asset).

\begin{keybox}[: four marginal contributions]
\begin{enumerate}
  \item \textbf{Expected return.} $\E[R_p]=R_f+\sum_n w_n(\E[R_n]-R_f)$
    $\Rightarrow$ $\dfrac{\partial\E[R_p]}{\partial w_n}=\E[R_n]-R_f$ (the risk premium).
  \item \textbf{Total contribution to variance.} Row $n$ of the table:
    $\sum_m w_nw_m\sigma_{nm}=w_n\sigma_{np}$, and $\sum_n w_n\sigma_{np}=\sigma_p^2$.
  \item \textbf{Marginal contribution to variance.}
    $\dfrac{\partial\sigma_p^2}{\partial w_n}=2w_n\sigma_n^2+2\sum_{m\ne n}w_m\sigma_{nm}=2\sum_m w_m\sigma_{nm}=2\,\text{Cov}[R_n,R_p]=2\sigma_{np}$.
  \item \textbf{Marginal contribution to volatility.} From $\frac{\partial\sigma_p^2}{\partial w_n}=2\sigma_p\frac{\partial\sigma_p}{\partial w_n}$:
    $\dfrac{\partial\sigma_p}{\partial w_n}=\dfrac{\sigma_{np}}{\sigma_p}$.
\end{enumerate}
In matrix form: the vector of covariances with the portfolio is $V\w$.
\end{keybox}
\begin{plainbox}
An asset's risk inside a portfolio is \textbf{not} its own volatility; it is how
much it \emph{moves with the portfolio} ($\sigma_{np}/\sigma_p$). A very volatile
asset that is uncorrelated with the portfolio adds almost no risk at the margin.
\end{plainbox}

\subsection*{Python exercise 3: where does the risk of a 60/40 portfolio come from?}
\begin{lstlisting}
import numpy as np
import pandas as pd
rng = np.random.default_rng(33)
T = 240
vol = np.array([0.16, 0.06]) / np.sqrt(12)          # stocks, bonds
C = np.array([[1, 0.1], [0.1, 1]])
R = rng.multivariate_normal([0.006, 0.003], C*np.outer(vol, vol), size=T)
V = np.cov(R, rowvar=False)
w = np.array([0.6, 0.4])
var_p = w @ V @ w
cov_np = V @ w                       # covariance of each asset with the portfolio
marginal = cov_np / np.sqrt(var_p)   # d sigma_p / d w_n
share = w * cov_np / var_p           # share of portfolio risk (sums to 1)
print(pd.DataFrame({"weight": w, "marginal": marginal, "risk share": share},
                   index=["stocks", "bonds"]).round(3))
\end{lstlisting}
\begin{lstlisting}[style=out]
        weight  marginal  risk share
stocks     0.6     0.044       0.925
bonds      0.4     0.005       0.075
\end{lstlisting}
\textbf{Stocks are 60\% of the money but 92.5\% of the risk}: a ``balanced'' 60/40
portfolio is almost a pure bet on stocks. Risk-parity portfolios (Chapter 6)
equalise risk shares.

\subsection{The beta relation}
\begin{definitionbox}{Return-to-risk ratio (RRR)}
\[
RRR_n=\frac{\text{marginal return}}{\text{marginal risk}}=\frac{\partial\E[R_p]/\partial w_n}{\partial\sigma_p/\partial w_n}=\frac{\E[R_n]-R_f}{\sigma_{np}/\sigma_p}.
\]
\end{definitionbox}
\textbf{Claim:} for any frontier portfolio $p$, every asset has the same RRR, equal to the portfolio's:
\[
\frac{\E[R_n]-R_f}{\sigma_{np}/\sigma_p}=\frac{\E[R_p]-R_f}{\sigma_{pp}/\sigma_p}=\frac{\E[R_p]-R_f}{\sigma_p}.
\]
\begin{intuition}
If one asset had a higher return per unit of marginal risk, you could buy a
bit more of it and sell a bit of another: the portfolio's return would rise for
the same risk, so the portfolio would not have been on the frontier. At the
optimum, all marginal ``prices of risk'' are equalised (like equalising marginal
utility per euro across goods in microeconomics).
\end{intuition}
\begin{keybox}[: beta pricing relation]
Rearranging,
\[
\E[R_n]-R_f=\beta_{np}\big(\E[R_p]-R_f\big),\qquad \beta_{np}=\frac{\sigma_{np}}{\sigma_p^2}=\frac{\text{Cov}[R_n,R_p]}{\text{Var}[R_p]},
\]
for any frontier portfolio $p$ (other than $R_f$).
$\beta_{np}$ = amount of systematic risk; $\E[R_p]-R_f$ = price per unit of risk.
\textbf{The CAPM is this relation with $p=$ market}; it holds only if the market
is on the frontier.
\end{keybox}
\begin{addnote}
\textbf{Quick proof with the tangency portfolio.} $\w_T\propto V^{-1}\bm e$ means
$V\w_T=c\,\bm e$ for some $c>0$, i.e.\ $\sigma_{nT}=c\,(\E[R_n]-R_f)$ for every $n$.
Multiply by $w_{T,n}$ and sum: $\sigma_T^2=c\,(\E[R_T]-R_f)$. Divide the two:
$\frac{\sigma_{nT}}{\sigma_T^2}=\frac{\E[R_n]-R_f}{\E[R_T]-R_f}$, i.e.\
$\E[R_n]-R_f=\beta_{nT}(\E[R_T]-R_f)$. $\square$
\end{addnote}

\subsection*{Python exercise 4: checking the beta relation}
\begin{lstlisting}
import numpy as np
import pandas as pd
N = 5
mu = np.array([0.05, 0.06, 0.08, 0.09, 0.11])          # excess means
vol = np.array([0.12, 0.15, 0.18, 0.22, 0.25])
C = 0.3*np.ones((N, N)) + 0.7*np.eye(N)
V = C * np.outer(vol, vol)
x = np.linalg.solve(V, mu)
w = {"tangency": x / x.sum(), "1/N": np.full(N, 1/N)}
beta = {p: V @ w[p] / (w[p] @ V @ w[p]) for p in w}    # beta of each asset w.r.t. p
prem = {p: w[p] @ mu for p in w}                       # premium of p
print(pd.DataFrame({"E[R]-Rf": mu,
                    "beta x prem (tangency)": beta["tangency"]*prem["tangency"],
                    "beta x prem (1/N)": beta["1/N"]*prem["1/N"]}).round(4))
\end{lstlisting}
\begin{lstlisting}[style=out]
   E[R]-Rf  beta x prem (tangency)  beta x prem (1/N)
1     0.05                    0.05             0.0443
2     0.06                    0.06             0.0586
3     0.08                    0.08             0.0743
4     0.09                    0.09             0.0971
5     0.11                    0.11             0.1157
\end{lstlisting}
Exact for the tangency portfolio (on the frontier); not for 1/N (not on the frontier).

% ---------------------------------------------------------------------
\section{Portfolio optimisation with respect to a benchmark}
Fund managers are judged against a \textbf{benchmark} (the investor's
alternative is a cheap index fund). Because mean returns are noisy, investors
focus on \textbf{tracking error volatility}.
\begin{keybox}[: tracking error]
\begin{gather*}
TE=R_A-R_B=(\w_A-\w_B)\T\R,\qquad
\E[TE]=(\w_A-\w_B)\T\mub,\\
\text{Vol}(TE)=\sqrt{(\w_A-\w_B)\T V(\w_A-\w_B)} .
\end{gather*}
Manager's problem (Roll 1992, \textbf{TEV criterion}):
\[
\max_{\w_A}\ (\w_A-\w_B)\T\mub\quad\text{s.t.}\quad\sqrt{(\w_A-\w_B)\T V(\w_A-\w_B)}\le\sigma_{\max},\quad \w_A\T\one=1 .
\]
\end{keybox}
Same as Markowitz, but in ``active'' space: active return instead of total
return, tracking-error volatility instead of total volatility.

\begin{center}
\begin{tabularx}{\textwidth}{@{}lXX@{}}
\toprule
Ratio & Formula & Counterpart of \\
\midrule
\textbf{Information ratio} & $\dfrac{(\w_A-\w_B)\T\mub}{\sqrt{(\w_A-\w_B)\T V(\w_A-\w_B)}}=\dfrac{\text{expected TE}}{\text{TE volatility}}$ & Sharpe ratio (benchmark instead of $R_f$) \\
\textbf{Appraisal ratio} & $\dfrac{\alpha}{\sigma(\varepsilon)}$: alpha from a factor regression over the residual volatility & Treynor-like numerator, but risk = unexplained risk \\
\bottomrule
\end{tabularx}
\end{center}

\subsection*{Python exercise 5: tracking error, IR and appraisal ratio}
Fund with true alpha 2.4\%/yr and beta 1.4 vs.\ a benchmark with beta 1.
\begin{lstlisting}
import numpy as np
rng = np.random.default_rng(35)
T = 240
m = rng.normal(0.005, 0.045, T)       # benchmark excess return
e = rng.normal(0, 0.01, T)            # manager's own risk
a = 0.002 + 1.4*m + e                 # fund excess return
active = a - m
te_vol = active.std(ddof=1)
ir = np.sqrt(12) * active.mean() / te_vol
beta, alpha = np.polyfit(m, a, 1)     # slope, intercept of the regression
resid_vol = (a - alpha - beta*m).std(ddof=2)
print("IR:", round(ir, 2), " Appraisal:", round(np.sqrt(12)*alpha/resid_vol, 2))
\end{lstlisting}
\begin{lstlisting}[style=out]
Beta: 1.39  alpha (annual): 0.0263
Residual volatility (annual): 0.0319   Tracking-error volatility (annual): 0.071
Information ratio: 1.04   Appraisal ratio: 0.82
\end{lstlisting}
TE volatility (7.1\%) is more than twice the residual volatility (3.2\%) because
it includes the extra beta. The IR (1.04) rewards the fund also for the extra
market exposure; the \textbf{appraisal ratio isolates skill} (0.82).
\begin{addnote}
\texttt{np.polyfit(x, y, 1)} fits $y=b\,x+a$ and returns \texttt{[b, a]}
(slope first!). \texttt{ddof=2} because two parameters were estimated.
\end{addnote}

% ---------------------------------------------------------------------
\section{Other constraints and transaction costs}
\paragraph{Other constraints.} Add to the problem: (3) $w_n\ge w_{\min}$
(long-only if $w_{\min}=0$), (4) $w_n\le w_{\max}$. No closed form; solve with a
numerical optimiser (\texttt{scipy.optimize.minimize}, Chapter 4 appendix).

\paragraph{Three ways to handle transaction costs.}
\begin{enumerate}
  \item Ignore them in the optimisation, subtract them from the returns:
    $r^{net}_{t+1}=R_{t+1}-\sum_n|w_{t+1,n}-w_{t^+,n}|\times TC$.
  \item Include them in a \textbf{myopic} optimisation (theoretically incorrect but easy):
    \[ \max_{\w}\ \w\T\mub-\frac1\gamma\w\T V\w-TC\,\|\w-\w_0\|_1\quad\text{s.t. }\sum_n w_n=1, \]
    with $\w_0$ the portfolio before trading and $\|\w-\w_0\|_1=\sum_n|w_n-w_{0,n}|$.
  \item Include them in a \textbf{dynamic} optimisation (correct but hard;
    DeMiguel--Uppal 2005; not in the course).
\end{enumerate}
\begin{plainbox}
The $\|\w-\w_0\|_1$ penalty makes it costly to move away from what you already
hold, so the optimiser trades only when the expected gain beats the cost. It
creates a ``no-trade zone'' around the current portfolio.
\end{plainbox}

% ---------------------------------------------------------------------
\section{Summary of modern portfolio theory}
\begin{enumerate}
  \item Investors hold \textbf{frontier portfolios}; more assets improve the frontier.
  \item With a risk-free asset, frontier portfolios are combinations of $R_f$ and
    the \textbf{tangency portfolio} (two-fund separation).
  \item An asset's risk is its \textbf{contribution to portfolio risk}: covariances
    matter, individual variances barely do.
  \item For any asset, risk $=$ beta with respect to a frontier portfolio; reward
    per unit of risk $=$ that portfolio's premium.
\end{enumerate}
\textbf{Libraries:} \texttt{PyPortfolioOpt} (MV optimisation, Black--Litterman),
\texttt{Riskfolio-Lib}; EDHEC's Coursera code.


\begin{mistakes}
\begin{enumerate}
  \item Writing the portfolio variance as $\sum w_n^2\sigma_n^2$ (forgetting the covariances).
  \item Confusing $V^{-1}\one$ (raw GMV weights) with $V^{-1}\mub$; or forgetting to
    normalise GMV and tangency weights so that they sum to one.
  \item Thinking the tangency portfolio depends on $\gamma$ or on $\mu_{targ}$: it
    depends only on $\mub$, $V$ and $R_f$.
  \item Using gross returns and a net risk-free rate in the same formula: $\mub-R_f\one$
    must be both gross or both net (the difference is the same).
  \item Reading an asset's own volatility as its contribution to portfolio risk:
    the contribution is $\sigma_{np}/\sigma_p$.
  \item In code, mixing column vectors \texttt{(N,1)} and flat arrays \texttt{(N,)}:
    results become matrices of the wrong shape. Pick one convention.
\end{enumerate}
\end{mistakes}
\section{Self-test questions}
\begin{enumerate}
  \item Two assets, $\sigma_1=0.2$, $\sigma_2=0.3$, $\rho=0.5$, $w=(0.6,0.4)$. Portfolio vol?
    \emph{($0.36\cdot0.04+0.16\cdot0.09+2\cdot0.24\cdot0.03=0.0432\Rightarrow\sigma=0.208$.)}
  \item Equal-weighted, average variance 0.09, average covariance 0.02, $N=10$.
    Portfolio variance? \emph{($0.009+0.9\cdot0.02=0.027$.)} Limit as $N\to\infty$? \emph{(0.02.)}
  \item Write $a_1,a_2,a_3$ and the GMV weights in terms of them. \emph{($\w_{GMV}=V^{-1}\one/a_3$.)}
  \item Why does $\gamma$ not appear in the tangency portfolio? \emph{(It cancels when normalising.)}
  \item With $R_f$, what is the Sharpe ratio of every frontier portfolio?
    \emph{(The same, maximum, Sharpe ratio:} $\sqrt{(\mub-R_f\one)\T V^{-1}(\mub-R_f\one)}$\emph{.)}
  \item In-sample SR 0.95, true max 0.49, true SR of estimated portfolio 0.34. Explain each.
  \item An asset has $\sigma_{np}=0$. What is its marginal contribution to portfolio volatility? \emph{(Zero.)}
  \item Difference between IR and appraisal ratio? \emph{(IR uses total active
    return/TE; appraisal uses alpha/residual vol, so it removes beta exposure.)}
  \item Write by hand the code for the tangency portfolio and its Sharpe ratio.
\begin{lstlisting}
ex = mu - rf * np.ones(len(mu))
x = np.linalg.solve(V, ex)
w_tan = x / x.sum()
sr_max = np.sqrt(ex @ x)                 # = sqrt(ex' V^{-1} ex)
\end{lstlisting}
  \item Write by hand the frontier weights for a target return \texttt{m} (no $R_f$).
\begin{lstlisting}
o = np.ones(len(mu)); Vi_o = np.linalg.solve(V, o); Vi_mu = np.linalg.solve(V, mu)
a1, a2, a3 = mu @ Vi_mu, mu @ Vi_o, o @ Vi_o; D = a1*a3 - a2**2
w = ((a1*Vi_o - a2*Vi_mu) + (a3*Vi_mu - a2*Vi_o)*m) / D
\end{lstlisting}
\end{enumerate}

\section{To do for next class}
Readings: Markowitz (1952, 1959), Merton (1972), Roll (1977, 1980);
Huang--Litzenberger (1988, ch.\,3); Moreira (2021, ch.\,7--8); Tidy Finance,
``Constrained Optimization and Backtesting''. Start Assignment \#2, Question 1.

\section*{References (Chapter 3)}
\begin{itemize}[leftmargin=*]\small
\item DeMiguel, V., and R.~Uppal. 2005. Portfolio investment with the exact tax basis via nonlinear programming. \emph{Management Science} 51(2): 277--290.
\item Huang, C.-F., and R.~Litzenberger. 1988. \emph{Foundations for Financial Economics}. Prentice Hall.
\item Markowitz, H.\,M. 1952. Portfolio selection. \emph{JF} 7(1): 77--91; 1959. \emph{Portfolio Selection}. Wiley.
\item Merton, R.\,C. 1972. An analytic derivation of the efficient portfolio frontier. \emph{JFQA} 7(4): 1851--1872.
\item Roll, R. 1977. A critique of the asset pricing theory's tests. \emph{JFE} 4(2): 129--176; 1980. Orthogonal portfolios. \emph{JFQA} 15(5): 1005--1023; 1992. A mean/variance analysis of tracking error. \emph{JPM} 18(4): 13--22.
\end{itemize}

% ===== end ch3 =====


% ===== begin ch4 =====
\chapter{Mean--Variance Portfolios that Adjust for Estimation Error}
\label{ch:4}

\begin{plainbox}
Chapter 3 gave the perfect recipe $\w=\frac1\gamma V^{-1}(\mub-R_f\one)$, assuming
we \emph{know} $\mub$ and $V$. Chapter 4 shows that with \emph{estimated} inputs
the recipe produces crazy, unstable portfolios that lose to the naive 1/N rule.
Then it studies the fixes. The single idea that connects all of them is
\textbf{shrinkage}: pull noisy estimates towards something simple and sensible.
\end{plainbox}

\section{What we want to do in Chapter 4}
\begin{enumerate}
  \item Show that the \textbf{out-of-sample} performance of sample-based MV portfolios is poor.
  \item Study methods to improve it: first those that fix the error in $\E[\R]$
    (Bayesian, short-sale constraints), then those that fix $V[\R]$ (GMV,
    Jagannathan--Ma, Ledoit--Wolf, norm constraints).
\end{enumerate}

% ---------------------------------------------------------------------
\section{The effect of estimation error on portfolio weights (focus)}
\[
\w=\underbrace{\frac1\gamma}_{\text{risk tolerance}}\ \underbrace{(V[\R])^{-1}}_{\text{ill-conditioned}}\ \underbrace{(\E[\R]-R_f\one)}_{\text{low precision}} .
\]

\subsection{Error in $\E[\R]$}
Merton (1980): sampling more often does not improve the mean (only the first and
last prices matter; Chapter 1). With annual volatility $\sigma_n$ and $T$ years:
\[
SE(\hat\mu_n)=\frac{\sigma_n}{\sqrt T}.
\]
\begin{keybox}[: how bad is it?]
Individual stock: $\sigma_n\approx30\%$. With 9 years, $SE=0.30/\sqrt9=10\%$, so the
95\% CI is $\pm20\%$ (a band 40\% wide). With expected returns around 10\%:
$0.10\pm0.20=[-10\%,\,30\%]$. \textbf{Useless in practice.}
\emph{Bottom line: beware of any model that uses expected returns as an input.}
\end{keybox}

\subsection*{Python exercise 1: means are hard to estimate, volatilities are not}
True $\mu=10\%$, $\sigma=30\%$; 5{,}000 samples of 5, 10, 40 years of monthly data.
\begin{lstlisting}
import numpy as np
import pandas as pd
rng = np.random.default_rng(41)
mu, sigma = 0.10, 0.30
S = 5000
sim = lambda T: rng.normal(mu/12, sigma/np.sqrt(12), size=(S, T))
est = lambda R: pd.DataFrame({"mean": 12 * R.mean(axis=1),
                              "volatility": np.sqrt(12) * R.std(axis=1, ddof=1)})
ranges = pd.concat({f"{T//12} years": est(sim(T)).quantile([0.025, 0.975])
                    for T in [60, 120, 480]})
print(ranges.round(3))
print("Years needed for +-2% (95%):", round((2 * sigma / 0.02)**2))
\end{lstlisting}
\begin{lstlisting}[style=out]
                 mean  volatility
5 years  0.025 -0.166       0.245
         0.975  0.360       0.354
10 years 0.025 -0.084       0.262
         0.975  0.285       0.337
40 years 0.025  0.006       0.281
         0.975  0.195       0.319
Years needed to estimate the mean within 2% (95% of the time): 900
\end{lstlisting}
\begin{plainbox}
With 10 years, the estimated mean can be anything between $-8\%$ and $+28\%$,
while the volatility is between 26\% and 34\%. To pin the mean down to $\pm2\%$
you need $(2\sigma/0.02)^2=(0.6/0.02)^2=900$ years of data.
\end{plainbox}

\subsection{Error in $V[\R]$}
$N(N+1)/2$ parameters: 1{,}275 for $N=50$; 5{,}050 for $N=100$.
\begin{itemize}
  \item \textbf{Singular} (not invertible) when $T\le N$: for $N=100$, 5 years of monthly data ($T=60$) are not enough.
  \item \textbf{Ill-conditioned} unless $T\gg N$: even $T=240$ gives a poorly conditioned matrix for $N=100$.
\end{itemize}
Result: sample-based weights are \textbf{unreasonable} (extreme), \textbf{unstable}
(fluctuate over time), \textbf{inefficient} (poor out of sample).

\subsection{A two-asset example}
Same setting as Chapter 3: $\gamma=5$, $R_f=1.01$, both assets gross mean 1.09, vol
20\%, correlation 0.99. True answer: MV weights $(0.20,0.20)$, $w_f=0.60$,
tangency $(0.5,0.5)$. Now suppose asset 1's mean is \textbf{estimated} as 1.10 instead of 1.09:
\begin{center}
\begin{tabular}{@{}lcc@{}}
\toprule
 & True parameters & Mean of asset 1 $=$ 1.10 \\
\midrule
MV weights & $(0.201,\ 0.201)$ & $(2.714,\ -2.286)$ \\
Risk-free weight & 0.598 & 0.573 \\
Tangency weights & $(0.5,\ 0.5)$ & $(6.353,\ -5.353)$ \\
\bottomrule
\end{tabular}
\end{center}
\begin{plainbox}
A 1\% error in one mean turns a 50/50 portfolio into $+635\%/-535\%$. Why? The
assets are 99\% correlated, so ``long asset 1, short asset 2'' is almost riskless
and earns the 1\% difference. The optimiser loves it, but the 1\% is probably
just estimation noise. With real data and more assets the weights are even more
extreme, and out-of-sample Sharpe ratios can be \emph{negative} (worse than
putting everything in the risk-free asset).
\end{plainbox}


\begin{workedex}[: why the weights explode, with a formula]
Two assets with the same volatility $\sigma$ and correlation $\rho$. Using the
$2\times2$ inverse,
\[
w_1-w_2=\frac{1}{\gamma}\,\frac{\mu_1-\mu_2}{\sigma^2(1-\rho)} .
\]
The difference between the two weights is the difference in means divided by
$\sigma^2(1-\rho)$, the variance of the long--short position. With $\gamma=5$,
$\sigma=0.2$, $\rho=0.99$ and $\mu_1-\mu_2=0.01$:
$w_1-w_2=0.01/(5\times0.04\times0.01)=5$, exactly $2.714-(-2.286)$. With $\rho=0.5$
the same 1\% error would give $w_1-w_2=0.1$ only. High correlation turns small
errors into huge positions.
\end{workedex}
\subsection*{Python exercise 2: small changes in means, large changes in weights}
Five identical assets (8\% excess mean, 20\% vol, correlation $\rho$); tangency $=1/N$.
Raise asset 1's mean by 1\% (about one-sixth of its SE with 10 years).
\begin{lstlisting}
import numpy as np
import pandas as pd
N, vol = 5, 0.20
mu = np.full(N, 0.08)
mu_new = mu + 0.01 * np.eye(N)[0]                    # asset 1: +1%
V = lambda rho: vol**2 * ((1 - rho)*np.eye(N) + rho)
tang = lambda m, rho: (lambda x: x / x.sum())(np.linalg.solve(V(rho), m))
res = pd.DataFrame({rho: tang(mu_new, rho) for rho in [0.0, 0.5, 0.8, 0.95]},
                   index=[f"asset {n}" for n in range(1, N+1)])
print(res.round(2))
\end{lstlisting}
\begin{lstlisting}[style=out]
correlation  0.00  0.50  0.80  0.95
asset 1      0.22  0.32  0.61  2.07
asset 2      0.20  0.17  0.10 -0.27
asset 3-5    0.20  0.17  0.10 -0.27
Standard error of the mean, 10 years of data: 0.063
\end{lstlisting}
\textbf{The more similar the assets, the more the optimiser bets on differences in
means that are pure estimation error.}

% ---------------------------------------------------------------------
\section{Performance of sample-based mean--variance portfolios}

\subsection{Methodology for out-of-sample evaluation (focus)}
\begin{keybox}[: the rolling-window backtest (used in almost every paper and in Assignment 2)]
\begin{enumerate}
  \item Choose an estimation window $T^{est}<T$ (monthly data: usually 60 or 120 months).
  \item Estimate the parameters on the window.
  \item Compute the optimal weights with the estimated parameters.
  \item Hold them for the \textbf{next} period and record the return (out of sample).
  \item Roll: drop the oldest observation, add the newest, repeat until the end.
  \item Compute performance metrics on the out-of-sample returns (SR, turnover, drawdowns, \dots).
\end{enumerate}
\end{keybox}
At month $t=60,\dots,T-1$, with the previous 60 months:
\[
\hat\mub_t=\frac1{60}\sum_{s=t-59}^{t}\R_s,\qquad
\hat V_t=\frac{1}{60-1}\sum_{s=t-59}^{t}(\R_s-\hat\mub_t)(\R_s-\hat\mub_t)\T,
\]
\begin{gather*}
\w_t^{MV}=\frac1\gamma\hat V_t^{-1}(\hat\mub_t-R_f\one),\qquad
R^{MV}_{p,t+1}=(1-\one\T\w_t^{MV})R_f+(\w_t^{MV})\T\R_{t+1},\\
\w_t^{tang}=\frac{\w_t^{MV}}{\one\T\w_t^{MV}},\qquad
R^{tang}_{p,t+1}=(\w_t^{tang})\T\R_{t+1},
\end{gather*}
\[
\w_t^{EW}=\tfrac1N\one,\qquad R^{EW}_{p,t+1}=(\w^{EW})\T\R_{t+1},\qquad
SR=\frac{12\,\E[R-R_f]}{\sqrt{12}\,\sigma[R-R_f]}\ \text{(annualised)} .
\]
Then test $H_0:SR^{MV}=SR^{EW}$ with Jobson--Korkie or a bootstrap (Chapter 2).

\subsection*{Python exercise 3: an out-of-sample backtest}
50 years of monthly excess returns on 10 assets (one-factor model, CAPM holds);
window $M=120$; three rules: MV ($\gamma=5$), GMV, 1/N.
\begin{lstlisting}
import numpy as np
import pandas as pd
rng = np.random.default_rng(45)
N, T, M = 10, 600, 120
gamma = 5
beta = np.linspace(0.6, 1.4, N)
mu = 0.005 * beta
V = 0.045**2*np.outer(beta, beta) + np.diag(np.full(N, 0.06**2))
R = rng.multivariate_normal(mu, V, size=T)

inv = lambda X, b: np.linalg.solve(np.cov(X, rowvar=False), b)
rules = {"mean-variance": lambda X: inv(X, X.mean(axis=0)) / gamma,
         "GMV": lambda X: (lambda x: x / x.sum())(inv(X, np.ones(N))),
         "1/N": lambda X: np.full(N, 1 / N)}
W = {k: np.array([f(R[t-M:t]) for t in range(M, T)]) for k, f in rules.items()}
oos = {k: (w * R[M:]).sum(axis=1) for k, w in W.items()}       # out-of-sample
sr = lambda r: np.sqrt(12) * r.mean() / r.std(ddof=1)
tsr = lambda w: (lambda m, v: np.sqrt(12) * m.mean() / np.sqrt(v.mean() + m.var()))(
    w @ mu, np.einsum("tn,nk,tk->t", w, V, w))                 # true SR
to = lambda w: np.abs(np.diff(w, axis=0)).sum(axis=1).mean()   # turnover
print(pd.DataFrame({"realized SR": {k: sr(r) for k, r in oos.items()},
                    "true SR": {k: tsr(w) for k, w in W.items()},
                    "turnover": {k: to(w) for k, w in W.items()}}).round(2))
\end{lstlisting}
\begin{lstlisting}[style=out]
True maximum Sharpe ratio: 0.36
               realized SR  true SR  turnover
mean-variance         0.09     0.12      0.34
GMV                   0.24     0.28      0.06
1/N                   0.30     0.35      0.00
\end{lstlisting}
MV gets a third of the maximum SR and trades the most. GMV (no means) is much
better. 1/N is best here because it is close to the true tangency portfolio.
Each realised SR has an SE of about $1/\sqrt{40}\approx0.16$ (40 years out of sample).
\begin{addnote}
\textbf{Reading the code.} \texttt{R[t-M:t]} contains months $t-M,\dots,t-1$, so
the weights for month $t$ use only the past (no look-ahead). \texttt{(w * R[M:]).sum(axis=1)}
multiplies each row of weights by the \emph{same row} of returns: row $j$ of
\texttt{W} was estimated with data up to $M+j-1$ and is applied to the return of
month $M+j$. \texttt{np.einsum("tn,nk,tk->t", w, V, w)} computes $\w_t\T V\w_t$
for every $t$ at once.
\end{addnote}

\subsection{Evidence: DeMiguel, Garlappi and Uppal (2009)}
Six datasets (Ken French, MSCI, S\&P sectors), $T^{est}=120$ months.
\begin{center}\small
\resizebox{\textwidth}{!}{%
\begin{tabular}{@{}lcccccc@{}}
\toprule
 & S\&P sectors & Industry & Int'l & Mkt+SMB/HML & FF 1-factor & FF 4-factor\\
 & $N=11$ & $N=11$ & $N=9$ & $N=3$ & $N=21$ & $N=24$ \\
\midrule
MV (in sample) & 0.3848 & 0.2124 & 0.2090 & 0.2851 & 0.5098 & 0.5364 \\
MV (out of sample) & 0.0794 & $-0.0363$ & $-0.0719$ & 0.2186 & $-0.0684$ & $-0.0031$ \\
1/N & 0.1876 & 0.1353 & 0.1277 & 0.2240 & 0.1623 & 0.1753 \\
$p$-value & (0.12) & (0.01) & (0.01) & (0.46) & (0.00) & (0.01) \\
\midrule
MV turnover / 1/N turnover & 38.99 & 607{,}479.61 & 4{,}236.10 & 2.83 & 9{,}408.16 & 2{,}672.71 \\
\bottomrule
\end{tabular}}
\end{center}
\begin{keybox}[: conclusions]
\begin{enumerate}
  \item Out-of-sample MV is far worse than in-sample; \textbf{negative SR in 4 of 6 datasets}.
  \item \textbf{1/N beats MV in all six}, significantly at 5\% in four.
  \item MV turnover is enormous (up to 600{,}000 times that of 1/N).
\end{enumerate}
\end{keybox}


\begin{addnote}
\textbf{How long would we need to wait for MV to beat 1/N?} DGU (2009) compute that
the estimation window needed for the sample-based MV strategy (and its extensions)
to outperform 1/N is around \textbf{3{,}000 months (250 years) for 25 assets} and
about \textbf{6{,}000 months for 50 assets}. With realistic data, 1/N is simply
very hard to beat.
\end{addnote}
% ---------------------------------------------------------------------
\section{Models that adjust for estimation error: the menu}
The problem has been known since the 1970s.
\begin{center}
\begin{tabularx}{\textwidth}{@{}lX@{}}
\toprule
Bayesian & Diffuse priors (Barry 1974; Bawa--Brown--Klein 1979); shrinkage of means (Jobson--Korkie--Ratti 1979; Jobson--Korkie 1980; Jorion 1985, 1986); covariance estimation (Best--Grauer 1992; Chan--Karceski--Lakonishok 1999; Ledoit--Wolf 2004); Black--Litterman (1990, 1992; Chapter 5). \\
Non-Bayesian & Short-sale constraints (Frost--Savarino 1988; Chopra 1993; Jagannathan--Ma 2003); robust optimisation (Goldfarb--Iyengar 2003; Uppal--Wang 2003; Garlappi--Uppal--Wang 2007; Boyle et al.\ 2012); optimal diversification across market and estimation risk (Kan--Zhou 2005); resampling (Michaud 1998). \\
\bottomrule
\end{tabularx}
\end{center}
\begin{keybox}[: the central idea of the whole course]
\textbf{Shrinkage.} Almost every method can be read as shrinking $\hat\mub$ and
$\hat V$ towards a sensible target:
\begin{itemize}
  \item means $\to$ the average of all means, or the market's expected return;
  \item covariance $\to$ a one-factor (market-model) matrix, or a matrix with all
    correlations equal, or a multiple of the identity.
\end{itemize}
\end{keybox}
\begin{plainbox}
Think of a new restaurant with three 5-star reviews. You do not believe it is the
best restaurant in town: you ``shrink'' the score towards the average restaurant,
because three reviews are too few. The less data you have, the more you shrink.
Exactly the same logic applies to sample means and covariances.
\end{plainbox}

% ---------------------------------------------------------------------
\section{Bayesian approaches}
Bayesian estimates of $\mub$ and $V$ come from the \textbf{predictive (posterior)
distribution}, obtained by integrating the likelihood $f(\R\mid\mub,V)$ over the
parameters with respect to a prior $p(\mub,V)$.

\paragraph{1. Diffuse prior} (Barry 1974; Klein--Bawa 1976; Brown 1979). With
$p(\mub,V)\propto|V|^{-(N+1)/2}$ and Normal likelihood, the predictive
distribution is Student-$t$ with mean $\hat\mub$ and variance $\hat V(1+1/T^{est})$.
With $T^{est}=60$ the correction is $1/60<2\%$: \textbf{practically the same as
sample MV}.

\paragraph{2. Bayes--Stein shrinkage} (Stein 1956; James--Stein 1961; Jorion 1985, 1986).
\begin{keybox}[: Bayes--Stein estimator of the means]
\begin{gather*}
\hat\mub^{BS}_t=(1-\hat\psi_t)\,\hat\mub_t+\hat\psi_t\,\bar\mu_t\one,\\
\hat\psi_t=\frac{N+2}{(N+2)+T^{est}\,(\hat\mub_t-\mu^{GMV}_t\one)\T\hat V_t^{-1}(\hat\mub_t-\mu^{GMV}_t\one)},
\end{gather*}
with $0<\hat\psi_t<1$, $\hat V_t=\frac{1}{T^{est}-N-2}\sum(\R_s-\hat\mub_t)(\R_s-\hat\mub_t)\T$,
and grand mean $\bar\mu_t=\mu^{GMV}_t$ (Jorion: the mean return of the GMV portfolio).
\end{keybox}
\begin{intuition}
$\hat\psi$ is large (more shrinkage) when: $T^{est}$ is small (little data); the
sample means are close to each other relative to their noise (the quadratic form
is small, i.e.\ the differences between means are not statistically convincing);
$N$ is large (many means to estimate). In practice the error is so large that
the optimal weight on the grand mean is often very high (the slides mention
around 95\% in applications).
\end{intuition}


\begin{workedex}[: size of the Bayes--Stein shrinkage]
$N=10$ assets, $T^{est}=60$ months, so $N+2=12$ and
$\hat\psi=12/(12+60\,q)$ where $q$ is the quadratic form (how far the sample means
are from the grand mean, measured in units of their risk):
\begin{center}
\begin{tabular}{@{}lccc@{}}
\toprule
$q$ (dispersion of means) & 0.05 & 0.2 & 1.0 \\
\midrule
$\hat\psi$ (weight on grand mean) & 0.80 & 0.50 & 0.17 \\
\bottomrule
\end{tabular}
\end{center}
When the sample means are close together (small $q$) we cannot trust their
differences, so we shrink a lot. Only very dispersed means survive.
\end{workedex}
\subsection*{Python exercise 4: Bayes--Stein shrinkage of the means}
2{,}000 samples, 5 years, 10 assets.
\begin{lstlisting}
import numpy as np
import pandas as pd
rng = np.random.default_rng(44)
N, T, S = 10, 60, 2000
beta = np.linspace(0.6, 1.4, N)
mu = 0.005 * beta
V = 0.045**2*np.outer(beta, beta) + np.diag(np.full(N, 0.06**2))
R = rng.multivariate_normal(mu, V, size=(S, T))        # S samples of T x N
m = R.mean(axis=1)                                     # S x N sample means
D = R - m[:, None, :]
Vh = np.einsum("stn,stk->snk", D, D) / (T - N - 2)     # S covariance matrices
Vi = np.linalg.inv(Vh)
wg = Vi.sum(axis=2) / Vi.sum(axis=(1, 2))[:, None]     # GMV weights
m0 = (wg * m).sum(axis=1)                              # GMV mean = grand mean
d = m - m0[:, None]
q = np.einsum("sn,snk,sk->s", d, Vi, d)
psi = (N + 2) / ((N + 2) + T * q)                      # shrinkage factor
m_bs = (1 - psi)[:, None] * m + (psi * m0)[:, None]

mse = lambda e: (144 * ((e - mu)**2).sum(axis=1)).mean()
tsr = lambda e: (lambda w: np.sqrt(12) * (w @ mu) / np.sqrt(np.einsum("sn,nk,sk->s", w, V, w)))(
    np.einsum("snk,sk->sn", Vi, e))
est = {"sample": m, "Bayes-Stein": m_bs}
print(pd.DataFrame({"MSE": {k: mse(e) for k, e in est.items()},
                    "median true SR": {k: np.median(tsr(e)) for k, e in est.items()}}).round(3))
\end{lstlisting}
\begin{lstlisting}[style=out]
Average shrinkage factor psi: 0.6     True maximum Sharpe ratio: 0.36
              mean squared error   median true SR
sample                     0.138            0.085
Bayes-Stein                0.068            0.148
\end{lstlisting}
On average 60\% weight on the grand mean; shrinkage \textbf{halves the MSE} of the
means; the true SR rises from 0.09 to 0.15, still far below 0.36.
\textbf{Shrinking the means helps but does not solve the problem.}

\paragraph{3. ``Data-and-Model'' (DM)} (P\'astor 2000; P\'astor--Stambaugh 2000).
The shrinkage target comes from an \textbf{asset-pricing model} (CAPM, FF3, FF4),
and the amount of shrinkage from how strongly the investor believes it. The prior
on mispricing $\alpha$ is Normal around 0 with $\sigma_\alpha=1\%$ p.a.: with 95\%
probability mispricing is between $-2\%$ and $+2\%$. This removes the arbitrary
choice of target and of $\psi$.

\paragraph{Performance (DGU 2009).}
\begin{center}\small
\resizebox{\textwidth}{!}{%
\begin{tabular}{@{}lcccccc@{}}
\toprule
 & S\&P & Ind. & Int'l & Mkt+SMB/HML & FF1 & FF4\\
\midrule
1/N & 0.1876 & 0.1353 & 0.1277 & 0.2240 & 0.1623 & 0.1753 \\
Bayes--Stein & 0.0811 & $-0.0319$ & $-0.0528$ & \textbf{0.2536} & $-0.0636$ & $-0.0042$ \\
Data-Model ($\sigma_\alpha=1\%$) & 0.1410 & $-0.0475$ & 0.1032 & 0.0181 & 0.0745 & \textbf{0.2355} \\
\midrule
Turnover / 1/N: Bayes--Stein & 22.41 & 10{,}092 & 1{,}761 & 1.85 & 10{,}510 & 2{,}784 \\
Turnover / 1/N: Data-Model & 1.72 & 20{,}468 & 45.24 & 76.85 & 762.83 & 32.46 \\
\bottomrule
\end{tabular}}
\end{center}
Bayes--Stein: negative SR in 4 of 6. Each Bayesian strategy beats 1/N in only one
dataset. Turnover lower than MV but much higher than 1/N.
\textbf{Conclusion: out of sample, 1/N usually beats the Bayesian strategies.}

% ---------------------------------------------------------------------
\section{Short-sale-constrained portfolios (focus)}
In the two-asset example the problem was $+271\%/-228\%$. Requiring $w_n\ge0$
limits how extreme weights can be. It is ad hoc but popular (and often required:
pension funds by regulation; mutual funds find shorting costly). It can be shown
formally to be \textbf{shrinkage}.

\paragraph{Strategies.} MV-C (sample MV with $w\ge0$), BS-C (Bayes--Stein with
$w\ge0$), GMV-C (next section).

\subsection{Why short-sale constraints shrink the means}
\[
\max_{\w}\ R_f+\w\T(\mub-R_f\one)-\frac\gamma2\w\T V\w\qquad\text{s.t. } w_n\ge0\ \forall n.
\]
Lagrangian with multipliers $\bm\lambda_{ss}=(\lambda_1,\dots,\lambda_N)\T\ge\bm0$:
\begin{align*}
L&=R_f+\w\T(\mub-R_f\one)-\frac\gamma2\w\T V\w+\w\T\bm\lambda_{ss}\\
&=R_f+\w\T\big(\underbrace{\mub+\bm\lambda_{ss}}_{\text{adjusted mean}}-R_f\one\big)-\frac\gamma2\w\T V\w .
\end{align*}
\begin{keybox}[: interpretation]
The constrained solution equals the \emph{unconstrained} solution with means
$\mub+\bm\lambda_{ss}$. The constraint binds ($\lambda_n>0$) for assets with
\textbf{low} $\mu_n$ (those the optimiser wanted to short), and their mean is
\textbf{raised} to $\mu_n+\lambda_n$. Low means are pushed up towards the others:
\textbf{shrinkage of expected returns towards the average.}
\end{keybox}
\begin{addnote}
\textbf{Complementary slackness (KKT).} For each asset either $w_n>0$ and
$\lambda_n=0$ (constraint not binding, mean untouched) or $w_n=0$ and
$\lambda_n\ge0$ (just enough to make the optimiser stop wanting to go short).
\end{addnote}

\paragraph{Performance (DGU 2009).}
\begin{center}\small
\resizebox{\textwidth}{!}{%
\begin{tabular}{@{}lcccccc@{}}
\toprule
 & S\&P & Ind. & Int'l & Mkt+SMB/HML & FF1 & FF4\\
\midrule
MV (out of sample) & 0.0794 & $-0.0363$ & $-0.0719$ & 0.2186 & $-0.0684$ & $-0.0031$ \\
1/N & 0.1876 & 0.1353 & 0.1277 & 0.2240 & 0.1623 & 0.1753 \\
MV-C & 0.0892 & 0.0678 & 0.0848 & 0.1084 & 0.1977 & 0.2024 \\
BS-C & 0.1075 & 0.0819 & 0.0848 & 0.1514 & 0.1955 & 0.2062 \\
\midrule
Turnover / 1/N: MV-C & 4.53 & 7.17 & 7.23 & 4.12 & 17.53 & 13.82 \\
Turnover / 1/N: BS-C & 3.64 & 7.22 & 6.10 & 3.65 & 17.32 & 13.07 \\
\bottomrule
\end{tabular}}
\end{center}
\begin{keybox}[: conclusions on short-sale constraints]
\begin{itemize}
  \item Constrained beats unconstrained in 5 of 6 datasets (not with $N=3$).
  \item \textbf{Very effective at reducing turnover} (from thousands of times 1/N to 4--18 times).
  \item Less successful at improving SR: 1/N still wins in 4 of 6 (constraints win only in FF1 and FF4).
\end{itemize}
\end{keybox}

% ---------------------------------------------------------------------
\section{Global-minimum-variance portfolios (focus)}
\begin{plainbox}
Even after fixing them, the estimated means are too noisy. So \textbf{drop them
completely}: assume all assets have the same expected return. Then the
target-return constraint disappears and we only minimise risk. Risk (covariances)
can be estimated much more precisely than means. GMV $=$ shrinking all expected
returns 100\% towards the same value.
\end{plainbox}
Evidence that even adjusted means perform poorly: DGU (2009), Jacobs--M\"uller--Weber (2014).

\paragraph{Five variants.}
\begin{enumerate}
  \item \textbf{GMV} with the sample covariance (Jorion 1985, 1986).
  \item \textbf{Factor-structure covariance} (Green--Hollifield 1992; Chan--Karceski--Lakonishok 1999):
    $V$ built from a $K\times K$ factor covariance, $N\times K$ loadings, $N$ residual variances.
    Performs worse than 3--5 (Jagannathan--Ma 2003); factor models later.
  \item \textbf{GMV with short-sale constraints} (Jagannathan--Ma 2003): indirect shrinkage of $V$.
  \item \textbf{Ledoit--Wolf shrinkage} (2003, 2004a,b, 2017, 2022): explicit shrinkage of $V$.
  \item \textbf{Norm constraints} (DeMiguel, Garlappi, Nogales, Uppal 2009): generalised shrinkage.
\end{enumerate}

\subsection{Variant 1: the GMV portfolio}
\begin{keybox}[: GMV weights]
\[
\min_{\w}\ \w\T V\w\quad\text{s.t.}\quad\w\T\one=1
\qquad\Rightarrow\qquad
\w_{GMV}=\frac{V^{-1}\one}{\one\T V^{-1}\one}.
\]
Numerator: ``raw'' weights, inversely related to risk. Denominator: their sum, so that they add up to 1.
\end{keybox}
\begin{addnote}
\textbf{Derivation.} $L=\w\T V\w+\lambda(1-\w\T\one)$; FOC $2V\w=\lambda\one\Rightarrow\w=\frac\lambda2V^{-1}\one$;
the constraint gives $\frac\lambda2=1/(\one\T V^{-1}\one)$.
\end{addnote}

\paragraph{Example.} Two assets with means 1.10 and 1.09 (the ``estimation error''
case), vol 20\%, correlation 0.99: $\w_{GMV}=(0.5,0.5)$, exactly the true
answer. \textbf{GMV is immune to errors in expected returns.}

\paragraph{Zero correlations.} If $V$ is diagonal, $V^{-1}=\text{diag}(1/\sigma_1^2,\dots,1/\sigma_N^2)$ and
\[
w_n=\frac{1/\sigma_n^2}{\sum_{m=1}^N1/\sigma_m^2}\qquad\text{(inverse-variance weights).}
\]
Example: $V=\text{diag}(0.0625,0.1225,0.36)$ $\Rightarrow$ $\w=(0.594,\,0.303,\,0.103)$.
In code:
\begin{lstlisting}
w = np.diag(Vinv) / np.sum(np.diag(Vinv))    # only valid when V is diagonal
\end{lstlisting}

\paragraph{When does GMV earn a high Sharpe ratio?} It overweights low-volatility
(and, with correlations, low-beta) assets (Scherer 2011; Clarke--de Silva--Thorley 2010).
So GMV does well when \textbf{low-volatility / low-beta stocks earn high returns}
(the ``low-volatility anomaly''; Dangl--Kashofer 2013). Results are very sensitive
to how returns and volatility are measured (Plyakha, Uppal, Vilkov 2021).

\paragraph{When are GMV weights reasonable?} When $V$ is estimated precisely.
Recent realised volatility (e.g.\ last 30 days) forecasts future volatility well
(Ghysels--Santa-Clara--Valkanov 2006; Andersen--Benzoni 2010); option-implied
volatilities and correlations help too (Driessen--Maenhout--Vilkov 2009;
DeMiguel--Plyakha--Uppal--Vilkov 2013). Remaining problem: $\hat V$ is
\textbf{ill-conditioned} $\Rightarrow$ variants 3, 4, 5. DGNU (2009) show all
three are forms of shrinkage/regularisation. Practitioners mostly use
Ledoit--Wolf; Amundi sells an ETF based on DGNU (2009).

\subsection{Variant 3: GMV with short-sale constraints (Jagannathan--Ma 2003)}
\[
\min_{\w}\ \w\T\hat V\w\qquad\text{s.t.}\quad\w\T\one=1,\quad\w\ge\bm0 .
\]
\begin{keybox}[: Jagannathan--Ma shrinkage interpretation]
The solution coincides with the \emph{unconstrained} GMV computed with
\[
\hat V_{JM}=\hat V-\bm\lambda_{ss}\one\T-\one\bm\lambda_{ss}\T,
\]
where $\lambda_{ss,n}=0$ if $w_n>0$ and $\lambda_{ss,n}\ge0$ if the constraint binds ($w_n=0$).
If asset $n$'s constraint binds, its covariance with every other asset is reduced by
$\lambda_n$ (and its variance by $2\lambda_n$): it looks \textbf{less risky}, so its weight
rises from negative to exactly zero.
\end{keybox}
\paragraph{Example ($N=3$, only asset 3 wants to go short).} $\lambda_1=\lambda_2=0$, $\lambda_3>0$:
\[
\hat V_{JM}=\begin{pmatrix}\sigma_{11}&\sigma_{12}&\sigma_{13}-\lambda_3\\
\sigma_{12}&\sigma_{22}&\sigma_{23}-\lambda_3\\
\sigma_{13}-\lambda_3&\sigma_{23}-\lambda_3&\sigma_{33}-2\lambda_3\end{pmatrix}.
\]
\begin{intuition}
Why does it help? Assets that the unconstrained GMV shorts are typically those
with high estimated covariances with the others, and high sample covariances are
partly \emph{over}-estimated noise. Cutting them is exactly what a good estimator
would do. Jagannathan and Ma's title says it: ``why imposing the wrong
constraints helps''. \textbf{Drawback:} many weights end up exactly zero (poor
diversification); in the slides' picture, the constrained optimum sits on the
edge of the triangle of feasible weights.
\end{intuition}

\subsection*{Python exercise 5: short-sale constraints as shrinkage}
200 samples of 5 years, 25 assets.
\begin{lstlisting}
import numpy as np
import pandas as pd
from scipy.optimize import minimize
rng = np.random.default_rng(46)
N, T, S = 25, 60, 200
beta = rng.uniform(0.5, 1.5, N)
V = 0.045**2*np.outer(beta, beta) + np.diag(rng.uniform(0.04, 0.10, N)**2)
one = np.ones(N)
gmv = lambda C: (lambda x: x / x.sum())(np.linalg.solve(C, one))
budget = {"type": "eq", "fun": lambda w: w.sum() - 1}
gmv_c = lambda C: minimize(lambda w: w @ C @ w, one / N, method="SLSQP",
                           bounds=[(0, 1)] * N, constraints=[budget],
                           options={"ftol": 1e-12}).x
Vh = [np.cov(rng.multivariate_normal(np.zeros(N), V, size=T), rowvar=False) for s in range(S)]
W = {"GMV": np.array([gmv(C) for C in Vh]),
     "GMV, no short sales": np.array([gmv_c(C) for C in Vh])}
vol = lambda w: np.sqrt(12 * np.einsum("sn,nk,sk->s", w, V, w))   # TRUE volatility
print(pd.DataFrame({"true vol": {k: vol(w).mean() for k, w in W.items()},
                    "min weight": {k: w.min(axis=1).mean() for k, w in W.items()},
                    "max weight": {k: w.max(axis=1).mean() for k, w in W.items()},
                    "zeros": {k: (np.abs(w) < 1e-6).sum(axis=1).mean() for k, w in W.items()}}).round(3))
\end{lstlisting}
\begin{lstlisting}[style=out]
True GMV volatility: 0.113      1/N volatility: 0.161
                     true vol  min weight  max weight  zeros
GMV                     0.148      -0.169       0.295   0.00
GMV, no short sales     0.131       0.000       0.243  14.22
\end{lstlisting}
Estimated GMV: 14.8\% true vol (true minimum 11.3\%) because it shorts to exploit
estimation error. No short sales: 13.1\%, with about 14 of 25 weights at zero.

\paragraph{Performance (DGU 2009).}
\begin{center}\small
\resizebox{\textwidth}{!}{%
\begin{tabular}{@{}lcccccc@{}}
\toprule
 & S\&P & Ind. & Int'l & Mkt+SMB/HML & FF1 & FF4\\
\midrule
MV (out of sample) & 0.0794 & $-0.0363$ & $-0.0719$ & 0.2186 & $-0.0684$ & $-0.0031$ \\
MV-C & 0.0892 & 0.0678 & 0.0848 & 0.1084 & 0.1977 & 0.2024 \\
1/N & 0.1876 & 0.1353 & 0.1277 & 0.2240 & 0.1623 & 0.1753 \\
GMV & 0.0820 & 0.1554 & 0.1490 & 0.2493 & 0.2778 & $-0.0183$ \\
GMV-C & 0.0834 & 0.1425 & 0.1501 & 0.2493 & 0.1546 & 0.3580 \\
\midrule
Turnover / 1/N: GMV & 6.54 & 21.65 & 7.30 & 1.11 & 45.47 & 6.83 \\
Turnover / 1/N: GMV-C & 2.47 & 2.58 & 2.27 & 1.11 & 3.93 & 1.76 \\
\bottomrule
\end{tabular}}
\end{center}
\begin{keybox}[: conclusions on GMV]
\begin{itemize}
  \item \textbf{Ignoring expected returns pays}: GMV almost always beats MV; GMV-C beats MV-C in 4 of 6.
  \item Effect of the short-sale constraint on GMV's SR is mixed (small in 1--4,
    negative with $N=21$, strongly positive with $N=24$).
  \item But it \textbf{cuts turnover}: GMV-C is only 1.1--3.9 times 1/N.
\end{itemize}
\end{keybox}

\subsection{Variant 4: Ledoit--Wolf shrinkage of the covariance matrix}
\begin{keybox}[: Ledoit--Wolf estimator]
\[
\hat V_{LW}=(1-\psi)\,\hat V+\psi\,\hat V_{target},\qquad 0\le\psi\le1 .
\]
$\psi$ (often called $\alpha$ in papers and code) is chosen to minimise the expected
squared \textbf{Frobenius norm} $\|\hat V_{LW}-V\|_F^2$, where
$\|A\|_F=\sqrt{\sum_n\sum_m a_{nm}^2}$.
\end{keybox}
\paragraph{Shrinkage targets.}
\begin{enumerate}
  \item \textbf{LW 2003}: the one-factor (market-model) covariance matrix.
  \item \textbf{LW 2004b} (``Honey, I shrunk the sample covariance matrix''): constant correlation
    (all pairwise correlations $=$ average correlation; sample variances unchanged).
  \item \textbf{LW 2004a}: $\hat V_{target}=\frac{\text{trace}(\hat V)}{N}I_N=\bar\lambda I_N$,
    the average variance on the diagonal ($\bar\lambda$ = average eigenvalue, since the trace is the sum of eigenvalues).
\end{enumerate}
\begin{keybox}[: why identity-target shrinkage fixes the condition number]
Shrinking towards $\bar\lambda I_N$ changes every eigenvalue
$\lambda_n\mapsto(1-\psi)\lambda_n+\psi\bar\lambda$ and leaves the eigenvectors unchanged.
Large eigenvalues fall, small ones rise $\Rightarrow$ \textbf{lower condition number};
the average eigenvalue (total variance) does not change. Equivalent to an
$\ell_2$-penalised maximum-likelihood estimator.
\end{keybox}
\paragraph{$N=3$ example (LW 2004a).}
\[
\hat V_{LW}=(1-\psi)\begin{pmatrix}\hat\sigma_{11}&\hat\sigma_{12}&\hat\sigma_{13}\\\hat\sigma_{12}&\hat\sigma_{22}&\hat\sigma_{23}\\\hat\sigma_{13}&\hat\sigma_{23}&\hat\sigma_{33}\end{pmatrix}
+\psi\,\frac{\hat\sigma_{11}+\hat\sigma_{22}+\hat\sigma_{33}}{3}\,I_3 .
\]
Off-diagonal covariances are multiplied by $(1-\psi)$; the diagonal is pulled towards the average variance.
\begin{plainbox}
With $\psi=0$ you get the sample GMV; with $\psi=1$ (and target $\bar\lambda I$) all
assets look identical and uncorrelated, so GMV becomes \textbf{1/N}. Ledoit--Wolf
picks the point in between that is, on average, closest to the truth. So GMV-LW
lies between the sample GMV and 1/N.
\end{plainbox}
\begin{lstlisting}
from sklearn.covariance import LedoitWolf
lw = LedoitWolf().fit(X)          # X: T x N array of returns
V_lw = lw.covariance_             # shrunk covariance (identity-type target)
psi = lw.shrinkage_               # estimated shrinkage intensity
w = np.linalg.solve(V_lw, np.ones(X.shape[1])); w /= w.sum()
\end{lstlisting}
Other code: OAS (Oracle Approximating Shrinkage, Chen--Wiesel--Eldar--Hero 2010) in scikit-learn.


\begin{workedex}[: Ledoit--Wolf with two assets]
Sample covariance $\hat V=\begin{pmatrix}0.04&0.018\\0.018&0.01\end{pmatrix}$
(vols 20\% and 10\%, correlation 0.9). Target: average variance
$(0.04+0.01)/2=0.025$ times $I$. With $\psi=0.5$:
\[
\hat V_{LW}=0.5\hat V+0.5\begin{pmatrix}0.025&0\\0&0.025\end{pmatrix}
=\begin{pmatrix}0.0325&0.009\\0.009&0.0175\end{pmatrix}.
\]
\begin{center}
\begin{tabular}{@{}lccc@{}}
\toprule
 & implied correlation & condition number & GMV weights \\
\midrule
Sample $\hat V$ & 0.90 & 30.9 & $(-0.57,\ 1.57)$ \\
Ledoit--Wolf & 0.38 & 2.8 & $(0.27,\ 0.73)$ \\
\bottomrule
\end{tabular}
\end{center}
Shrinkage lowers the correlation, makes the matrix well conditioned, and removes
the short position. In practice $\psi$ is not chosen by hand: Ledoit--Wolf estimate
it from the data (\texttt{LedoitWolf().fit(X).shrinkage\_}).
\end{workedex}
\subsection*{Python exercise 6: Ledoit--Wolf shrinkage with 50 assets}
\begin{lstlisting}
import numpy as np
import pandas as pd
from sklearn.covariance import LedoitWolf
rng = np.random.default_rng(47)
N, S = 50, 200
beta = rng.uniform(0.5, 1.5, N)
V = 0.045**2*np.outer(beta, beta) + np.diag(rng.uniform(0.04, 0.10, N)**2)
one = np.ones(N)
gmv = lambda C: (lambda x: x / x.sum())(np.linalg.solve(C, one))
vol = lambda w: np.sqrt(12 * w @ V @ w)
one_run = lambda X: (lambda C, L: [vol(gmv(C)), vol(gmv(L)), np.linalg.cond(C), np.linalg.cond(L)])(
    np.cov(X, rowvar=False), LedoitWolf().fit(X).covariance_)
cols = ["vol, sample", "vol, LW", "cond, sample", "cond, LW"]
res = pd.DataFrame({f"T = {T}": np.mean([one_run(rng.multivariate_normal(np.zeros(N), V, size=T))
                                         for s in range(S)], axis=0) for T in [60, 120, 240]},
                   index=cols).T
print(res)
\end{lstlisting}
\begin{lstlisting}[style=out]
True GMV volatility: 0.093   1/N volatility: 0.169   Condition number of true V: 74
          vol, sample  vol, LW  cond, sample  cond, LW
T = 60          0.235    0.121          3887       101
T = 120         0.122    0.111           263       112
T = 240         0.105    0.103           132        98
\end{lstlisting}
With $T=60\approx N$: sample condition number $\approx3{,}900$ and GMV true vol 23.5\%
(\emph{worse than 1/N}). LW: condition number $\approx100$, vol 12.1\% (close to 9.3\%).
The gain shrinks as $T$ grows.

\paragraph{Nonlinear shrinkage} (LW 2017, 2020, 2022): shrinks each eigenvalue
differently. Gains over linear shrinkage are modest, so not in the course
(possible master's project).

\subsection{Variant 5: norm-constrained portfolios (DGNU 2009)}
\[
\min_{\w}\ \w\T\hat V\w\qquad\text{s.t.}\quad\w\T\one=1,\quad\|\w\|\le\delta .
\]
\begin{center}
\begin{tabularx}{\textwidth}{@{}lXX@{}}
\toprule
Norm & Definition & Effect \\
\midrule
1-norm & $\|\w\|_1=\sum_n|w_n|\le\delta$ & Limits total gross exposure (longs $+$ shorts). Feasible set is a diamond/polygon $\Rightarrow$ corner solutions, \textbf{some weights exactly zero}. With $\delta=1$ it is the short-sale constraint. \\
2-norm & $\|\w\|_2=\big(\sum_n w_n^2\big)^{1/2}\le\delta$ & Feasible set is a circle/ball $\Rightarrow$ \textbf{rarely exact zeros}, well diversified. Smallest possible 2-norm with $\sum w=1$ is 1/N. \\
A-norm & $\|\w\|_A=(\w\T A\w)^{1/2}\le\delta$ & General version. \\
\bottomrule
\end{tabularx}
\end{center}
\begin{addnote}
\textbf{Why the 2-norm constraint is shrinkage.} Adding the constraint
$\w\T\w\le\delta^2$ with multiplier $\nu\ge0$ gives the Lagrangian
$\w\T\hat V\w+\nu(\w\T\w-\delta^2)+\dots=\w\T(\hat V+\nu I)\w+\dots$, so the
solution is the unconstrained GMV with $\hat V+\nu I$: the sample covariance plus a
multiple of the identity, which is the same idea as Ledoit--Wolf's identity target.
(Same family as ridge regression; the 1-norm is the LASSO family.)
\end{addnote}

\paragraph{Choosing $\delta$: cross-validation.} Unlike LW (distributional
assumptions), let the data choose $\delta$, \emph{using only historical data}
(no look-ahead):
\begin{enumerate}
  \item split the estimation sample into 5 folds; define a grid of $\delta$ values;
  \item for each $\delta$, estimate on 4 folds and evaluate on the held-out fold
    (e.g.\ minimum volatility or maximum SR);
  \item rotate the folds; pick the $\delta$ with the best average hold-out performance.
\end{enumerate}
Code in scikit-learn; you do not need to learn cross-validation for this course.

\subsection{Performance of Ledoit--Wolf and norm-constrained portfolios}
\begin{center}\small
\resizebox{\textwidth}{!}{%
\begin{tabular}{@{}lccccc@{}}
\toprule
Sharpe ratio & 10Ind & 48Ind & 6FF & 25FF & 500CRSP\\
\midrule
1/N & 0.2541 & 0.2508 & 0.2563 & 0.2565 & 0.3326 \\
GMV-C & 0.2852 & 0.2914 & 0.2629 & 0.2720 & 0.3985 \\
LW-diagonal & 0.2962 & 0.2620 & 0.3226 & 0.3974 & 0.4086 \\
LW-1-factor & 0.2902 & 0.2544 & 0.3296 & 0.3927 & 0.4500 \\
NC-1-norm & 0.2890 & 0.2831 & 0.3374 & 0.3553 & 0.3706 \\
NC-2-norm & \textbf{0.3193} & 0.2891 & \textbf{0.3922} & \textbf{0.4278} & \textbf{0.4672} \\
\midrule
Turnover & & & & & \\
1/N & 0.0232 & 0.0311 & 0.0155 & 0.0174 & 0.0595 \\
GMV & 0.1656 & 0.8286 & 0.2223 & 0.7953 & 0.7769 \\
GMV-C & 0.0552 & 0.0741 & 0.0461 & 0.0841 & 0.4222 \\
LW-diagonal & 0.1132 & 0.4029 & 0.0905 & 0.3144 & 0.5594 \\
LW-1-factor & 0.1428 & 0.4034 & 0.1455 & 0.4893 & 0.5463 \\
NC-1-norm & 0.6013 & 0.8232 & 1.0064 & 0.9767 & 0.9753 \\
NC-2-norm & 1.0177 & 2.7556 & 1.6594 & 3.6275 & 1.0443 \\
\bottomrule
\end{tabular}}
\end{center}
\textbf{2-norm-constrained portfolios have the best SR, close to LW. Turnover
(excluding 1/N): lowest for GMV-C, then LW-diagonal, LW-1-factor; norm-constrained
portfolios have the highest.} Here the GMV-type strategies beat 1/N.

% ---------------------------------------------------------------------
\section{The big comparison at a glance}
\begin{center}\small
\begin{tabularx}{\textwidth}{@{}lXlll@{}}
\toprule
Strategy & Uses & Fixes & Out-of-sample SR & Turnover \\
\midrule
MV (sample) & $\hat\mub$, $\hat V$ & nothing & very poor, often $<0$ & enormous \\
Bayes--Stein & shrunk $\hat\mub$, $\hat V$ & means (partly) & poor & very high \\
Data-Model & model-based $\hat\mub$ & means & mixed & high \\
MV-C / BS-C & $\hat\mub$, $\hat V$, $w\ge0$ & means (shrinkage) & better than MV & moderate \\
GMV & $\hat V$ only & means (100\%) & good & moderate \\
GMV-C & $\hat V$, $w\ge0$ & means $+$ covariances & good & low \\
GMV-LW & $\hat V_{LW}$ & means $+$ covariances & good & low--moderate \\
GMV-NC (2-norm) & $\hat V$, $\|\w\|_2\le\delta$ & means $+$ covariances & best & high \\
1/N & nothing & everything (no estimation) & hard to beat & minimal \\
\bottomrule
\end{tabularx}
\end{center}

% ---------------------------------------------------------------------
\section{Appendix: numerical optimisation with \texttt{scipy}}
\begin{plainbox}
When there are inequality constraints ($w_n\ge0$, $w_n\le0.5$) there is no
formula: we ask a numerical solver to search for the best weights. The solver
only \textbf{minimises}, so to maximise utility we minimise \emph{minus} utility.
You need: an objective function, a starting guess (1/N), and optionally
\texttt{bounds} (one interval per weight) and \texttt{constraints} (equalities such
as $\sum w=1$).
\end{plainbox}

\subsection{Mean--variance utility (with risk-free asset)}
\begin{lstlisting}
import numpy as np
from scipy.optimize import minimize
ones = np.ones((3, 1))
gamma = 2.715745          # equivalent to mu_targ = 30% in Chapter 3
Rf = 1.065
V = np.array([[0.0625, 0.0700, 0.1050],
              [0.0700, 0.1225, 0.0840],
              [0.1050, 0.0840, 0.3600]])
mu = np.array([[1.20], [1.30], [1.40]])
Vinverse = np.linalg.inv(V)

# Analytic (no constraints)
w_mvp = (1/gamma) * Vinverse @ (mu - Rf * ones)    # [-1.0041, 1.0052, 0.4010]
w_riskFree = 1 - ones.T @ w_mvp                    # 0.5979
w_tangency = w_mvp / (ones.T @ w_mvp)              # [-2.4972, 2.5000, 0.9972]

# Numerical: minimise MINUS utility
def objective(weights):
    portfolio_return = weights @ (mu - Rf * ones).ravel()
    portfolio_risk = weights @ V @ weights
    return (gamma/2) * portfolio_risk - portfolio_return
initial_weights = np.ones(3) / 3
result = minimize(objective, initial_weights, method="SLSQP", options={"ftol": 1e-12})
print(result.x)                                    # same as analytic

# With short-sale constraint: add bounds (0, upper)
bounds = tuple((0, 999) for _ in range(3))
result = minimize(objective, initial_weights, method="SLSQP", bounds=bounds,
                  options={"ftol": 1e-12})
w = result.x                                       # [0.000, 0.5612, 0.2117]
print(1 - w.sum(), w / w.sum())                    # rf 0.2271; tangency [0, 0.7261, 0.2739]
\end{lstlisting}
\begin{addnote}
In the slides \texttt{weights.T @ (mu - Rf*ones)} mixes a 1-D array with a
$3\times1$ column; it works but returns a 1-element array. Using \texttt{.ravel()}
(as above) keeps everything 1-D and returns a scalar, which is what
\texttt{minimize} expects. With MV utility no budget constraint is needed: the
rest goes to the risk-free asset.
\end{addnote}

\subsection{GMV: analytic, numerical, with bounds}
Here we need an explicit \textbf{equality constraint} $\sum w=1$.
\begin{lstlisting}
w_gmv = (Vinverse @ ones) / (ones.T @ Vinverse @ ones)   # [1.5636, -0.3206, -0.2430]

def objective(weights, V):
    return weights @ V @ weights
constraints = ({"type": "eq", "fun": lambda weights: np.sum(weights) - 1})
initial_weights = np.ones(len(V)) / len(V)
result = minimize(objective, initial_weights, args=(V,), method="SLSQP",
                  constraints=constraints, options={"ftol": 1e-12})
# -> [1.5636, -0.3206, -0.2430], variance 0.04977

bounds = tuple((0, 999) for _ in range(len(V)))           # no short sales
result = minimize(objective, initial_weights, args=(V,), method="SLSQP",
                  bounds=bounds, constraints=constraints, options={"ftol": 1e-12})
# -> [1, 0, 0], variance 0.0625  (everything in asset 1: not diversified!)

bounds = tuple((0, 0.50) for _ in range(len(V)))          # 0 <= w <= 50%
result = minimize(objective, initial_weights, args=(V,), method="SLSQP",
                  bounds=bounds, constraints=constraints, options={"ftol": 1e-12})
# -> [0.5, 0.4944, 0.0056], variance 0.0812  (better diversified)
\end{lstlisting}
\begin{keybox}[: \texttt{minimize} cheat sheet]
\begin{itemize}
  \item \texttt{minimize(fun, x0, args=(...), method="SLSQP", bounds=..., constraints=..., options={"ftol":1e-12})}
  \item \texttt{bounds}: list/tuple of \texttt{(low, high)}, one per weight; \texttt{None} means no bound.
  \item \texttt{constraints}: dict(s) with \texttt{"type": "eq"} (must equal 0) or \texttt{"ineq"} (must be $\ge0$).
  \item Result: \texttt{result.x} (weights), \texttt{result.fun} (objective value), \texttt{result.success}.
  \item Always check \texttt{result.success} and compare with the analytic solution when one exists.
\end{itemize}
\end{keybox}

\paragraph{Norms (aside).} $\|\bm x\|_2=\sqrt{x_1^2+\dots+x_N^2}$ (Euclidean);
$\|\bm x\|_1=|x_1|+\dots+|x_N|$ (Manhattan/taxicab);
$\|\bm x\|_p=(\sum|x_n|^p)^{1/p}$; $\|\bm x\|_\infty=\max_n|x_n|$.
Frobenius norm of a matrix: $\|A\|_F=\sqrt{\sum_{n,m}a_{nm}^2}$, useful to measure
the distance between two matrices. In NumPy: \texttt{np.linalg.norm(x, 1)},
\texttt{np.linalg.norm(x)}, \texttt{np.linalg.norm(A, "fro")}.

% ---------------------------------------------------------------------
\section{Takeaways of Chapter 4}
\begin{enumerate}
  \item Means are estimated so imprecisely ($SE=\sigma/\sqrt T$) that MV weights built on them are extreme, unstable, and lose out of sample.
  \item Sample covariance matrices are ill-conditioned when $N$ is large relative to $T$; $V^{-1}$ amplifies errors.
  \item Out of sample, \textbf{1/N beats sample MV} in all six DGU datasets, and Bayesian fixes do not change that.
  \item \textbf{Shrinkage} is the unifying idea: Bayes--Stein (means), short-sale constraints (means and covariances), GMV (100\% shrinkage of means), Ledoit--Wolf and norm constraints (covariances).
  \item GMV-type strategies (GMV-C, LW, 2-norm) are the ones that can beat 1/N; constraints are excellent at cutting turnover.
  \item With inequality constraints, use \texttt{scipy.optimize.minimize}.
\end{enumerate}


\begin{mistakes}
\begin{enumerate}
  \item Believing an in-sample Sharpe ratio: it is biased upwards (overfitting).
  \item Forgetting that the GMV needs \emph{no} expected returns: putting $\hat\mub$
    in the GMV formula.
  \item Imposing $w\ge0$ and concluding the problem is solved: constraints mainly
    cut turnover; they do not make the means precise.
  \item Using \texttt{minimize} without the equality constraint $\sum w=1$ for GMV
    (the solver then returns $\w=\bm0$, which has zero variance).
  \item Not checking \texttt{result.success}, or starting the solver from a bad point.
  \item Choosing $\delta$ (norm constraint) or $\psi$ using the whole sample,
    including the out-of-sample period (look-ahead bias).
  \item Comparing strategies on Sharpe ratios alone, without turnover and costs.
\end{enumerate}
\end{mistakes}
\section{Self-test questions}
\begin{enumerate}
  \item A stock has $\sigma=40\%$. Years of data for a 95\% CI of $\pm5\%$ on the mean?
    \emph{($(2\cdot0.4/0.05)^2=256$ years.)}
  \item Why do higher correlations make MV weights more sensitive to errors in means?
    \emph{(Long--short spreads become almost riskless, so the optimiser levers up on tiny mean differences.)}
  \item What does the short-sale constraint do to the expected return of an asset whose constraint binds?
    \emph{(Raises it by $\lambda_n$: shrinkage towards the average.)}
  \item What does Jagannathan--Ma's $\hat V_{JM}$ do to asset $n$ when its constraint binds?
    \emph{(Reduces its covariances by $\lambda_n$ and its variance by $2\lambda_n$.)}
  \item Ledoit--Wolf with target $\bar\lambda I$ and $\psi=1$: what is the GMV? \emph{(1/N.)}
  \item What happens to the eigenvalues under LW identity shrinkage? \emph{($\lambda\mapsto(1-\psi)\lambda+\psi\bar\lambda$.)}
  \item 1-norm vs.\ 2-norm constraint: which gives exact zeros? \emph{(1-norm.)}
  \item Why must $\delta$ be chosen without look-ahead bias? \emph{(Otherwise the
    backtest uses future information and overstates performance.)}
  \item Why does GMV overweight low-volatility stocks? When does it earn a high SR?
  \item Write by hand the long-only GMV with \texttt{scipy}.
\begin{lstlisting}
from scipy.optimize import minimize
N = V.shape[0]
res = minimize(lambda w: w @ V @ w, np.ones(N)/N, method="SLSQP",
               bounds=[(0, 1)]*N,
               constraints=[{"type": "eq", "fun": lambda w: w.sum() - 1}])
w_gmv_c = res.x
\end{lstlisting}
  \item Write by hand the Bayes--Stein means given \texttt{mu\_hat}, \texttt{V\_hat}, \texttt{T}.
\begin{lstlisting}
N = len(mu_hat); o = np.ones(N)
Vi = np.linalg.inv(V_hat)
w_g = Vi @ o / (o @ Vi @ o); mu_g = w_g @ mu_hat
d = mu_hat - mu_g
psi = (N + 2) / ((N + 2) + T * d @ Vi @ d)
mu_bs = (1 - psi) * mu_hat + psi * mu_g
\end{lstlisting}
\end{enumerate}

\section{To do for next class}
\begin{itemize}
  \item Read DeMiguel, Garlappi, Uppal (2009), ``Optimal versus naive diversification:
    How inefficient is the 1/N portfolio strategy?'' \emph{RFS} 22(5).
  \item Moreira (2021), ch.\,9--10; Tidy Finance, ``Constrained Optimization and Backtesting''.
  \item Work on Assignment \#2 (due 16 Nov, 22:00 CET).
\end{itemize}

\section*{References (Chapter 4, selection)}
\begin{itemize}[leftmargin=*]\small
\item DeMiguel, V., L.~Garlappi, F.\,J.~Nogales, and R.~Uppal. 2009. A generalized approach to portfolio optimization: improving performance by constraining portfolio norms. \emph{Management Science} 55(5): 798--812.
\item DeMiguel, V., L.~Garlappi, and R.~Uppal. 2009. Optimal versus naive diversification: how inefficient is the 1/N portfolio strategy? \emph{RFS} 22(5): 1915--1953.
\item Jagannathan, R., and T.~Ma. 2003. Risk reduction in large portfolios: why imposing the wrong constraints helps. \emph{JF} 58(4): 1651--1684.
\item James, W., and C.~Stein. 1961. Estimation with quadratic loss. \emph{Proc.\ 4th Berkeley Symposium}: 361--379.
\item Jorion, P. 1985. International portfolio diversification with estimation risk. \emph{J.\ Business} 58(3): 259--278; 1986. Bayes--Stein estimation for portfolio analysis. \emph{JFQA} 21(3): 279--292.
\item Ledoit, O., and M.~Wolf. 2003. \emph{JEF} 10(5): 603--621; 2004a. \emph{J.\ Multivariate Analysis} 88: 365--411; 2004b. Honey, I shrunk the sample covariance matrix. \emph{JPM} 30(4): 110--119; 2017. \emph{RFS} 30(12): 4349--4388; 2022. \emph{J.\ Fin.\ Econometrics} 20(1): 187--218.
\item Merton, R.\,C. 1980. On estimating the expected return on the market. \emph{JFE} 8(4): 323--361.
\item P\'astor, L. 2000. Portfolio selection and asset pricing models. \emph{JF} 55: 179--223.
\item Pástor, L., and R.\,F.~Stambaugh. 2000. Comparing asset pricing models: an investment perspective. \emph{JFE} 56(3): 335--381.
\end{itemize}

% ===== end ch4 =====


% =====================================================================
\chapter*{References (Chapter 1)}
\addcontentsline{toc}{chapter}{References (Chapter 1)}
\markboth{References}{}
\begin{itemize}[leftmargin=*]
\item Brugi\`ere, P. 2020. \emph{Quantitative Portfolio Management}. Springer.
\item Chen, A.\,Y., and T.~Zimmermann. 2022. Open source cross-sectional asset pricing. \emph{Critical Finance Review} 11(2): 207--264.
\item Dick-Nielsen, J., P.~Feldh\"utter, L.\,H.~Pedersen, and C.~Stolborg. 2023. Corporate bond factors: replication failures and a new framework. SSRN 4586652.
\item Dick-Nielsen, J., and T.\,K.~Poulsen. 2019. How to clean academic TRACE data. SSRN 3456082.
\item Dickerson, A., P.~Mueller, and C.~Robotti. 2023. Priced risk in corporate bonds. \emph{Journal of Financial Economics}.
\item Gu, S., B.~Kelly, and D.~Xiu. 2020. Empirical asset pricing via machine learning. \emph{Review of Financial Studies} 33(5): 2223--2273.
\item Hou, K., C.~Xue, and L.~Zhang. 2015. Digesting anomalies: an investment approach. \emph{Review of Financial Studies} 28(3): 650--705.
\item Scheuch, C., S.~Voigt, P.~Weiss, and C.~Frey. 2024. \emph{Tidy Finance with Python}. Chapman \& Hall/CRC.
\item Welch, I., and A.~Goyal. 2008. A comprehensive look at the empirical performance of equity premium prediction. \emph{Review of Financial Studies} 21(4): 1455--1508.
\end{itemize}

\appendix
% ===== begin assign =====
\chapter{Assignment guides}
\label{ch:assign}

\begin{plainbox}
These guides explain, question by question, \emph{what} is asked, \emph{which
theory} you need, \emph{how} to code it, and \emph{what to watch out for}. The
assignments are not graded (full credit if submitted), but each in-class quiz
only counts if the matching assignment is submitted, and the final exam has one
question per assignment, \textbf{with code written by hand}. So use the code
below to understand, then make sure you can rewrite the key lines yourself and
explain the results in your own words.
\end{plainbox}

\section{Submission checklist (applies to every assignment)}
\begin{itemize}
  \item Groups of exactly 3 or 4. Email to \texttt{Raman.Uppal@edhec.edu}, subject
    \texttt{QPM 2026-2027: Assignment n}. Deadline 22:00 CET; late $=0$.
  \item File name \texttt{QPM-2026-Assignment-n-A-B-C-D.ipynb} (last names).
  \item Before sending: Colab \emph{Runtime $\to$ Restart session and run all};
    Jupyter \emph{Kernel $\to$ Restart Kernel and Run All Cells}. No errors.
  \item Mention any outside help (AI included): no penalty.
\end{itemize}
\paragraph{Section 0 template (Markdown cell).}
\begin{lstlisting}[language={},numbers=none]
# 0. Submission information
**Submission date:** 10 November 2026
**QPM 2026-2027: Assignment 1**
**Group members:** BIANCHI Anna, MASTROGIACOMO Paolo, MULLER Jonas, ROSSI Luca
**Comments/challenges:** Some S&P 500 tickers failed to download from Yahoo
(ignored, as allowed). We used ChatGPT/Claude to debug the download code.
\end{lstlisting}
Then, for each question: a Markdown cell with the question text, followed by
code cells and a Markdown cell with your interpretation.

% =====================================================================
\section{Assignment 1 (Lectures 1--2; due 11 November)}

\subsection*{Q1.1 Download daily FAANG prices, January 2015 -- December 2020}
\textbf{Theory:} Chapter 1, ``Stock-price data''. Use \textbf{adjusted} prices
(splits and dividends).
\begin{lstlisting}
import numpy as np
import pandas as pd
import yfinance as yf
from scipy import stats

faang = ["META", "AMZN", "AAPL", "NFLX", "GOOG"]
prices = yf.download(tickers=faang, start="2015-01-01", end="2021-01-01",
                     auto_adjust=False, progress=False)["Adj Close"]
prices = prices[faang]           # keep the column order
print(prices.shape)
prices.head()
\end{lstlisting}
\begin{addnote}
\begin{itemize}
  \item \texttt{end} is \textbf{exclusive} in \texttt{yfinance}: \texttt{end="2021-01-01"} includes 31 Dec 2020.
  \item \texttt{auto\_adjust=False} keeps the \texttt{"Adj Close"} column; with the
    default (\texttt{True}) the \texttt{"Close"} column is already adjusted.
  \item This is \emph{wide} format (dates $\times$ tickers), which is fine here because
    the next questions are matrix calculations. For tidy format:
    \texttt{prices.stack().rename("adjusted").reset\_index()}.
\end{itemize}
\end{addnote}

\subsection*{Q1.2 First and second moments of returns}
\textbf{Theory:} Chapter 1, returns and sample moments. \emph{Means},
\emph{variances}, \emph{covariances} of (daily, simple) returns.
\begin{lstlisting}
returns = prices.pct_change(fill_method=None).dropna()   # daily simple returns
mean_d = returns.mean()                                 # daily means
var_d = returns.var()                                   # daily variances (ddof=1)
cov_d = returns.cov()                                   # 5 x 5 covariance matrix
corr_d = returns.corr()

moments = pd.DataFrame({"mean (daily)": mean_d, "variance (daily)": var_d,
                        "mean (ann.)": 252 * mean_d, "vol (ann.)": np.sqrt(252 * var_d)})
print(moments.round(5)); print(cov_d.round(6)); print(corr_d.round(2))
\end{lstlisting}
\begin{plainbox}
Report daily numbers (what is asked) and also annualised ones (easier to read:
mean $\times252$, volatility $\times\sqrt{252}$). Comment on the correlation matrix:
FAANG stocks are positively correlated (common ``tech'' factor), which limits
diversification among them (Chapter 3).
\end{plainbox}

\subsection*{Q1.3 Skewness, excess kurtosis; are daily returns Normal?}
\textbf{Theory:} Chapter 1, ``Are returns Normal?''. Under Normality
$\hat S\sim\mathcal N(0,6/T)$, $\hat K\sim\mathcal N(3,24/T)$; Jarque--Bera.
\begin{lstlisting}
T = len(returns)
dist = pd.DataFrame({
    "skewness": returns.skew(),
    "excess kurtosis": returns.kurt(),               # pandas .kurt() is EXCESS kurtosis
    "z skew": returns.skew() / np.sqrt(6 / T),
    "z ex-kurt": returns.kurt() / np.sqrt(24 / T),
    "JB stat": [stats.jarque_bera(returns[c]).statistic for c in returns],
    "JB p-value": [stats.jarque_bera(returns[c]).pvalue for c in returns]})
print(dist.round(3))
\end{lstlisting}
\begin{keybox}[: what you should find and how to answer]
Daily returns are \textbf{not Normal}: excess kurtosis is large (fat tails), the
$z$-statistics are far above 1.96 and the JB $p$-values are $\approx0$. Expect
some skewness too (sign can differ across stocks; big single-day jumps after
earnings, e.g.\ Netflix, Meta). The Normal distribution underestimates the
frequency of extreme days (Chapter 1, Python exercise 5).
\end{keybox}
\begin{addnote}
\texttt{pandas .kurt()} and \texttt{scipy.stats.kurtosis(x)} (default
\texttt{fisher=True}) both return \textbf{excess} kurtosis (Normal $=0$).
\texttt{scipy.stats.kurtosis(x, fisher=False)} returns kurtosis (Normal $=3$).
pandas uses bias-corrected estimators; differences with scipy are tiny for large $T$.
\end{addnote}

\subsection*{Q2.1 S\&P 500 constituents, month-end adjusted prices 2000--2022}
\textbf{Theory:} Chapter 1, ``From a single stock to multiple stocks''.
\begin{lstlisting}
import requests
from io import StringIO

url = "https://en.wikipedia.org/wiki/List_of_S%26P_500_companies"
html = requests.get(url, headers={"User-Agent": "Mozilla/5.0"}).text  # avoids 403
sp500 = pd.read_html(StringIO(html))[0]
tickers = sp500["Symbol"].str.replace(".", "-", regex=False).tolist()  # BRK.B -> BRK-B
print(len(tickers), tickers[:10])

start_date, end_date = "2000-01-01", "2022-12-31"
raw = yf.download(tickers=tickers, start=start_date, end=end_date,
                  auto_adjust=False, progress=False)["Adj Close"]

monthly = raw.resample("ME").last()                       # month-end prices
monthly = monthly.dropna(axis=1, how="all")               # drop all-NaN columns
monthly = monthly.loc[:, monthly.isna().sum() <= 100]     # drop > 100 missing months
print(monthly.shape)                                      # (276 months, ~4xx stocks)
\end{lstlisting}
\begin{addnote}
\begin{itemize}
  \item Yahoo uses \texttt{-} instead of \texttt{.} in tickers (BRK-B, BF-B).
  \item Downloading about 500 tickers $\times$ 23 years takes a few minutes; failed
    tickers are printed and can be ignored (the assignment allows it).
  \item \texttt{"ME"} = month end (pandas $\ge$ 2.2; in older versions use \texttt{"M"}).
    Alternative: the option \texttt{interval="1mo"} of \texttt{yf.download}, but its dates are the
    \emph{first} day of each month; resampling daily data is cleaner.
  \item Save the raw download too (\texttt{raw.to\_csv("sp500\_daily.csv")}) so you do
    not need to re-download for Assignment 2.
\end{itemize}
\end{addnote}

\subsection*{Q2.2 Simple monthly returns (and save them)}
\begin{lstlisting}
ret = monthly.pct_change(fill_method=None).iloc[1:]   # first row is NaN
print(ret.index[0], ret.index[-1], ret.shape)         # 2000-02-29 ... 2022-12-31
ret.to_csv("sp500_monthly_returns.csv")
# Later: ret = pd.read_csv("sp500_monthly_returns.csv", index_col=0, parse_dates=True)
\end{lstlisting}
\begin{addnote}
In Colab, files in the session disappear when it restarts. Save to Google Drive
(\texttt{from google.colab import drive; drive.mount("/content/drive")}) or download the CSV.
\end{addnote}

\subsection*{Q2.3 Annualised mean, volatility and Sharpe ratio}
\textbf{Theory:} Chapter 2, Sharpe ratio; Chapter 1, Fama--French data. The
Sharpe ratio uses \textbf{excess} returns $r_t-RF_t$ (monthly RF from Ken French).
\begin{lstlisting}
import pandas_datareader as pdr
ff = pdr.DataReader("F-F_Research_Data_Factors", "famafrench",
                    start=start_date, end=end_date)[0] / 100     # percent -> decimal
rf = ff["RF"]                                                    # PeriodIndex (monthly)

ret_p = ret.copy()
ret_p.index = ret_p.index.to_period("M")          # align on months
excess = ret_p.sub(rf, axis=0)                    # subtract RF from every column

stats_sp = pd.DataFrame({
    "mean (ann.)": 12 * ret.mean(),
    "vol (ann.)": np.sqrt(12) * ret.std(),
    "Sharpe (ann.)": np.sqrt(12) * excess.mean() / excess.std()})
print(stats_sp.describe().round(3))                         # cross-sectional summary
print(stats_sp.sort_values("Sharpe (ann.)", ascending=False).head(10).round(3))
\end{lstlisting}
\begin{addnote}
Ken French dates are monthly \emph{periods} (\texttt{2000-01}), Yahoo month-end
dates are timestamps (\texttt{2000-01-31}): converting both to periods with
\texttt{.to\_period("M")} is the safest way to align them. Annualise SR with
$\sqrt{12}$ (Chapter 2; exact only without autocorrelation).
\end{addnote}

\subsection*{Q2.4 Should weights be chosen only from the stocks' Sharpe ratios?}
\textbf{No.} Arguments to develop (use the course):
\begin{enumerate}
  \item \textbf{Correlations are ignored.} The SR is meant for efficient portfolios,
    not individual assets (Chapter 2). Portfolio risk depends on covariances
    (Chapter 3: diversification formula, marginal contribution
    $\sigma_{np}/\sigma_p$). Two high-SR stocks that are highly correlated add
    little; a low-SR stock that hedges the others can be valuable. The optimal
    rule is $\w\propto V^{-1}(\mub-R_f\one)$, not ``rank by SR''. (Only if all
    correlations were zero would weights be $\propto(\mu_n-R_f)/\sigma_n^2=SR_n/\sigma_n$.)
  \item \textbf{Estimation error.} Each SR is estimated with
    $SE\approx\sqrt{(1+SR^2/2)/T}$; with about 23 years the SE of an annual SR is
    about $0.21$, so most differences between stocks are not significant
    (Chapter 2). Ranking by \emph{in-sample} SR is fitting noise (Chapter 4:
    in-sample vs.\ out-of-sample).
  \item \textbf{Look-ahead bias.} These SRs use data from 2000--2022; choosing
    weights with them and evaluating on the same period is in-sample.
  \item \textbf{Survivorship bias.} We used \emph{today's} S\&P 500 constituents:
    firms that survived and grew. Their past SRs are biased upwards (Chapter 1,
    Python exercise 7).
  \item \textbf{Multiple testing.} Picking the best of $\sim$480 stocks will find
    high SRs by luck alone (Chapter 2, Python exercise 5).
  \item Other: non-Normal returns (SR ignores skewness/kurtosis), turnover and
    transaction costs.
\end{enumerate}
\begin{addnote}
\textbf{Where $SE\approx0.21$ comes from.} $T=275$ months, monthly SR about 0.1--0.2:
$SE_{monthly}=\sqrt{(1+0.5\cdot0.15^2)/275}\approx0.061$; annualised
$\times\sqrt{12}\approx0.21$.
\end{addnote}

% =====================================================================
\section{Assignment 2 (Lectures 3--4; due 16 November)}
\begin{plainbox}
You run a real \textbf{rolling-window backtest} (Chapter 4 methodology) on 10
stocks with four strategies: MVP and GMV without constraints (Question 1), and
MVP-C and GMV-C with $w\ge0$ (Question 2). Then you compare out-of-sample Sharpe
ratios and explain them with the theory of Chapter 4.
\end{plainbox}

\subsection*{Q1.1 Data: 10 stocks, February 2000 -- December 2022, $R_f=0$}
\begin{lstlisting}
import numpy as np
import pandas as pd
from scipy.optimize import minimize

ret = pd.read_csv("sp500_monthly_returns.csv", index_col=0, parse_dates=True)
tickers = ["MMM", "AOS", "ABT", "ADM", "ADBE", "ADP", "AES", "AFL", "A", "AKAM"]
R = ret.loc["2000-02":"2022-12", tickers]
print(R.shape, R.isna().sum().sum())      # (275, 10) and 0 missing values
\end{lstlisting}
With $R_f=0$, excess returns $=$ returns and $SR=\sqrt{12}\,\bar r/\sigma_r$.

\subsection*{Q1.2 Estimation window of 60 months: MVP and GMV, fully invested}
\textbf{Theory:} Chapter 3 (tangency portfolio), Chapter 4 (GMV).
\[
\text{MVP:}\ \ \w=\frac1\gamma\hat V^{-1}\hat\mub\ \xrightarrow{\text{rescale}}\ \w_{MVP}=\frac{\hat V^{-1}\hat\mub}{\one\T\hat V^{-1}\hat\mub}\quad(\gamma\text{ cancels}),\qquad
\text{GMV:}\ \ \w_{GMV}=\frac{\hat V^{-1}\one}{\one\T\hat V^{-1}\one}.
\]
\begin{lstlisting}
T_est = 60
N = R.shape[1]
ones = np.ones(N)

def w_mvp(mu, V):                       # tangency with rf = 0
    x = np.linalg.solve(V, mu)
    return x / x.sum()

def w_gmv(V):
    x = np.linalg.solve(V, ones)
    return x / x.sum()

first = R.iloc[:T_est]                  # Feb 2000 - Jan 2005: estimation sample
mu_hat, V_hat = first.mean().values, first.cov().values
print(pd.DataFrame({"MVP": w_mvp(mu_hat, V_hat), "GMV": w_gmv(V_hat)},
                   index=tickers).round(3))
\end{lstlisting}

\subsection*{Q1.3 -- Q1.4 Rolling window and out-of-sample returns}
At each month $t=60,\dots,T-1$: estimate on months $t-60,\dots,t-1$, compute the
weights, apply them to the return of month $t$ (i.e.\ weights chosen at the end
of month $t-1$ earn the return of month $t$: this is ``$t+1$'' in the
assignment's notation).
\begin{lstlisting}
def run_backtest(R, strategies, T_est=60):
    weights = {k: [] for k in strategies}
    oos = {k: [] for k in strategies}
    dates = []
    for t in range(T_est, len(R)):
        window = R.iloc[t - T_est:t]          # past 60 months only (no look-ahead)
        mu_hat = window.mean().values
        V_hat = window.cov().values
        r_next = R.iloc[t].values             # next month's returns
        for k, rule in strategies.items():
            w = rule(mu_hat, V_hat)
            weights[k].append(w)
            oos[k].append(w @ r_next)         # out-of-sample portfolio return
        dates.append(R.index[t])
    oos = pd.DataFrame(oos, index=dates)
    W = {k: pd.DataFrame(v, index=dates, columns=R.columns) for k, v in weights.items()}
    return oos, W

strategies_q1 = {"MVP": lambda mu, V: w_mvp(mu, V),
                 "GMV": lambda mu, V: w_gmv(V)}
oos1, W1 = run_backtest(R, strategies_q1)
print(oos1.shape)                              # 215 out-of-sample months (Feb 2005 - Dec 2022)
\end{lstlisting}


\begin{plainbox}
\textbf{The loop, line by line.}
\begin{itemize}
  \item \texttt{for t in range(T\_est, len(R))}: $t$ goes from 60 (the 61st month,
    February 2005) to the last month. At each $t$ we pretend we are at the
    \emph{end of month $t-1$}.
  \item \texttt{window = R.iloc[t - T\_est:t]}: the 60 rows $t-60,\dots,t-1$.
    Python slices exclude the end, so month $t$ is \emph{not} in the window: no
    look-ahead.
  \item \texttt{mu\_hat}, \texttt{V\_hat}: estimates from the window only.
  \item \texttt{rule(mu\_hat, V\_hat)}: the strategy's weights, decided with past data.
  \item \texttt{w @ r\_next}: the return of month $t$ earned by those weights,
    i.e.\ the out-of-sample return.
  \item The loop runs $275-60=215$ times, so you get 215 out-of-sample returns,
    from February 2005 to December 2022.
\end{itemize}
\end{plainbox}
\subsection*{Q1.5 Sharpe ratios: which is higher, and why?}
\begin{lstlisting}
sharpe = lambda r: np.sqrt(12) * r.mean() / r.std(ddof=1)
def summarize(oos, W):
    return pd.DataFrame({
        "mean (ann.)": 12 * oos.mean(),
        "vol (ann.)": np.sqrt(12) * oos.std(ddof=1),
        "Sharpe (ann.)": oos.apply(sharpe),
        "avg max |w|": {k: W[k].abs().max(axis=1).mean() for k in W},
        "avg turnover": {k: W[k].diff().abs().sum(axis=1).mean() for k in W}})
print(summarize(oos1, W1).round(3))
\end{lstlisting}
\begin{keybox}[: how to explain the result (expected: GMV $>$ MVP)]
\begin{itemize}
  \item MVP uses $\hat\mub$, which is estimated very imprecisely
    ($SE=\sigma/\sqrt5\approx13\%$ per year for a 30\%-vol stock with a 5-year window).
    $\hat V^{-1}$ amplifies those errors $\Rightarrow$ extreme, unstable weights
    (look at ``avg max $|w|$'' and turnover).
  \item GMV ignores means (shrinks them 100\% to a common value) and relies only on
    $\hat V$, which is estimated much more precisely.
  \item This is the DGU (2009) evidence and Chapter 4, Python exercise 3. Optionally test
    the difference with the Jobson--Korkie--Memmel test (Chapter 2).
\end{itemize}
\end{keybox}
\begin{keybox}[: a trap specific to rescaling the MVP weights]
The rescaled MVP divides by $\one\T\hat V^{-1}\hat\mub$. When this number is
\textbf{close to zero} the weights explode; when it is \textbf{negative} the
``tangency'' portfolio flips sign (you end up \emph{short} the efficient
portfolio, i.e.\ on the inefficient lower half of the frontier). With 5-year
windows this happens in some months. In a simulation of this exact design,
the sum was negative in 25 of 215 months, the average largest weight was
$\pm950\%$, and the annual volatility of MVP returns was above 2{,}000\%. Check it
in your data and \textbf{mention it} in your answer: it is a direct consequence
of estimation error in the means.
\end{keybox}
\begin{lstlisting}
denoms = [np.linalg.solve(R.iloc[t-60:t].cov().values, R.iloc[t-60:t].mean().values).sum()
          for t in range(T_est, len(R))]
print("months with negative 1'V^{-1}mu:", sum(d < 0 for d in denoms),
      " smallest |denominator|:", round(min(abs(d) for d in denoms), 4))
\end{lstlisting}

\subsection*{Q2.1 -- Q2.3 MVP-C and GMV-C (no short sales), rolling window}
\textbf{Theory:} Chapter 4, short-sale constraints and appendix (\texttt{scipy}).
\[
\text{MVP-C: }\max_{\w\ge\bm0}\ \w\T\hat\mub-\frac\gamma2\w\T\hat V\w\ \text{ then rescale};\qquad
\text{GMV-C: }\min_{\w}\ \w\T\hat V\w\ \text{ s.t. }\ \one\T\w=1,\ \w\ge\bm0 .
\]
\begin{lstlisting}
def w_gmv_c(V):
    cons = {"type": "eq", "fun": lambda w: w.sum() - 1}
    res = minimize(lambda w: w @ V @ w, ones / N, method="SLSQP",
                   bounds=[(0, 1)] * N, constraints=[cons],
                   options={"ftol": 1e-12, "maxiter": 1000})
    return res.x

def w_mvp_c(mu, V, gamma=1.0):
    obj = lambda w: gamma / 2 * w @ V @ w - w @ mu       # minus MV utility
    res = minimize(obj, ones / N, method="SLSQP", bounds=[(0, None)] * N,
                   options={"ftol": 1e-12, "maxiter": 1000})
    w = np.clip(res.x, 0, None)
    if w.sum() < 1e-10:            # all estimated means <= 0: optimum is 100% risk-free
        return w_gmv_c(V)          # fallback needed to be "fully invested" (say so!)
    return w / w.sum()

strategies_q2 = {"MVP-C": lambda mu, V: w_mvp_c(mu, V),
                 "GMV-C": lambda mu, V: w_gmv_c(V)}
oos2, W2 = run_backtest(R, strategies_q2)
print(summarize(oos2, W2).round(3))
print("average number of zero weights:",
      {k: (W2[k] < 1e-6).sum(axis=1).mean() for k in W2})
\end{lstlisting}
\begin{addnote}
\textbf{Why $\gamma$ does not matter for MVP-C.} The feasible set $\{\w\ge\bm0\}$ is a
cone: if $\w^*$ solves the problem for $\gamma=1$, then $\w^*/\gamma$ solves it for
any $\gamma$ (multiply the FOCs by $1/\gamma$). After rescaling to sum to one, the
weights are identical. So MVP-C $=$ the long-only maximum-Sharpe portfolio.
(Do not add the constraint $\sum w=1$ inside the MV-utility problem: with $R_f$
available the leftover goes to the risk-free asset; rescale afterwards as asked.)
Always check \texttt{res.success}.
\end{addnote}

\subsection*{Q2.4 MVP-C vs.\ GMV-C}
Often (not always) GMV-C has the higher SR, and almost always the lower turnover
and volatility. If your data say the opposite, that is a legitimate result:
report it and explain it (see the caveats in Q2.5). Arguments: MVP-C still
uses the noisy means (the constraint only limits the damage, Chapter 4:
``constraints are less successful at improving SR''); GMV-C uses no means and the
constraint shrinks the covariance matrix (Jagannathan--Ma). MVP-C tends to
concentrate in the few stocks with the highest recent average returns (many
zero weights, little diversification).

\subsection*{Q2.5 Which of the four is best, and why?}
Put the four in one table:
\begin{lstlisting}
oos_all = pd.concat([oos1, oos2], axis=1)
W_all = {**W1, **W2}
table = summarize(oos_all, W_all)
print(table.round(3))
print("1/N Sharpe (benchmark):", round(sharpe(R.iloc[T_est:].mean(axis=1)), 3))
\end{lstlisting}
\begin{keybox}[: structure of a strong answer]
\begin{enumerate}
  \item \textbf{Report} the four SRs (and vol, turnover, max weight). MVP is
    almost surely the worst. Among the other three the ranking depends on the
    sample: GMV-C or GMV is the most common winner (DGU 2009; Jagannathan--Ma
    2003), but MVP-C can win in a period when the stocks with high past average
    returns kept doing well. Do not force the ``textbook'' answer: explain yours.
  \item \textbf{Means vs.\ no means:} strategies that ignore $\hat\mub$ (GMV, GMV-C)
    beat those that use it (MVP, MVP-C): expected returns cannot be estimated
    precisely (Merton 1980; $SE=\sigma/\sqrt T$).
  \item \textbf{Constraints as shrinkage:} $w\ge0$ is equivalent to shrinking
    the means (MV) or the covariances (GMV; $\hat V_{JM}=\hat V-\bm\lambda\one\T-\one\bm\lambda\T$),
    and cuts turnover a lot.
  \item \textbf{Caveats:} 10 stocks, one sample period (2005--2022, includes
    2008 and 2020), survivorship bias (today's constituents), no transaction
    costs; SR differences may not be significant (SE of each annual SR
    $\approx\sqrt{12\cdot(1+SR_m^2/2)/215}\approx0.24$; use JKM or a bootstrap).
  \item Optional: add 1/N and Ledoit--Wolf GMV as benchmarks.
\end{enumerate}
\end{keybox}
\begin{addnote}
\textbf{Optional extension: GMV with Ledoit--Wolf.}
\texttt{from sklearn.covariance import LedoitWolf}; inside the loop use
\texttt{V\_hat = LedoitWolf().fit(window.values).covariance\_}. With $N=10$ and
$T^{est}=60$ the gain is small (the matrix is not badly conditioned); it matters
much more for large $N$.
\end{addnote}

% ===== end assign =====

% ===== begin back =====
\chapter{Formula sheet (Chapters 1--4)}
\label{ch:formulas}
{\small
\section*{Chapter 1 -- Data and returns}
\begin{tabularx}{\textwidth}{@{}lX@{}}
\toprule
Gross / net / excess & $R_t=P_t/P_{t-1}$;\quad $r_t=R_t-1$;\quad $R_{t,n}-R_{t,f}$ (gross excess $=$ net excess) \\
With dividends & $R_t=(P_t+D_t)/P_{t-1}$ \\
Compound & $R_t(h)=\prod_{i=0}^{h-1}R_{t-i}$;\quad $\ln R_t(h)=\sum_{i=0}^{h-1}\ln R_{t-i}=p_t-p_{t-h}$ \\
Aggregation & across assets: $R_p=\sum w_nR_n$ (gross); over time: sum of log returns \\
Annualise (monthly) & mean $\times12$, vol $\times\sqrt{12}$, variance $\times12$ (daily: 252) \\
Arith.\ vs geom.\ & $\bar r_{arith}-\bar r_{geom}\approx\sigma^2/2$ \\
Lognormal & $\ln R\sim\mathcal N(\mu,\sigma^2)$: $\E[R]=e^{\mu+\sigma^2/2}$, $V[R]=e^{2\mu+\sigma^2}(e^{\sigma^2}-1)$, $\E[X^m]=e^{m\mu+m^2\sigma^2/2}$ \\
Skew / kurtosis & $S=\E[(X-\mu)^3]/\sigma^3$, $K=\E[(X-\mu)^4]/\sigma^4$; Normal: $S=0$, $K=3$; $\hat S\sim\mathcal N(0,6/T)$, $\hat K\sim\mathcal N(3,24/T)$ \\
Jarque--Bera & $JB=T[\hat S^2/6+(\hat K-3)^2/24]\sim\chi^2_2$ \\
SE of mean & $\sigma/\sqrt{\text{years}}$ (independent of sampling frequency) \\
Cov.\ parameters & $N(N+1)/2$ \\
Condition number & $\lambda_{\max}/\lambda_{\min}$ \\
\bottomrule
\end{tabularx}

\section*{Chapter 2 -- Performance}
\begin{tabularx}{\textwidth}{@{}lX@{}}
\toprule
Out of sample & $R_{t+1,p}=\w_t\T\R_{t+1}$, $\w_t$ uses info up to $t$ only \\
Sharpe & $SR=\E[R_p-R_f]/\sigma(R_p-R_f)$; annual $=\sqrt{12}\times$ monthly (IID) \\
Lo (2002) & $\eta(q)=q/\sqrt{q+2\sum_{k=1}^{q-1}(q-k)\rho^k}$ \\
Sortino & $(\E[R_p]-R_{target})/\sqrt{DV}$, $DV=\frac1T\sum\min(R_{t,p}-R_{target},0)^2$ \\
Treynor & $\E[R_p-R_f]/\beta_{p,mkt}$,\quad $\beta=C[R_p,R_m]/V[R_m]$ \\
M$^2$ / RAP & $SR_p\times\sigma(R_B)$ \\
Information ratio & $\E[R_p-R_B]/\sigma(R_p-R_B)$ \\
Jensen's alpha & $(R_p-R_f)-\beta_{p,mkt}(R_{mkt}-R_f)$ (or multifactor) \\
Max drawdown & $\max_{\tau_1<\tau_2}\{0,\,1-V(\tau_2)/V(\tau_1)\}$ \\
VaR / CVaR & $\text{VaR}_\alpha=-q_\alpha$;\quad $\text{CVaR}_\alpha=-\E[R\mid R\le q_\alpha]$ \\
Drifted weights & $w_{t^+,n}=w_{t,n}(1+r_{t+1,n})/(1+r_{t+1,p})$ \\
Turnover & $\sum_n|w_{t+1,n}-w_{t^+,n}|$; annual $=12\times$ average monthly \\
Net return & $R_{t+1,p}\big(1-\kappa\sum_n|w_{t+1,n}-w_{t^+,n}|\big)-1$ \\
Test of a mean & $z=(\bar X-\mu_0)/(\sigma/\sqrt T)$; reject if $|z|>1.96$ (5\%, two-sided); $p=2\Phi(-|z|)$ \\
Delta method & $\sqrt T(f(\hat\theta)-f(\theta))\to\mathcal N(0,\nabla f\T\Omega\nabla f)$ \\
SE of SR & $\sqrt{(1+SR^2/2)/T}$ \\
JKM test & $V[\widehat{SR}_n-\widehat{SR}_m]=\frac1T[2-2\rho+\frac12(SR_n^2+SR_m^2-2SR_nSR_m\rho^2)]$ \\
Bootstrap $p$ & $\hat p=(1+\sum_b\mathbb 1\{|S_b^*|\ge|S_{obs}|\})/(B+1)$, $S_b^*=(\widehat{\Delta SR}^{*b}-\widehat{\Delta SR})/\widehat{SE}^{*b}$ \\
Bonferroni & $p_{adj}=\min(1,K\times p)$ \\
\bottomrule
\end{tabularx}

\section*{Chapter 3 -- Mean--variance theory}
\begin{tabularx}{\textwidth}{@{}lX@{}}
\toprule
Portfolio moments & $\mu_p=\w\T\mub$, $\sigma_p^2=\w\T V\w$ \\
Diversification & $\sigma_p^2=\frac1N\overline{\text{var}}+(1-\frac1N)\overline{\text{cov}}\to\overline{\text{cov}}$ \\
``a'' numbers & $a_1=\mub\T V^{-1}\mub$, $a_2=\mub\T V^{-1}\one$, $a_3=\one\T V^{-1}\one$, $D=a_1a_3-a_2^2$ \\
Frontier (no $R_f$) & $\w=\frac1D(a_1V^{-1}\one-a_2V^{-1}\mub)+\frac1D(a_3V^{-1}\mub-a_2V^{-1}\one)\mu_{targ}$ \\
Frontier variance & $\sigma^2=(a_3\mu^2-2a_2\mu+a_1)/D$ \\
GMV & $\w=V^{-1}\one/a_3$, $\mu_{GMV}=a_2/a_3$, $\sigma^2_{GMV}=1/a_3$ \\
Frontier (with $R_f$) & $\w=\frac{\mu_{targ}-R_f}{(\mub-R_f\one)\T V^{-1}(\mub-R_f\one)}V^{-1}(\mub-R_f\one)$ \\
MV utility & $\max\ \w\T(\mub-R_f\one)-\frac\gamma2\w\T V\w\Rightarrow\w=\frac1\gamma V^{-1}(\mub-R_f\one)$ \\
Tangency & $\w_T=V^{-1}(\mub-R_f\one)/\one\T V^{-1}(\mub-R_f\one)$ (no $\gamma$) \\
Max Sharpe & $SR_T=\sqrt{(\mub-R_f\one)\T V^{-1}(\mub-R_f\one)}$ \\
Marginal contributions & $\partial\E[R_p]/\partial w_n=\E[R_n]-R_f$;\ $\partial\sigma_p^2/\partial w_n=2\sigma_{np}$;\ $\partial\sigma_p/\partial w_n=\sigma_{np}/\sigma_p$ \\
Risk share & $w_n\sigma_{np}/\sigma_p^2$ (sums to 1); vector of $\sigma_{np}$ is $V\w$ \\
Beta relation & $\E[R_n]-R_f=\beta_{np}(\E[R_p]-R_f)$, $\beta_{np}=\sigma_{np}/\sigma_p^2$, $p$ on the frontier \\
Tracking error & $TE=(\w_A-\w_B)\T\R$; vol $=\sqrt{(\w_A-\w_B)\T V(\w_A-\w_B)}$ \\
Appraisal ratio & $\alpha/\sigma(\varepsilon)$ \\
Market weights & $w_n=num_nP_n/\sum_m num_mP_m$ \\
\bottomrule
\end{tabularx}

\section*{Chapter 4 -- Estimation error}
\begin{tabularx}{\textwidth}{@{}lX@{}}
\toprule
Years for $\pm x$ (95\%) & $(2\sigma/x)^2$ \\
Diffuse prior & predictive variance $\hat V(1+1/T^{est})$ \\
Bayes--Stein & $\hat\mub^{BS}=(1-\hat\psi)\hat\mub+\hat\psi\mu^{GMV}\one$,\ $\hat\psi=\frac{N+2}{(N+2)+T(\hat\mub-\mu^{GMV}\one)\T\hat V^{-1}(\hat\mub-\mu^{GMV}\one)}$ \\
Short sales $\Rightarrow$ means & constrained MV $=$ unconstrained MV with $\mub+\bm\lambda_{ss}$ \\
Jagannathan--Ma & $\hat V_{JM}=\hat V-\bm\lambda\one\T-\one\bm\lambda\T$ \\
GMV diagonal $V$ & $w_n=(1/\sigma_n^2)/\sum_m(1/\sigma_m^2)$ \\
Ledoit--Wolf & $\hat V_{LW}=(1-\psi)\hat V+\psi\hat V_{target}$; target $\bar\lambda I$: $\lambda_n\mapsto(1-\psi)\lambda_n+\psi\bar\lambda$ \\
Frobenius norm & $\|A\|_F=\sqrt{\sum a_{nm}^2}$ \\
Norm constraints & $\|\w\|_1=\sum|w_n|\le\delta$;\ $\|\w\|_2=\sqrt{\sum w_n^2}\le\delta$;\ $\|\w\|_A=\sqrt{\w\T A\w}\le\delta$ \\
\bottomrule
\end{tabularx}
}

% =====================================================================
\chapter{Glossary (English -- Italian)}
\label{ch:glossary}
{\small
\begin{longtable}{@{}p{4.2cm}p{4.2cm}p{6.6cm}@{}}
\toprule
English & Italiano & Meaning \\
\midrule
\endhead
Asset universe & universo investibile & set of assets one can choose from \\
Backtest & backtest / verifica storica & simulate a strategy on past data \\
Benchmark & benchmark / indice di riferimento & portfolio against which a manager is judged \\
Beta & beta & sensitivity to a portfolio (market) return; systematic risk \\
Bootstrap & bootstrap / ricampionamento & resampling the data to get the distribution of a statistic \\
Condition number & numero di condizionamento & $\lambda_{\max}/\lambda_{\min}$; how close a matrix is to singular \\
Covariance matrix & matrice di varianza-covarianza & $V$ or $\Sigma$ \\
Diversification & diversificazione & reducing risk by combining assets \\
Drawdown & drawdown / perdita dal massimo & fall from a previous peak of wealth \\
Efficient frontier & frontiera efficiente & portfolios with max return for each level of risk \\
Estimation error & errore di stima & difference between estimate and true parameter \\
Estimation window & finestra di stima & past observations used to estimate ($T^{est}$) \\
Excess return & rendimento in eccesso & return minus risk-free rate \\
Expanding window & finestra espandibile & estimation sample that grows over time \\
Expected shortfall / CVaR & expected shortfall / CVaR & average loss in the worst $q\%$ of cases \\
Fat tails & code spesse & more extreme outcomes than the Normal \\
Global minimum-variance (GMV) & portafoglio a varianza minima globale & lowest-risk fully invested portfolio \\
Gross / net return & rendimento lordo / netto & $P_t/P_{t-1}$ vs.\ $P_t/P_{t-1}-1$ \\
In-sample / out-of-sample & nel campione / fuori campione & evaluated on estimation data vs.\ on new data \\
Kurtosis & curtosi & fourth standardised moment (Normal $=3$) \\
Lagrange multiplier & moltiplicatore di Lagrange & shadow price of a constraint \\
Look-ahead bias & distorsione da informazione futura & using information not available at decision time \\
Market capitalisation & capitalizzazione di mercato & shares outstanding $\times$ price \\
Mean--variance utility & utilità media-varianza & $\mu_p-\frac\gamma2\sigma_p^2$ \\
Multiple testing & test multipli & many tests inflate false discoveries \\
Myopic / dynamic portfolio & portafoglio miope / dinamico & one-period vs.\ multi-period optimisation \\
Norm constraint & vincolo di norma & limit on $\|\w\|$ \\
Power of a test & potenza del test & $1-P(\text{Type II error})$ \\
Proportional transaction cost & costo di transazione proporzionale & cost $\kappa$ per unit traded \\
Risk aversion / tolerance & avversione / tolleranza al rischio & $\gamma$ / $1/\gamma$ \\
Risk-free asset & attività priva di rischio & $R_f$ \\
Rolling window & finestra mobile & fixed-length estimation sample that moves forward \\
Sharpe ratio & indice di Sharpe & excess return per unit of volatility \\
Short sale & vendita allo scoperto & negative weight \\
Short-sale constraint & vincolo di non negatività & $w_n\ge0$ (long-only) \\
Shrinkage & shrinkage / contrazione & pulling an estimate towards a target \\
Skewness & asimmetria & third standardised moment \\
Standard error & errore standard & standard deviation of an estimator \\
Survivorship bias & distorsione da sopravvivenza & using only firms that survived \\
Systematic / unsystematic risk & rischio sistematico / idiosincratico & non-diversifiable / diversifiable risk \\
Tangency portfolio & portafoglio di tangenza & max-Sharpe portfolio of risky assets \\
Tidy (long) / wide format & formato lungo / largo & one row per (date, asset) vs.\ one column per asset \\
Tracking error & tracking error / errore di replica & portfolio minus benchmark return \\
Turnover & turnover / rotazione del portafoglio & amount traded at rebalancing \\
Two-fund separation & separazione a due fondi & every efficient portfolio mixes two funds \\
Type I / Type II error & errore di I / II tipo & false rejection / failure to reject \\
Value at Risk (VaR) & valore a rischio & loss exceeded with probability $\alpha$ \\
Volatility & volatilità & standard deviation of returns \\
\bottomrule
\end{longtable}
}

% ===== end back =====


\end{document}
